% Options for packages loaded elsewhere
\PassOptionsToPackage{unicode}{hyperref}
\PassOptionsToPackage{hyphens}{url}
\PassOptionsToPackage{dvipsnames,svgnames,x11names}{xcolor}
%
\documentclass[
  a4paper,
  13pt]{extreport}
\usepackage{amsmath,amssymb}
\usepackage{setspace}
\usepackage{iftex}
\ifPDFTeX
  \usepackage[T1]{fontenc}
  \usepackage[utf8]{inputenc}
  \usepackage{textcomp} % provide euro and other symbols
\else % if luatex or xetex
  \usepackage{unicode-math} % this also loads fontspec
  \defaultfontfeatures{Scale=MatchLowercase}
  \defaultfontfeatures[\rmfamily]{Ligatures=TeX,Scale=1}
\fi
\usepackage{lmodern}
\ifPDFTeX\else
  % xetex/luatex font selection
  \setmainfont[]{TeX Gyre Termes}
\fi
% Use upquote if available, for straight quotes in verbatim environments
\IfFileExists{upquote.sty}{\usepackage{upquote}}{}
\IfFileExists{microtype.sty}{% use microtype if available
  \usepackage[]{microtype}
  \UseMicrotypeSet[protrusion]{basicmath} % disable protrusion for tt fonts
}{}
\usepackage{xcolor}
\usepackage[left=3cm,right=2cm,top=2cm,bottom=2cm,headsep=1cm,footskip=1cm]{geometry}
\usepackage{longtable,booktabs,array}
\usepackage{calc} % for calculating minipage widths
% Correct order of tables after \paragraph or \subparagraph
\usepackage{etoolbox}
\makeatletter
\patchcmd\longtable{\par}{\if@noskipsec\mbox{}\fi\par}{}{}
\makeatother
% Allow footnotes in longtable head/foot
\IfFileExists{footnotehyper.sty}{\usepackage{footnotehyper}}{\usepackage{footnote}}
\makesavenoteenv{longtable}
\usepackage{graphicx}
\makeatletter
\def\maxwidth{\ifdim\Gin@nat@width>\linewidth\linewidth\else\Gin@nat@width\fi}
\def\maxheight{\ifdim\Gin@nat@height>\textheight\textheight\else\Gin@nat@height\fi}
\makeatother
% Scale images if necessary, so that they will not overflow the page
% margins by default, and it is still possible to overwrite the defaults
% using explicit options in \includegraphics[width, height, ...]{}
\setkeys{Gin}{width=\maxwidth,height=\maxheight,keepaspectratio}
% Set default figure placement to htbp
\makeatletter
\def\fps@figure{htbp}
\makeatother
\setlength{\emergencystretch}{3em} % prevent overfull lines
\providecommand{\tightlist}{%
  \setlength{\itemsep}{0pt}\setlength{\parskip}{0pt}}
\setcounter{secnumdepth}{-\maxdimen} % remove section numbering
\ifLuaTeX
\usepackage[bidi=basic]{babel}
\else
\usepackage[bidi=default]{babel}
\fi
\babelprovide[main,import]{vietnamese}
\ifPDFTeX
\else
\babelfont{rm}[]{TeX Gyre Termes}
\fi
% get rid of language-specific shorthands (see #6817):
\let\LanguageShortHands\languageshorthands
\def\languageshorthands#1{}
% ctu_xelatex_header.tex — header-includes cho build LaTeX CTU (xelatex).
% Ánh xạ các ký tự Unicode mà TeX Gyre Termes thiếu glyph (chỉ số dưới/trên,
% ký tự kẻ hộp) sang lệnh LaTeX tương đương để không bị rớt glyph khi biên dịch.
\usepackage{newunicodechar}
\usepackage{pifont}
\usepackage{amssymb}
\newunicodechar{✓}{\ding{51}}
\newunicodechar{✗}{\ding{55}}
\newunicodechar{≥}{\ensuremath{\geq}}
\newunicodechar{≤}{\ensuremath{\leq}}
\newunicodechar{≈}{\ensuremath{\approx}}
\newunicodechar{↔}{\ensuremath{\leftrightarrow}}
\newunicodechar{₀}{\textsubscript{0}}
\newunicodechar{₁}{\textsubscript{1}}
\newunicodechar{₂}{\textsubscript{2}}
\newunicodechar{₃}{\textsubscript{3}}
\newunicodechar{₄}{\textsubscript{4}}
\newunicodechar{₅}{\textsubscript{5}}
\newunicodechar{₆}{\textsubscript{6}}
\newunicodechar{₇}{\textsubscript{7}}
\newunicodechar{₈}{\textsubscript{8}}
\newunicodechar{₉}{\textsubscript{9}}
\newunicodechar{²}{\textsuperscript{2}}
\newunicodechar{³}{\textsuperscript{3}}
\newunicodechar{─}{\rule[0.5ex]{0.9em}{0.4pt}}

% --- Bảng rộng: thu nhỏ cỡ chữ + khoảng cột, cho phép ngắt token dài để
%     bảng không tràn lề phải (overfull hbox). Áp dụng cho mọi longtable
%     (pandoc sinh toàn bộ bảng dạng longtable).
\usepackage{etoolbox}
\usepackage{seqsplit}      % ngắt được token rất dài (mã biến, số dài)
% Lưu ý: KHÔNG đổi \tabcolsep — pandoc tính bề rộng cột theo công thức
% (\columnwidth - 4\tabcolsep)*frac (tham chiếu \tabcolsep trực tiếp), nên
% giảm \tabcolsep lại làm cột RỘNG hơn và tràn nặng hơn.
\AtBeginEnvironment{longtable}{%
  \footnotesize
  \emergencystretch=3em
  \hbadness=10000          % không cảnh báo underfull do RaggedRight
}
% Cho phép TeX bẻ dòng linh hoạt hơn trong ô p{} hẹp
\setlength{\emergencystretch}{2em}

% URL/DOI dài: cho phép ngắt ở bất kỳ ký tự nào (tránh tràn lề ở danh mục TLTK).
% Kết hợp với pandoc +autolink_bare_uris để URL trần được bọc \url{}.
\usepackage{xurl}
\ifLuaTeX
  \usepackage{selnolig}  % disable illegal ligatures
\fi
\IfFileExists{bookmark.sty}{\usepackage{bookmark}}{\usepackage{hyperref}}
\IfFileExists{xurl.sty}{\usepackage{xurl}}{} % add URL line breaks if available
\urlstyle{same}
\hypersetup{
  pdflang={vi},
  colorlinks=true,
  linkcolor={black},
  filecolor={Maroon},
  citecolor={black},
  urlcolor={blue},
  pdfcreator={LaTeX via pandoc}}

\author{}
\date{}

\begin{document}

\setstretch{1.2}
\hypertarget{chux1b0ux1a1ng-4-kux1ebft-quux1ea3-nghiuxean-cux1ee9u}{%
\chapter{CHƯƠNG 4: KẾT QUẢ NGHIÊN
CỨU}\label{chux1b0ux1a1ng-4-kux1ebft-quux1ea3-nghiuxean-cux1ee9u}}

\begin{center}\rule{0.5\linewidth}{0.5pt}\end{center}

\hypertarget{giux1edbi-thiux1ec7u-chux1b0ux1a1ng}{%
\section{Giới thiệu chương}\label{giux1edbi-thiux1ec7u-chux1b0ux1a1ng}}

Chương 4 trình bày toàn bộ kết quả nghiên cứu của luận án theo logic đi
từ bức tranh thực trạng mô tả (Mục 4.1) đến bằng chứng tổng hợp ở cấp
phân tích tổng hợp (Mục 4.2), sau đó đến bằng chứng quốc gia cụ thể và
cuối cùng là kiểm định toàn mẫu đa quốc gia. Cấu trúc trình bày gồm mười
phần chính, bắt đầu bằng thực trạng quốc tế hóa và hiệu quả hoạt động
kinh doanh tại châu Á theo các phân nhóm con thể chế ICRV, tiếp theo là:
(ii) kết quả phân tích tổng hợp ba tầng (P6); (iii) bối cảnh thể chế
tiên tiến--đổi mới tại Singapore (P4); (iv) bối cảnh thể chế trung bình
cao tại Trung Quốc (P5); (v) bối cảnh chuyển đổi bậc thấp--trung tại
Việt Nam (P3); (vi) kiểm định toàn mẫu châu Á đa bối cảnh (P7); (vii)
trường hợp biên là các nền kinh tế đảo nhỏ đang phát triển ở Thái Bình
Dương (P8); và (viii) tổng hợp kết quả theo hệ giả thuyết.

\begin{center}\rule{0.5\linewidth}{0.5pt}\end{center}

\hypertarget{thux1ef1c-trux1ea1ng-quux1ed1c-tux1ebf-huxf3a-vuxe0-hiux1ec7u-quux1ea3-houx1ea1t-ux111ux1ed9ng-kinh-doanh-tux1ea1i-chuxe2u-uxe1}{%
\section{4.1 Thực trạng quốc tế hóa và hiệu quả hoạt động kinh doanh tại
châu
Á}\label{thux1ef1c-trux1ea1ng-quux1ed1c-tux1ebf-huxf3a-vuxe0-hiux1ec7u-quux1ea3-houx1ea1t-ux111ux1ed9ng-kinh-doanh-tux1ea1i-chuxe2u-uxe1}}

\hypertarget{bux1ee9c-tranh-tux1ed5ng-thux1ec3-vux1ec1-hiux1ec7u-quux1ea3-houx1ea1t-ux111ux1ed9ng-kinh-doanh-cux1ee7a-doanh-nghiux1ec7p-chuxe2u-uxe1}{%
\subsection{4.1.1 Bức tranh tổng thể về hiệu quả hoạt động kinh doanh
của doanh nghiệp châu
Á}\label{bux1ee9c-tranh-tux1ed5ng-thux1ec3-vux1ec1-hiux1ec7u-quux1ea3-houx1ea1t-ux111ux1ed9ng-kinh-doanh-cux1ee7a-doanh-nghiux1ec7p-chuxe2u-uxe1}}

Phân tích mô tả trên khung 50 nền kinh tế châu Á và Thái Bình Dương (gồm
Nhật Bản, thành viên thứ sáu của Nhóm I được WBES khảo sát lần đầu năm
2025; xem Mục 4.1.5) cho thấy bức tranh hiệu quả hoạt động kinh doanh có
sự phân tán lớn và không hội tụ. Các bảng 4.1--4.2 dưới đây được
\textbf{tái lập trực tiếp từ dữ liệu thô} trên khung mô tả 50 nền có
năng suất hợp lệ (84.998 doanh nghiệp; 107 cặp nền--năm), bằng cùng quy
trình hài hòa hóa với pipeline kinh tế lượng P7
(\texttt{scripts/relock\_descriptives\_canonical.py}); pool phân loại
ICRV rộng hơn là 96.415 doanh nghiệp / 52 nhãn nền, khung phân tích đa
quốc gia là 88.869 / 50 nền gồm Nhật Bản, và mẫu hồi quy chính P7 M2 là
N = 81.022. Vì doanh thu WBES được báo cáo bằng nội tệ, độ phân tán năng
suất được tính \textbf{trong từng cặp nền kinh tế × năm} (bất biến với
đơn vị tiền tệ): độ lệch chuẩn của ln(doanh thu/lao động) trong đợt tăng
từ 0,96 (Nhóm I, Advanced đổi mới) lên 1,24--1,37 (Nhóm III--V), và tỷ
số P90/P10 leo từ khoảng 11 lần (Nhóm I) lên khoảng 27 lần (Nhóm V) theo
chiều từ nhóm thể chế cao đến nhóm thể chế thấp.

\textbf{Bảng 4.1.} \emph{Phân tán năng suất lao động và đặc điểm quốc tế
hóa theo phân nhóm con ICRV (khung mô tả tái lập 50 nền có năng suất hợp
lệ, 84.998 doanh nghiệp, WBES 2006--2026; pool phân loại ICRV rộng hơn
96.415/52 nhãn nền; mẫu hồi quy chính P7 M2 N = 81.022).}

\begin{longtable}[]{@{}
  >{\raggedright\arraybackslash}p{(\columnwidth - 10\tabcolsep) * \real{0.1667}}
  >{\raggedright\arraybackslash}p{(\columnwidth - 10\tabcolsep) * \real{0.1667}}
  >{\raggedright\arraybackslash}p{(\columnwidth - 10\tabcolsep) * \real{0.1667}}
  >{\raggedright\arraybackslash}p{(\columnwidth - 10\tabcolsep) * \real{0.1667}}
  >{\raggedright\arraybackslash}p{(\columnwidth - 10\tabcolsep) * \real{0.1667}}
  >{\raggedright\arraybackslash}p{(\columnwidth - 10\tabcolsep) * \real{0.1667}}@{}}
\toprule\noalign{}
\begin{minipage}[b]{\linewidth}\raggedright
Phân nhóm con ICRV
\end{minipage} & \begin{minipage}[b]{\linewidth}\raggedright
n (LP-valid)
\end{minipage} & \begin{minipage}[b]{\linewidth}\raggedright
sd log LP (trong đợt)¹
\end{minipage} & \begin{minipage}[b]{\linewidth}\raggedright
FSTS (\%)
\end{minipage} & \begin{minipage}[b]{\linewidth}\raggedright
Xuất khẩu (\%)
\end{minipage} & \begin{minipage}[b]{\linewidth}\raggedright
FDI≥10\% (\%)
\end{minipage} \\
\midrule\noalign{}
\endhead
\bottomrule\noalign{}
\endlastfoot
Nhóm I: Advanced đổi mới (SG, HK, KOR, TWN, ISR, JPN) & 5.688 & 0,96 &
10,3 & 24,3 & 7,7 \\
Nhóm II: Advanced tài nguyên (GCC) & 2.082 & \textbf{0,79}² & 3,9 & 11,6
& 13,6 \\
Nhóm III: trung bình cao (CHN, MYS, THA, KAZ\ldots) & 15.174 & 1,30 &
10,3 & 21,6 & 8,4 \\
Nhóm IV: chuyển đổi thu nhập trung bình thấp (VNM, IND, IDN, PHL, BGD,
PAK, MNG) & 42.798 & 1,24 & 8,1 & 14,1 & 3,1 \\
Nhóm V: đang nổi (LAO, KHM, NPL, LKA, JOR, AZE\ldots) & 17.340 & 1,37 &
10,2 & 18,1 & 7,0 \\
Nhóm VI: SIDS (FJI, PNG, SLB\ldots) & 1.916 & \textbf{1,40} & 6,7 & 14,4
& 20,0 \\
\textbf{Tổng (LP-valid, 50 nền)} & \textbf{84.998} & n/a & n/a & n/a &
n/a \\
\end{longtable}

\emph{Nguồn: Phân tích của tác giả từ World Bank Enterprise Surveys
(WBES), 2006--2026.}

\emph{Ghi chú: Cột n là số doanh nghiệp có \textbf{năng suất lao động
hợp lệ (LP-valid)} trong khung mô tả tái lập gồm 50 nền / 107 cặp
nền--năm (tổng 84.998), tính trực tiếp từ dữ liệu thô bằng
\texttt{scripts/relock\_descriptives\_canonical.py} với cùng quy trình
hài hòa hóa như pipeline kinh tế lượng P7. Khung mô tả này là phần lõi
có năng suất hợp lệ của khung phân tích đa quốc gia (50 nền gồm Nhật
Bản); mẫu hồi quy chính P7 là N = 81.022 (M2) và 79.080 (M5; xem Mục
4.6.1), còn pool phân loại ICRV rộng hơn (gồm cả doanh nghiệp thiếu biến
năng suất/thể chế) là 96.415 / 52 nhãn nền. Vì vậy n ở bảng này nhỏ hơn
các đếm mẫu đầy đủ dùng ở Bảng 4.3. ¹ Độ lệch chuẩn ln(năng suất lao
động) khử trung bình trong từng cặp nền kinh tế × năm rồi tổng hợp trên
các quan sát của nhóm, bất biến với đơn vị tiền tệ (thừa số quy đổi
triệt tiêu trên thang log); năng suất được winsorize 1/99 trong cụm
country-year (xem Mục 3.5.3). ² Nhóm II có phân tán hẹp nhất toàn khung
(nén năng suất kiểu kinh tế nhà nước tô); hai nền lớn nhất khối hẹp nhất
toàn pool (Saudi Arabia 0,49; Qatar 0,33).}

\emph{Quy ước nhãn nhóm ICRV (căn theo nhãn dữ liệu P7): Nhóm I =
Advanced innovation-driven; Nhóm II = tiên tiến tài nguyên; Nhóm III =
trung bình cao; \textbf{Nhóm IV = chuyển đổi thu nhập trung bình thấp}
(\texttt{Lower\_mid\_transition}: Bangladesh, Ấn Độ, Indonesia, Mông Cổ,
Pakistan, Philippines, Việt Nam); \textbf{Nhóm V = đang nổi}
(\texttt{Emerging}: 17 nền gồm Campuchia, Lào, Nepal, Sri Lanka,
Jordan\ldots); Nhóm VI = SIDS (\texttt{SIDS\_small}). Hệ nhãn này được
dùng thống nhất trong toàn bộ luận án và hai chuyên đề.}

Phát hiện quan trọng nhất từ thực trạng mô tả là sự phân biệt nội bộ
trong nhóm Advanced. Hai phân nhóm cùng mức thu nhập cao nhưng đối lập
về cấu trúc: nhóm đổi mới sáng tạo dẫn dắt có R\&D 21,0\%, ISO 32,5\% và
24,3\% doanh nghiệp xuất khẩu, trong khi nhóm tài nguyên dẫn dắt chỉ có
R\&D 7,0\% và 11,6\% doanh nghiệp xuất khẩu (Chuyên đề 1). Về phân phối,
hai nền lớn nhất Vùng Vịnh có phân tán năng suất trong đợt hẹp nhất toàn
mẫu gộp (Saudi Arabia 0,49; Qatar 0,33, so với Singapore 1,08, chênh hơn
2 lần), biểu hiện của kinh tế nhà nước tô (rentier economy) nơi phân
phối lại từ xuất khẩu dầu mỏ thu hẹp khoảng cách năng suất giữa doanh
nghiệp, nhưng không nhất thiết phản ánh hiệu quả hoạt động kinh doanh
của doanh nghiệp đáng kể. Điều này gợi ý rằng việc gộp hai loại Advanced
vào một nhóm duy nhất (như các nghiên cứu trước thường làm) sẽ che khuất
cơ chế quan trọng này.

Sự gia tăng của độ phân tán năng suất trong đợt theo phổ thể chế ICRV
(từ \textasciitilde11 lần lên \textasciitilde27 lần ở tỷ số P90/P10)
không chỉ là một quan sát mô tả mà còn nhất quán với khung lý thuyết
phân bổ sai nguồn lực (resource misallocation) của Hsieh và Klenow
(2009): trong môi trường thể chế yếu, các rào cản thị trường, tiếp cận
tài chính hạn chế, thực thi hợp đồng kém, méo mó chính sách, ngăn nguồn
lực dịch chuyển từ doanh nghiệp kém hiệu quả sang doanh nghiệp hiệu quả
hơn, khiến phương sai năng suất giữa các doanh nghiệp nới rộng. Theo
logic này, độ phân tán năng suất cao ở Nhóm V--VI \textbf{có thể} là dấu
hiệu của tổn thất phân bổ (allocative loss) thay vì phản ánh sự đa dạng
lành mạnh. Cần lưu ý rằng độ lệch chuẩn của log năng suất lao động chỉ
là \textbf{biến đại diện gián tiếp} cho cơ chế phân tán TFPR (năng suất
doanh thu) mà Hsieh và Klenow (2009) dùng để định danh phân bổ sai nguồn
lực, vì nó còn chịu ảnh hưởng của dị biệt công nghệ, mức thâm dụng vốn
và điều kiện cầu; do đó đây là một cách diễn giải \emph{nhất quán với}
khung phân bổ sai nguồn lực chứ không phải bằng chứng phân rã trực tiếp.
Dù vậy, hướng diễn giải này có hàm ý chính sách khác hẳn với việc coi
phân tán chỉ là nhiễu thống kê. Cần nhấn mạnh rằng các so sánh giữa các
nhóm ICRV trong Mục 4.1 mang bản chất \textbf{mô tả}; xác nhận
\textbf{suy diễn} rằng khác biệt giữa các chế độ thể chế có ý nghĩa
thống kê được cung cấp ở hai cấp độ độc lập trong các phần sau của
chương: kiểm định điều tiết ICRV của phân tích tổng hợp P6
(\(Q_M = 17{,}35\), \(df = 4\), \(p = {,}002\); Mục 4.2.3) và phổ điểm
uốn theo ICRV của mô hình hồi quy đa quốc gia P7 (Mục 4.6.2). Hai nguồn
bằng chứng này bảo đảm bức tranh thực trạng mô tả được neo vào bằng
chứng suy diễn nhất quán, thay vì dừng lại ở so sánh trung bình--độ lệch
chuẩn.

\begin{figure}
\centering
\includegraphics{figures/fig_4_1_pool_composition.png}
\caption{Hình 4.1. Thành phần mẫu gộp theo nhóm ICRV}
\end{figure}

\begin{figure}
\centering
\includegraphics{figures/fig_4_0_pool_by_year.png}
\caption{Hình 4.2. Quy mô mẫu gộp theo năm khảo sát}
\end{figure}

\hypertarget{thux1ef1c-trux1ea1ng-quux1ed1c-tux1ebf-huxf3a-phuxe2n-cux1ef1c-tham-gia-xuux1ea5t-khux1ea9u}{%
\subsection{4.1.2 Thực trạng quốc tế hóa: phân cực tham gia xuất
khẩu}\label{thux1ef1c-trux1ea1ng-quux1ed1c-tux1ebf-huxf3a-phuxe2n-cux1ef1c-tham-gia-xuux1ea5t-khux1ea9u}}

Trung vị FSTS bằng 0\% trong tất cả các phân nhóm, chỉ khoảng 12--24\%
doanh nghiệp tham gia xuất khẩu tùy theo bối cảnh. Phân phối FSTS có
dạng phân cực mạnh: đại đa số không xuất khẩu, trong khi một thiểu số
nhỏ có tỷ lệ xuất khẩu rất cao. Phân cực này đặt ra vấn đề phương pháp
quan trọng cho nghiên cứu quốc tế hóa và hiệu quả: kết quả hồi quy trên
toàn mẫu bị kéo bởi hành vi của thiểu số xuất khẩu tích cực.

Cường độ quốc tế hóa (FSTS trung bình, tính trên toàn mẫu gồm cả doanh
nghiệp không xuất khẩu) dao động từ 6,7\% (SIDS) đến 10,3\% (trung bình
cao), với mức trung bình khu vực khoảng 9\%. Điều đáng chú ý là trung
bình cao và Advanced có cùng FSTS trung bình (\textasciitilde10\%) nhưng
cơ chế xuất khẩu hoàn toàn khác: Advanced chủ yếu dịch vụ và hàng hóa
hàm lượng tri thức cao; trung bình cao chủ yếu sản xuất theo chuỗi giá
trị toàn cầu với biên lợi nhuận khác nhau.

\begin{figure}
\centering
\includegraphics{figures/fig_4_4_growth_pathway.png}
\caption{Hình 4.3. Quỹ đạo tăng trưởng theo nhóm ICRV}
\end{figure}

\hypertarget{thux1ef1c-trux1ea1ng-nux103ng-lux1ef1c-sux1ed1-vuxe0-cuxf4ng-nghux1ec7}{%
\subsection{4.1.3 Thực trạng năng lực số và công
nghệ}\label{thux1ef1c-trux1ea1ng-nux103ng-lux1ef1c-sux1ed1-vuxe0-cuxf4ng-nghux1ec7}}

Phân tích năm thành phần năng lực (đổi mới sản phẩm, đổi mới quy trình,
R\&D, ISO, website) theo phân nhóm con ICRV cho thấy sự đa dạng đáng kể.

\textbf{Bảng 4.2.} \emph{Đổi mới sáng tạo và chấp nhận số theo phân nhóm
con ICRV (\%).}

\begin{longtable}[]{@{}
  >{\raggedright\arraybackslash}p{(\columnwidth - 10\tabcolsep) * \real{0.1667}}
  >{\raggedright\arraybackslash}p{(\columnwidth - 10\tabcolsep) * \real{0.1667}}
  >{\raggedright\arraybackslash}p{(\columnwidth - 10\tabcolsep) * \real{0.1667}}
  >{\raggedright\arraybackslash}p{(\columnwidth - 10\tabcolsep) * \real{0.1667}}
  >{\raggedright\arraybackslash}p{(\columnwidth - 10\tabcolsep) * \real{0.1667}}
  >{\raggedright\arraybackslash}p{(\columnwidth - 10\tabcolsep) * \real{0.1667}}@{}}
\toprule\noalign{}
\begin{minipage}[b]{\linewidth}\raggedright
Nhóm ICRV (nhãn dữ liệu)
\end{minipage} & \begin{minipage}[b]{\linewidth}\raggedright
Sản phẩm mới
\end{minipage} & \begin{minipage}[b]{\linewidth}\raggedright
Quy trình mới
\end{minipage} & \begin{minipage}[b]{\linewidth}\raggedright
R\&D
\end{minipage} & \begin{minipage}[b]{\linewidth}\raggedright
ISO
\end{minipage} & \begin{minipage}[b]{\linewidth}\raggedright
Website (DAI Tầng 1)
\end{minipage} \\
\midrule\noalign{}
\endhead
\bottomrule\noalign{}
\endlastfoot
I, Adv. đổi mới & 26,9 & 18,8 & \textbf{21,0} & 32,5 & \textbf{69,2} \\
II, Adv. tài nguyên (GCC) & 17,0 & 8,6 & 7,0 & 15,2 & 50,2 \\
III, Trung bình cao & 26,7 & 18,9 & 21,5 & 31,4 & 56,9 \\
IV, Chuyển đổi TB-thấp & 17,1 & 14,7 & 15,5 & 25,3 & 48,3 \\
V, Đang nổi & 24,0 & 15,9 & 17,6 & 13,0 & 40,6 \\
VI, SIDS & \textbf{34,2} & \textbf{24,7} & 10,2 & 12,2 & 41,5 \\
\end{longtable}

\emph{Nguồn: Phân tích của tác giả từ WBES, 2006--2026 (tái lập:
\texttt{scripts/relock\_descriptives\_canonical.py}).}

\emph{Ghi chú: Bảng 4.2 được tái lập trực tiếp từ dữ liệu thô WBES trên
cùng khung mô tả canonical 50 nền (một cross-section chuẩn cho mỗi cặp
nền kinh tế × năm, đợt khảo sát ≥ 2006) như Bảng 4.1, với các chỉ báo
chuẩn của bộ công cụ WBES toàn cầu: đổi mới sản phẩm (h1), đổi mới quy
trình (h5), chi R\&D (h8), chứng nhận chất lượng quốc tế (b8) và website
(c22b). Toàn bộ sáu nhóm ICRV đều có độ phủ dữ liệu đầy đủ, kể cả Nhóm
II (đủ 6/6 nền GCC). Phổ năng lực thể chế thể hiện rõ ở các biến nền
tảng: R\&D 21,0\% (Nhóm I) xuống 10,2\% (Nhóm VI); ISO 32,5\% xuống
12,2\%; website 69,2\% xuống 41,5\%.}

Phát hiện đáng chú ý là khuôn mẫu ``thích nghi và nhảy vọt'' (adaptive
leapfrogging) ở SIDS: với dữ liệu raw đủ 8/8 nền Nhóm VI, SIDS có
\textbf{tỷ lệ đổi mới sản phẩm cao nhất mọi nhóm (34,2\%)} và đổi mới
quy trình cao nhất (24,7\%), trong khi website (41,5\%) ngang nhóm Đang
nổi (40,6\%) dù GNI thấp hơn đáng kể, phản ánh đổi mới bằng nguồn lực
hạn chế (frugal innovation) như công cụ sinh tồn chứ không phải tối ưu
hóa hiệu quả. Kết quả này hỗ trợ lập luận lý thuyết về sự phân biệt giữa
DAI (chấp nhận số bề mặt) và TCI (năng lực công nghệ chiều sâu): SIDS
đổi mới/áp dụng số ở bề mặt cao nhưng năng lực công nghệ chiều sâu (R\&D
10,2\%, ISO 12,2\%) thuộc nhóm thấp nhất, và đây là nhóm duy nhất nơi
dạng chữ U ngược hoàn toàn mất cấu trúc (FIP là trường hợp giới hạn lý
thuyết).

Một biểu hiện khác của bão hòa Tầng 1 nằm ở tương quan giữa website và
năng suất: phần bù tương quan của DAI Tầng 1 thấp nhất ở nhóm Advanced
đổi mới (+0,093, so với +0,183--0,202 ở các nhóm chuyển đổi và tài
nguyên, Chuyên đề 1, Bảng 2.3.8.1) dù tỷ lệ có website của nhóm này cao
(69,2\%). Khi một chỉ báo nhị phân tiệm cận bão hòa, phương sai của nó
co lại và nhóm ``không website'' còn lại mang tính chọn lọc đặc thù, nên
hệ số đo được phản ánh giới hạn của thước đo Tầng 1 nhiều hơn là giới
hạn của số hóa, một vấn đề đo lường cần phân biệt với quan hệ nhân quả
thực.

\begin{figure}
\centering
\includegraphics{figures/fig_4_6_firm_size.png}
\caption{Hình 4.4. Phân bố quy mô doanh nghiệp theo nhóm ICRV}
\end{figure}

\hypertarget{bux1ed1n-ruxe0o-cux1ea3n-cux1ea5u-truxfac-ux1ea3nh-hux1b0ux1edfng-ux111ux1ebfn-quux1ed1c-tux1ebf-huxf3a-vuxe0-hiux1ec7u-quux1ea3}{%
\subsection{4.1.4 Bốn rào cản cấu trúc ảnh hưởng đến quốc tế hóa và hiệu
quả}\label{bux1ed1n-ruxe0o-cux1ea3n-cux1ea5u-truxfac-ux1ea3nh-hux1b0ux1edfng-ux111ux1ebfn-quux1ed1c-tux1ebf-huxf3a-vuxe0-hiux1ec7u-quux1ea3}}

Phân tích thực trạng xác định bốn rào cản cấu trúc phổ biến nhất ảnh
hưởng đến khả năng quốc tế hóa và hiệu quả hoạt động kinh doanh của
doanh nghiệp tại châu Á: (i) tiếp cận tài chính hạn chế, đặc biệt nghiêm
trọng ở đang nổi và chuyển đổi thu nhập trung bình thấp (Nhóm V--IV);
(ii) thiếu hụt lao động có kỹ năng; (iii) cạnh tranh từ doanh nghiệp phi
chính thức; và (iv) nguồn cung điện không đáng tin cậy. Bốn rào cản này
hình thành ``khoảng trống thể chế'' (Khanna \& Palepu, 2010) tạo ra chi
phí giao dịch bổ sung cho doanh nghiệp khi mở rộng hoạt động ra ngoài
biên giới. Sự hiện diện đồng thời của nhiều khoảng trống tại đang nổi và
SIDS giải thích tại sao quan hệ quốc tế hóa và hiệu quả ở các nhóm này
yếu hoặc âm.

\hypertarget{nhux1eadt-bux1ea3n-2025-lux1ea7n-ux111ux1ea7u-tiuxean-ux111ux1b0ux1ee3c-khux1ea3o-suxe1t-vuxe0-vux1ecb-truxed-trong-khung-icrv}{%
\subsection{4.1.5 Nhật Bản 2025: Lần đầu tiên được khảo sát và vị trí
trong khung
ICRV}\label{nhux1eadt-bux1ea3n-2025-lux1ea7n-ux111ux1ea7u-tiuxean-ux111ux1b0ux1ee3c-khux1ea3o-suxe1t-vuxe0-vux1ecb-truxed-trong-khung-icrv}}

Một sự kiện dữ liệu đáng chú ý của giai đoạn nghiên cứu là việc Nhật Bản
lần đầu tiên được đưa vào Khảo sát Doanh nghiệp của Ngân hàng Thế giới
năm 2025, với mẫu \(N = 2.168\) doanh nghiệp. Nhật Bản là \textbf{thành
viên thứ sáu của Nhóm I} (Advanced đổi mới sáng tạo dẫn dắt) và là nền
kinh tế tiên tiến lớn cuối cùng gia nhập bộ dữ liệu WBES, cung cấp một
điểm neo quý giá cho biên trên của phổ thể chế. Nhật Bản được tích hợp
đầy đủ vào khung mô tả 50 nền của chuyên đề thực trạng (Chuyên đề 1) và
mọi thống kê Nhóm I trong luận án. Về phía ước lượng kinh tế lượng,
\textbf{điểm uốn toàn mẫu P7 đã được tái ước lượng trên khung 50 nền bao
gồm Nhật Bản} (2.168 doanh nghiệp Nhật trong khung phân tích 88.869; mẫu
hồi quy chính M2 N = 81.022, M5 N = 79.080; xem Mục 4.6), nên các kết
quả P7 trong luận án đã phản ánh Nhật Bản. Hai nghiên cứu thành phần còn
lại là mẫu gộp con riêng không liên quan Nhật Bản: P8 (FIP) trên 7 nền
đảo nhỏ Thái Bình Dương và P9 (ngưỡng) trên Ấn Độ. Ngoài ra, nghiên cứu
thành phần P10 kiểm định chuyên sâu quan hệ quốc tế hóa và hiệu quả tại
Nhật Bản (xem Mục 4.9). Toàn bộ số liệu dưới đây được tính trực tiếp từ
tệp vi mô gốc.

\textbf{Bảng 4.3.} \emph{Đặc trưng doanh nghiệp Nhật Bản 2025 so với năm
nền còn lại của Nhóm I (Advanced đổi mới).}

\begin{longtable}[]{@{}lrr@{}}
\toprule\noalign{}
Chỉ tiêu & Nhật Bản 2025 & 5 nền Nhóm I còn lại \\
\midrule\noalign{}
\endhead
\bottomrule\noalign{}
\endlastfoot
n doanh nghiệp & 2.168 & 4.222 \\
sd ln(LP) trong đợt & \textbf{0,89} & 1,04 \\
P90/P10 năng suất & \textbf{8,3} & 12,2 \\
FSTS trung bình (\%) & 4,1 & 13,4 \\
Doanh nghiệp xuất khẩu (\%) & 16,5 & 28,4 \\
FDI ≥ 10\% (\%) & \textbf{1,9} & 10,6 \\
Website (\%) & \textbf{83,8} & 61,7 \\
R\&D (\%) & 21,9 & 20,5 \\
ISO (\%) & 27,7 & 35,0 \\
Đổi mới sản phẩm (\%) & 32,2 & 24,3 \\
Nữ quản lý cấp cao (\%) & \textbf{7,3} & 27,5 \\
Tuổi doanh nghiệp trung bình (năm) & \textbf{50,4} & 23,0 \\
\end{longtable}

\emph{Nguồn: tính trực tiếp từ \texttt{Japan-2025-full-data.dta} (WBES)
và pool mô tả của năm nền còn lại Nhóm I. FSTS = d3b + d3c.}

Bốn đặc trưng nổi bật định vị Nhật Bản trong khung ICRV. Thứ nhất,
\textbf{phân tán năng suất hẹp nhất}: độ lệch chuẩn ln năng suất trong
đợt của Nhật (0,89) thấp hơn mức trung vị của năm nền còn lại Nhóm I
(1,04), và tỷ số P90/P10 chỉ 8,3 lần so với 12,2 lần, phản ánh một nền
kinh tế tiên tiến trưởng thành với sàn năng suất cao và đồng đều. Đặc
điểm này củng cố trực tiếp logic ``bão hòa Tầng 1'' của khung CDCM: ở
nền kinh tế đồng đều và số hóa cao, dư địa để hiện diện số cơ bản tạo
khác biệt năng suất là rất hẹp. Thứ hai, \textbf{số hóa trưởng thành
nhất}: tỷ lệ có website 83,8\% là cao nhất trong toàn bộ mẫu gộp, vượt
xa mức 61,7\% của năm nền còn lại Nhóm I. Thứ ba, \textbf{trường thọ
doanh nghiệp}: tuổi trung bình 50,4 năm (so với 23,0 năm của năm nền còn
lại Nhóm I) phản ánh cấu trúc doanh nghiệp lâu đời đặc thù của Nhật Bản.
Thứ tư, \textbf{nghịch lý đổi mới--quốc tế hóa và khoảng cách giới}: dù
cường độ R\&D và đổi mới cao, cường độ quốc tế hóa cấp cơ sở lại thấp
(FSTS 4,1\%) cùng tỷ lệ sở hữu nước ngoài rất thấp (1,9\%), cho thấy
hoạt động quốc tế hóa tập trung ở các tập đoàn đa quốc gia lớn chứ không
phân bố đều ở mẫu cơ sở; đồng thời tỷ lệ nữ quản lý cấp cao chỉ 7,3\%,
thấp hơn nhiều so với mức dẫn đầu khu vực Đông Á--Thái Bình Dương, nhất
quán với tài liệu về trần kính trong quản trị doanh nghiệp Nhật Bản. Các
đặc trưng này được khai thác đầy đủ trong nghiên cứu thành phần P10,
kiểm định chuyên sâu quan hệ quốc tế hóa và hiệu quả tại Nhật Bản: kết
quả xác nhận quan hệ gần tuyến tính dương (hiệu ứng tuyến tính +0,671, p
\textless{} 0,001; số hạng bậc hai không ý nghĩa; điểm uốn đại số
90--108\% nằm ngoài miền dữ liệu), TCI (+0,100, p \textless{} 0,001) và
DAI (+0,110, p = 0,028) nâng mặt bằng không uốn đường cong, cùng phần bù
trường thọ doanh nghiệp (+0,073, p = 0,007), đúng như dự đoán biên trên
của khung ICRV (trình bày đầy đủ ở Mục 4.9).

\begin{center}\rule{0.5\linewidth}{0.5pt}\end{center}

\hypertarget{so-suxe1nh-xuyuxean-nhuxf3m-icrv-phux1ed5-thux1ec3-chux1ebf-tux1b0ux1a1ng-quan-vuxe0-ux111ux1ed9-phuxe2n-tuxe1n-hiux1ec7u-quux1ea3}{%
\subsection{4.1.6 So sánh xuyên nhóm ICRV: phổ thể chế tương quan và độ
phân tán hiệu
quả}\label{so-suxe1nh-xuyuxean-nhuxf3m-icrv-phux1ed5-thux1ec3-chux1ebf-tux1b0ux1a1ng-quan-vuxe0-ux111ux1ed9-phuxe2n-tuxe1n-hiux1ec7u-quux1ea3}}

Khi đặt cạnh nhau toàn bộ các chỉ báo mô tả theo sáu nhóm ICRV, một bức
tranh có cấu trúc nổi lên trong đó cường độ của các tương quan thô suy
giảm dần (đến mức triệt tiêu ở SIDS) theo chiều suy yếu của thể chế. Hệ
số tương quan Pearson giữa năng suất lao động chuẩn hóa và cường độ quốc
tế hóa giảm đơn điệu theo phổ thể chế: từ \(+0{,}141\) ở Nhóm I, qua
\(+0{,}122\) ở Nhóm II, \(+0{,}061\) ở Nhóm III, \(+0{,}049\) ở Nhóm IV,
\(+0{,}015\) ở Nhóm V, rồi mất ý nghĩa thống kê ở SIDS (\(+0{,}002\),
không có ý nghĩa). Mô thức này cho thấy quan hệ thô giữa quốc tế hóa và
năng suất gần như biến mất ở các môi trường thể chế yếu, gợi ý một quan
hệ phi tuyến và bị điều tiết thay vì một hiệu ứng tuyến tính phổ quát;
SIDS là nhóm duy nhất nơi tương quan này triệt tiêu hoàn toàn (gần 0,
không ý nghĩa), nhất quán với giả thuyết chi phí buộc phải quốc tế hóa,
theo đó doanh nghiệp đảo nhỏ xuất khẩu vì bắt buộc chứ không phải vì
được chọn lọc theo năng suất.

\begin{figure}
\centering
\includegraphics{figures/fig_icrv_capability_gradient.png}
\caption{Hình 4.5. Phổ năng lực công nghệ theo sáu nhóm ICRV}
\end{figure}

Ngược với phổ thể chế của quốc tế hóa, năng lực công nghệ là yếu tố
tương quan mạnh và đồng đều nhất ở mọi nhóm, dao động từ \(+0{,}094\) ở
SIDS đến \(+0{,}180\) ở Nhóm II, củng cố cách đọc rằng năng lực công
nghệ nền tảng là một yếu tố dịch chuyển mức năng suất có tính phổ quát
chứ không phụ thuộc chế độ thể chế. Tương quan của sở hữu nước ngoài
dương ở mọi nhóm nhưng mạnh hơn ở hai cực của phổ thể chế (Nhóm I
\(+0{,}105\) và SIDS \(+0{,}104\)), phản ánh hai cơ chế đầu tư trực tiếp
khác nhau đã nêu ở Mục 4.1.1: trung tâm khu vực của doanh nghiệp đa quốc
gia tại nhóm tiên tiến so với đầu tư ngoại sinh hướng du lịch và viễn
thông tại các đảo nhỏ.

\begin{figure}
\centering
\includegraphics{figures/fig_4_2_productivity_density.png}
\caption{Hình 4.6. Phân bố mật độ năng suất lao động theo nhóm ICRV}
\end{figure}

Độ phân tán hiệu quả cũng thể hiện phổ thể chế rõ rệt khi đo bằng các tỷ
số phân vị bất biến với đơn vị tiền tệ. Tỷ số liên tứ phân vị
\(P75/P25\) của năng suất tăng từ \(3{,}2\) lần ở Nhóm I lên \(5{,}5\)
lần ở Nhóm V, một mô thức song hành với tỷ số \(P90/P10\) đã trình bày ở
Bảng 4.1. Bổ sung cho thước đo phân tán, mức sinh lời cũng mỏng dần theo
chiều thể chế yếu: tỷ suất sinh lời trên doanh thu trung vị giảm từ
\(0{,}50\) ở Nhóm I xuống \(0{,}45\) ở Nhóm III rồi \(0{,}35\) ở Nhóm
IV, nhất quán với lập luận của lý thuyết thể chế rằng chi phí giao dịch
cao trong môi trường thể chế yếu bào mòn biên lợi nhuận của doanh
nghiệp.

\begin{figure}
\centering
\includegraphics{figures/fig_4_7_spider_chart.png}
\caption{Hình 4.7. Biểu đồ radar năng lực và quốc tế hóa theo nhóm ICRV}
\end{figure}

Một điểm tương phản đáng chú ý về số hóa nằm ở chính nhóm tiên tiến.
Phần bù tương quan của chấp nhận số bề mặt (website, thuộc Tầng 1) lại
thấp nhất ở Nhóm I (\(+0{,}093\)), trong khi cao nhất ở Nhóm II
(\(+0{,}202\)) và Nhóm IV (\(+0{,}183\)), dù tỷ lệ có website ở Nhóm I
đã đạt \(69{,}2\%\). Đây là biểu hiện trực tiếp của hiện tượng bão hòa
thước đo Tầng 1: khi một chỉ báo nhị phân tiệm cận trần, phương sai của
nó co lại và lợi suất biên của hiện diện số cơ bản giảm dần, một vấn đề
đo lường cần phân biệt với giới hạn của số hóa thực chất. Cuối cùng, về
cấu trúc lãnh đạo, khu vực Đông Á và Thái Bình Dương dẫn đầu toàn cầu về
tỷ lệ nữ trong nhóm quản trị cấp cao (\(33{,}4\%\)) lẫn tỷ lệ nữ sở hữu
đa số (\(25{,}7\%\)), bỏ xa khu vực Trung Đông và Bắc Phi (tương ứng
\(3{,}3\%\) và \(1{,}7\%\)), một khác biệt thể chế kép gắn với bối cảnh
nhà nước tô ở khối tài nguyên dẫn dắt và làm phong phú thêm bức tranh đa
dạng nội bộ của các nhóm tiên tiến.

\hypertarget{bux1eb1ng-chux1ee9ng-nux1ec1n-tux1ea3ng-ux111a-quux1ed1c-gia-phuxe2n-tux1ea7ng-nux103ng-suux1ea5t-vuxf9ng-vuxe0-hiux1ec7u-ux1ee9ng-luxe1-chux1eafn-sux1ed1-p1}{%
\subsection{4.1.7 Bằng chứng nền tảng đa quốc gia: phân tầng năng suất
vùng và hiệu ứng lá chắn số
(P1)}\label{bux1eb1ng-chux1ee9ng-nux1ec1n-tux1ea3ng-ux111a-quux1ed1c-gia-phuxe2n-tux1ea7ng-nux103ng-suux1ea5t-vuxf9ng-vuxe0-hiux1ec7u-ux1ee9ng-luxe1-chux1eafn-sux1ed1-p1}}

Trước khi đi vào các kiểm định đường cong theo từng bối cảnh, nghiên cứu
thành phần P1 cung cấp một nền tảng thực nghiệm cấp doanh nghiệp về dị
biệt năng suất và vai trò của năng lực trong khu vực, trên mẫu phân tích
\textbf{40.633 doanh nghiệp} thuộc \textbf{17 nền kinh tế châu Á} khảo
sát trong giai đoạn 2009--2024. Mẫu được nhóm thành ba cụm vùng: Đại
Trung Hoa (5.251 doanh nghiệp), ASEAN (12.068 doanh nghiệp) và Nam Á
(23.314 doanh nghiệp). Biến phụ thuộc là năng suất lao động dạng
logarit, winsorize tại bách phân vị 1 và 99 để kiểm soát giá trị ngoại
lai.

Kết quả hồi quy bình phương nhỏ nhất phân tầng theo từng khối biến cho
thấy một cấu trúc rõ ràng. Mô hình chỉ với biến vùng (M1) xác nhận phân
tầng năng suất sâu sắc: so với ASEAN, doanh nghiệp Đại Trung Hoa có năng
suất cao hơn đáng kể (\(\hat{\beta} = 0{,}990\), \(p < {,}001\), tương
ứng khoảng 169,1\% cao hơn) và doanh nghiệp Nam Á cũng vượt ASEAN
(\(\hat{\beta} = 0{,}506\), \(p < {,}001\), phần bù khoảng 65,9\%);
riêng biến vùng giải thích 5,4\% phương sai năng suất
(\(R^2 = 0{,}054\)). Khi bổ sung khối năng lực doanh nghiệp (M2), chấp
nhận công nghệ gắn mạnh và dương với năng suất
(\(\hat{\beta} = 0{,}409\), \(p < {,}001\), tương ứng khoảng 50,5\% cao
hơn ở doanh nghiệp có áp dụng), và chứng nhận chất lượng quốc tế có hiệu
ứng tương tự (\(\hat{\beta} = 0{,}322\), \(p < {,}001\), khoảng 38,0\%);
độ phù hợp mô hình tăng gần gấp đôi, từ \(R^2 = 0{,}054\) lên
\(0{,}105\).

Phát hiện then chốt nằm ở mô hình tương tác (M3). Chỉ số rào cản kinh
doanh tác động âm và có ý nghĩa lên năng suất
(\(\hat{\beta} = -0{,}076\), \(p < {,}001\) ở hiệu ứng chính), nhưng
tương tác giữa rào cản và chấp nhận công nghệ dương và có ý nghĩa cao
(\(\hat{\beta} = +0{,}110\), \(p < {,}001\)). Đây là bằng chứng định
lượng trực tiếp cho \textbf{hiệu ứng lá chắn số}: chấp nhận công nghệ
làm giảm khoảng 83\% tổn thất năng suất do rào cản thể chế gây ra, theo
đó ở doanh nghiệp không áp dụng công nghệ, mỗi đơn vị tăng của rào cản
làm giảm năng suất khoảng 12,5\%, trong khi ở doanh nghiệp có áp dụng,
tổn thất biên gần như bị triệt tiêu. Khoảng cách năng suất của doanh
nghiệp nhỏ và vừa cũng phụ thuộc vùng: bất lợi quy mô rõ rệt ở ASEAN
(tương tác \(\hat{\beta} = -0{,}390\), \(p < {,}001\)) nhưng gần như
không tồn tại ở Đại Trung Hoa (\(\hat{\beta} = 0{,}043\), không ý
nghĩa). Bằng chứng nền tảng này định khung cho toàn bộ chương: dị biệt
năng suất giữa các nền kinh tế châu Á không phải ngẫu nhiên mà gắn với
năng lực doanh nghiệp và chất lượng thể chế, và năng lực số vận hành như
cơ chế bù đắp cho thiếu hụt thể chế, đúng theo logic của khung CDCM.

\hypertarget{phux1ed5-biuxean-lux1ee3i-nhuux1eadn-vuxe0-hux1ed3-sux1a1-ux111ux1eb7c-trux1b0ng-doanh-nghiux1ec7p-theo-nhuxf3m-icrv}{%
\subsection{4.1.8 Phổ biên lợi nhuận và hồ sơ đặc trưng doanh nghiệp
theo nhóm
ICRV}\label{phux1ed5-biuxean-lux1ee3i-nhuux1eadn-vuxe0-hux1ed3-sux1a1-ux111ux1eb7c-trux1b0ng-doanh-nghiux1ec7p-theo-nhuxf3m-icrv}}

Phân tích phân tán năng suất ở các mục trước dựa trên thước đo bất biến
tiền tệ trong từng cặp nền kinh tế và năm, nên không cho phép so sánh
\emph{mức} hiệu quả giữa các nhóm. Để bổ sung một so sánh mức hợp lệ
giữa các nền kinh tế, luận án sử dụng tỷ suất lợi nhuận trên doanh thu
(ROS), một tỷ số vô đơn vị không chịu ảnh hưởng của khác biệt đơn vị
tiền tệ. Bảng 4.4 trình bày trung vị ROS theo sáu nhóm ICRV.

\textbf{Bảng 4.4.} \emph{Trung vị tỷ suất lợi nhuận trên doanh thu (ROS)
theo nhóm ICRV (mẫu con có dữ liệu chi phí).}

\begin{longtable}[]{@{}lrr@{}}
\toprule\noalign{}
Nhóm ICRV & n & Trung vị ROS \\
\midrule\noalign{}
\endhead
\bottomrule\noalign{}
\endlastfoot
I, Tiên tiến đổi mới & 3.337 & 0,503 \\
II, Tiên tiến tài nguyên & 118¹ & 0,470 \\
III, Trung bình cao & 13.401 & 0,455 \\
IV, Chuyển đổi thu nhập trung bình thấp & 37.792 & 0,351 \\
V, Đang nổi & 12.523 & 0,408 \\
VI, Đảo nhỏ & 520¹ & 0,474 \\
\end{longtable}

\emph{Ghi chú: trung vị được ưu tiên vì trung bình ở Nhóm V và VI bị méo
bởi giá trị ngoại lai. ¹ Mẫu nhỏ do ít doanh nghiệp thuộc khối tài
nguyên và đảo nhỏ có dữ liệu chi phí đầy đủ. Nguồn: ước lượng của tác
giả từ dữ liệu thô WBES (tái lập: \texttt{scripts/relock\_ros.py}).}

Mô thức ROS bổ sung một chiều thông tin cho bức tranh phân tán. Biên lợi
nhuận cao nhất ở chế độ tiên tiến đổi mới (khoảng 0,50) và thấp nhất ở
chế độ chuyển đổi thu nhập trung bình thấp (khoảng 0,35); đây là phổ thể
chế mức hợp lệ giữa các nền kinh tế, nhất quán với lý thuyết thể chế
theo đó thể chế mạnh kéo theo chi phí giao dịch thấp và biên lợi nhuận
cao hơn. Cần lưu ý rằng phổ thể chế này không đơn điệu hoàn toàn: các
nhóm thể chế yếu nhất (V và VI) có trung vị ROS cao hơn Nhóm IV, một
phần do mẫu nhỏ và do thành phần ngành dịch vụ biên cao chiếm tỷ trọng
lớn ở các nền đảo nhỏ. Điểm cốt lõi là sự khác biệt mức giữa các chế độ
thể chế hiện ra ở thước đo lợi nhuận hợp lệ, trong khi ở thước đo năng
suất thô nó bị nhiễu tiền tệ che khuất, nên luận án đọc ROS như chỉ báo
mức và đọc độ phân tán năng suất như chỉ báo phân bổ.

Hồ sơ đặc trưng doanh nghiệp theo nhóm cũng phơi bày những tương phản
đáng chú ý không trùng với thứ bậc thu nhập. Cường độ xuất khẩu trung
bình cao nhất ở Nhóm I (13,3\%) và thấp nhất ở Nhóm II tài nguyên
(3,5\%), phản ánh định hướng nội địa và tài nguyên của khối sau. Tỷ lệ
doanh nghiệp có hiện diện số bề mặt giảm dần khá đều theo phổ thể chế,
từ 61,7\% ở Nhóm I xuống 41,2\% ở Nhóm V, phù hợp với khoảng cách hạ
tầng số. Đáng chú ý nhất, tỷ lệ nữ trong quản trị cấp cao không theo phổ
thể chế thu nhập mà cao ở cả Nhóm I (27,5\%), Nhóm III (29,1\%) và Nhóm
VI đảo nhỏ (31,8\%) trong khi rất thấp ở Nhóm II tài nguyên (4,0\%), một
khác biệt thể chế kép gắn với bối cảnh văn hóa và nhà nước tô của khối
tài nguyên dẫn dắt. Hồ sơ này củng cố lập luận xuyên suốt rằng các nhóm
ICRV khác nhau về chất chứ không chỉ về mức phát triển, và do đó đòi hỏi
cách đọc theo từng chế độ thay vì một thang tuyến tính duy nhất.

\begin{center}\rule{0.5\linewidth}{0.5pt}\end{center}

\hypertarget{kux1ebft-quux1ea3-tux1ed5ng-hux1ee3p-ux111ux1ecbnh-lux1b0ux1ee3ng-phuxe2n-tuxedch-tux1ed5ng-hux1ee3p-ba-tux1ea7ng-p6}{%
\section{4.2 Kết quả tổng hợp định lượng: phân tích tổng hợp ba tầng
(P6)}\label{kux1ebft-quux1ea3-tux1ed5ng-hux1ee3p-ux111ux1ecbnh-lux1b0ux1ee3ng-phuxe2n-tuxedch-tux1ed5ng-hux1ee3p-ba-tux1ea7ng-p6}}

\hypertarget{chux1ec9-sux1ed1-hiux1ec7u-ux1ee9ng-tux1ed5ng-hux1ee3p}{%
\subsection{4.2.1 Chỉ số hiệu ứng tổng
hợp}\label{chux1ec9-sux1ed1-hiux1ec7u-ux1ee9ng-tux1ed5ng-hux1ee3p}}

Phân tích phân tích tổng hợp tổng hợp ba tầng (three-level multilevel
meta-analysis, MARA) được tiến hành trên tập dữ liệu gồm \(k = 238\)
nghiên cứu với tổng cộng K = 288 cỡ ảnh hưởng đa dạng theo vùng địa lý,
giai đoạn và phương pháp. Kết quả cho thấy hiệu ứng tổng hợp (pooled
correlation) là \(r = 0{,}074\), với khoảng tin cậy 95\% dương và có ý
nghĩa thống kê. Giá trị \(r\) này khẳng định bức tranh tổng quát từ các
các phân tích tổng hợp trước đây của Bausch và Krist (2007), Kirca et
al.~(2012), Marano et al.~(2016) và Arte và Larimo (2022): mối quan hệ
tổng hợp giữa quốc tế hóa và hiệu quả hoạt động là dương nhưng có biên
độ khiêm tốn.

Đặc biệt đáng lưu ý, kết quả P6 hoàn toàn nhất quán với cơ sở phân tích
tổng hợp đã công bố tại hội nghị ICBEF 2024, với \(k = 113\) nghiên cứu
và \(r = 0{,}07\), sự hội tụ này tạo ra bằng chứng chéo vững chắc khi mở
rộng mẫu gộp lên gần gấp đôi số nghiên cứu.

\hypertarget{phuxe2n-tuxedch-dux1ecb-biux1ec7t-ba-tux1ea7ng}{%
\subsection{4.2.2 Phân tích dị biệt ba
tầng}\label{phuxe2n-tuxedch-dux1ecb-biux1ec7t-ba-tux1ea7ng}}

Kết quả quan trọng nhất của phân tích MARA không phải là giá trị điểm
của \(r\) mà là cấu trúc của dị biệt (\(I^2\)). Tổng \(I^2 = 62{,}4\%\),
cho thấy mức dị biệt đáng kể trong tập hợp 238 nghiên cứu không phải do
sai số lấy mẫu. Quan trọng hơn, khi phân tách theo cấu trúc ba tầng:

\begin{itemize}
\tightlist
\item
  \textbf{Tầng 2} (between-effect within-study, tức dị biệt giữa các
  effect-size trong cùng một nghiên cứu): chiếm \(54{,}1\%\) tổng phương
  sai, đây là nguồn dị biệt chủ đạo.
\item
  \textbf{Tầng 3} (between-study, tức dị biệt giữa các nghiên cứu): chỉ
  chiếm \(8{,}4\%\).
\end{itemize}

Phát hiện này có hàm ý lý thuyết sâu sắc: dị biệt trong tài liệu về quan
hệ quốc tế hóa và hiệu quả xuất phát chủ yếu từ sự biến thiên
\textbf{bên trong các nghiên cứu} (ví dụ: cùng một mẫu doanh nghiệp
nhưng dùng các thước đo khác nhau cho cùng cấu trúc cho ra kết quả khác
nhau), chứ không phải giữa các nghiên cứu với nhau. Điều này gợi ý rằng
sự không nhất quán trong tài liệu phần lớn là vấn đề \textbf{đo lường và
bối cảnh}, không phải là bằng chứng về sự thiếu nhất quán thực sự trong
quan hệ nhân quả giữa quốc tế hóa và hiệu quả.

\hypertarget{ux111iux1ec1u-tiux1ebft-bux1edfi-chux1ebf-ux111ux1ed9-thux1ec3-chux1ebf-icrv-ux111iux1ec1u-tiux1ebft}{%
\subsection{4.2.3 Điều tiết bởi chế độ thể chế: ICRV điều
tiết}\label{ux111iux1ec1u-tiux1ebft-bux1edfi-chux1ebf-ux111ux1ed9-thux1ec3-chux1ebf-icrv-ux111iux1ec1u-tiux1ebft}}

Kiểm định điều tiết bằng nhóm chế độ thể chế ICRV cho kết quả có ý nghĩa
thống kê: \(Q_M = 17{,}35\), \(df = 4\), \(p = {,}002\). Điều này xác
nhận rằng chế độ thể chế, được phân loại thành các nhóm từ tiên tiến đổi
mới (Nhóm I) đến nhóm biên thể chế yếu nhất, là biến điều tiết có ý
nghĩa thực sự cho mối quan hệ tổng hợp giữa quốc tế hóa và hiệu quả
trong tài liệu. Kiểm định \(Q_M\) vượt ngưỡng tới hạn (\(p = {,}002\)),
hàm ý rằng sự khác biệt giữa các nhóm ICRV không thể giải thích bằng sai
số lấy mẫu.

Tuy điều tiết có ý nghĩa thống kê, mô thức cỡ ảnh hưởng theo nhóm
\textbf{không} tạo thành một phổ thể chế đơn điệu. Cỡ ảnh hưởng trung
bình theo nhóm là: Nhóm I tiên tiến đổi mới \(\bar{r} = 0{,}079\)
(\(k = 139\)), Nhóm III \(\bar{r} = 0{,}069\) (\(k = 91\)), Nhóm II
\(\bar{r} = 0{,}065\) (\(k = 25\)), và nhóm hỗn hợp
\(\bar{r} = 0{,}053\) (\(k = 30\)); nhóm biên thể chế yếu nhất cho giá
trị bất thường cao (\(\bar{r} = 0{,}349\), \(k = 3\)) nhưng bị chi phối
bởi một giá trị ngoại lai duy nhất (Pouresmaeili và cộng sự, 2018) nên
không đáng tin cậy. Vì vậy, kết luận đúng từ P6 là chế độ thể chế điều
tiết cỡ ảnh hưởng một cách có ý nghĩa nhưng không theo trật tự đơn điệu
rõ ràng ở cấp tổng hợp; nhóm thể chế mạnh nhất có cỡ ảnh hưởng tin cậy
cao nhất, song các nhóm còn lại không xếp hạng theo phổ thể chế kỳ vọng,
và nhóm đảo nhỏ không có đủ cỡ ảnh hưởng để phân tích riêng. Cách đọc
này nhất quán với lý thuyết thể chế ở chiều tổng quát (North, 1990;
Marano et al., 2016) nhưng tinh chỉnh kỳ vọng quan trọng: bằng chứng về
một \emph{trật tự} thể chế đơn điệu không hiện ra ở cỡ ảnh hưởng cấp
tổng hợp mà ở \textbf{vị trí điểm uốn} trong phân tích hồi quy cấp doanh
nghiệp của P7, nơi cấu trúc ba vùng định hình rõ (Mục 4.6.2). Sự phân
công này, P6 xác nhận \emph{sự tồn tại} của điều tiết thể chế còn P7
định vị \emph{hình thức} của nó, chính là một lý do để ưu tiên bằng
chứng cấp doanh nghiệp trên khung phân loại thể chế nhất quán.
\emph{(Lưu ý phân loại: P6, phân tích tổng hợp cấp nghiên cứu, dùng hệ
ICRV năm chế độ rút gọn {[}I tiên tiến đổi mới, II, III, FR biên thể chế
yếu, MX hỗn hợp{]}, khác với hệ sáu nhóm cấp doanh nghiệp dùng ở Mục 4.1
và Mục 4.6; nhóm biên thể chế yếu nhất của P6 tương ứng Nhóm V và VI
trong hệ sáu nhóm.)}

\hypertarget{cuxe1c-biux1ebfn-ux111iux1ec1u-tiux1ebft-nux103ng-lux1ef1c-sux1ed1-pha-sux1ed1-huxf3a-vuxe0-thux1b0ux1edbc-ux111o}{%
\subsection{4.2.3.1 Các biến điều tiết năng lực số, pha số hóa và thước
đo}\label{cuxe1c-biux1ebfn-ux111iux1ec1u-tiux1ebft-nux103ng-lux1ef1c-sux1ed1-pha-sux1ed1-huxf3a-vuxe0-thux1b0ux1edbc-ux111o}}

Ngoài chế độ thể chế, P6 kiểm định sáu biến điều tiết khác, ba biến mới
của luận án và bốn biến chuẩn của tài liệu, qua đó định vị rõ điều tiết
nào là thực chất. Trái với kỳ vọng, chấp nhận số ngữ cảnh ở cấp tổng hợp
\textbf{không} điều tiết đáng kể cỡ ảnh hưởng (\(Q_M = 1{,}23\),
\(df = 2\), \(p = {,}541\); hệ số hồi quy phân tích tổng hợp
\(\beta = +0{,}003\), \(p = {,}744\)), nghĩa là phần thưởng số không
hiện ra như một điều tiết phổ quát ở dữ liệu thứ cấp, nhất quán với phát
hiện cấp doanh nghiệp rằng vai trò uốn đường cong của DAI phụ thuộc bối
cảnh chứ không phổ quát. Tương tự, pha phát triển số hóa không điều tiết
đáng kể (\(Q_M = 0{,}56\), \(df = 2\), \(p = {,}755\)), với cỡ ảnh hưởng
gần như đồng đều qua các pha (trước, trong và sau làn sóng số hóa).

Hai biến điều tiết về thước đo cho thấy lựa chọn vận hành hóa có ảnh
hưởng đáng kể đến cỡ ảnh hưởng quan sát, củng cố cách đọc rằng phần lớn
dị biệt nằm ở tầng nội bộ nghiên cứu. Theo cách đo mức độ quốc tế hóa,
cỡ ảnh hưởng dao động từ 0,059 khi đo bằng cường độ xuất khẩu trên doanh
thu (FSTS, \(K = 137\)) đến 0,097 khi đo bằng phạm vi địa lý
(\(K = 50\)) và 0,111 khi đo bằng chỉ số tổng hợp (\(K = 31\)). Theo
cách đo hiệu quả, cỡ ảnh hưởng theo thước đo kế toán là 0,078
(\(K = 246\)) nhưng giảm xuống gần không và đổi dấu nhẹ khi đo bằng năng
suất lao động (\(-0{,}019\), \(K = 12\)), một mẫu rất nhỏ. Mô thức này
có hàm ý phương pháp trực tiếp: việc luận án dùng nhất quán cường độ
xuất khẩu làm thước đo quốc tế hóa và năng suất lao động làm thước đo
hiệu quả ở cấp doanh nghiệp loại bỏ chính nguồn dị biệt do thước đo mà
phân tích tổng hợp cấp tổng hợp không thể kiểm soát, qua đó cho ước
lượng điều tiết thể chế sạch hơn.

\hypertarget{kiux1ec3m-ux111ux1ecbnh-lux1ec7ch-lux1ea1c-cuxf4ng-bux1ed1}{%
\subsection{4.2.4 Kiểm định lệch lạc công
bố}\label{kiux1ec3m-ux111ux1ecbnh-lux1ec7ch-lux1ea1c-cuxf4ng-bux1ed1}}

Phân tích lệch lạc công bố gồm ba kiểm định bổ sung nhau. \textbf{Kiểm
định Egger} cho kết quả \(b = 0{,}487\) (\(p = {,}052\)), sát ngưỡng
\(\alpha = 0{,}05\) nhưng không đạt ý nghĩa thống kê theo tiêu chuẩn
thông thường. \textbf{Kiểm định Begg--Mazumdar} (Begg \& Mazumdar, 1994)
cho kết quả \(\tau = -0{,}132\) (\(p = {,}001\)), có ý nghĩa thống kê rõ
ràng, cho thấy các nghiên cứu với sai số chuẩn lớn hơn (mẫu nhỏ hơn) có
xu hướng báo cáo cỡ ảnh hưởng thấp hơn, phù hợp với lệch lạc công bố ưu
tiên kết quả dương tính. \textbf{Phân tích cắt và điền} điều chỉnh ước
lượng từ \(\bar{r} = 0{,}074\) xuống \(\bar{r} = 0{,}035\) (với
\(k = 57\) nghiên cứu được điền khuyết), mức điều chỉnh đáng kể nhưng
không đổi chiều hướng tích cực của quan hệ tổng hợp. \textbf{Số an toàn
\(N\) = 44.782} vượt xa ngưỡng tới hạn \(5k + 10 = 1.200\), khẳng định
lệch lạc công bố không thể bác bỏ hoàn toàn kết quả tổng hợp. Tổng hợp
lại, bằng chứng lệch lạc công bố được xác nhận qua kiểm định
Begg--Mazumdar có ý nghĩa (\(p = {,}001\)): cỡ ảnh hưởng gộp
\(\bar{r} = 0{,}074\) có thể bị thổi phồng so với hiệu ứng tổng thể
thực, với ước lượng bảo thủ hơn là \(\bar{r} = 0{,}035\).

Phát hiện lệch lạc công bố này có một hàm ý lý thuyết quan trọng vượt ra
ngoài một cảnh báo phương pháp đơn thuần. Trong phân tách dị biệt ba
tầng, phần lớn dị biệt nằm ở \textbf{tầng nội bộ nghiên cứu}
(\(I^2_{\text{within}} = 54{,}1\%\)) hơn là tầng giữa các nghiên cứu
(\(I^2_{\text{between}} = 8{,}4\%\)), nghĩa là biến thiên hệ số chủ yếu
đến từ cách mỗi nghiên cứu vận hành hóa quốc tế hóa, lựa chọn thước đo
hiệu quả và thiết lập biến kiểm soát, chứ không phải từ khác biệt bối
cảnh quốc gia. Kết hợp với bằng chứng lệch lạc công bố, cấu trúc dị biệt
này gợi ý một cách diễn giải lại ``câu đố dị biệt'' trong tài liệu: phần
lớn phương sai chưa giải thích được có thể phản ánh lựa chọn ở phía công
bố và thiết lập mô hình nhiều hơn là các điều kiện bối cảnh thể chế hay
số hóa. Phát hiện này không phủ định vai trò điều tiết của thể chế, vốn
được xác nhận độc lập qua kiểm định \(Q_M\) giữa các nhóm ICRV, nhưng nó
định vị lại kỳ vọng: phân tích tổng hợp cấp tổng hợp khó nắm bắt được
điều tiết thể chế vì các nghiên cứu sơ cấp hiếm khi báo cáo đủ chi tiết
theo chế độ thể chế, và vì vậy bằng chứng cấp doanh nghiệp trực tiếp
trên một khung phân loại thể chế nhất quán, như luận án thực hiện, là
con đường hữu hiệu hơn để bóc tách dị biệt so với việc tích lũy thêm
phân tích tổng hợp trên dữ liệu thứ cấp không đồng nhất.

\hypertarget{phuxe2n-ruxe3-phux1b0ux1a1ng-sai-ba-tux1ea7ng-vuxe0-huxe0m-uxfd-vux1ec1-nguux1ed3n-dux1ecb-biux1ec7t}{%
\subsection{4.2.5 Phân rã phương sai ba tầng và hàm ý về nguồn dị
biệt}\label{phuxe2n-ruxe3-phux1b0ux1a1ng-sai-ba-tux1ea7ng-vuxe0-huxe0m-uxfd-vux1ec1-nguux1ed3n-dux1ecb-biux1ec7t}}

Cấu trúc ba tầng của mô hình MARA cho phép quy phương sai hệ thống về
đúng nguồn gốc của nó thay vì gộp chung như mô hình hiệu ứng ngẫu nhiên
một tầng. Hai thành phần phương sai ước lượng bằng REML là phương sai
trong nghiên cứu \(\hat{\sigma}^2_{(2)} = 0{,}00878\) và phương sai giữa
các nghiên cứu \(\hat{\sigma}^2_{(3)} = 0{,}00136\). Phương sai trong
nghiên cứu lớn gấp khoảng sáu lần phương sai giữa các nghiên cứu, một tỷ
lệ trực tiếp chuyển thành phần dị biệt \(I^2_{(2)} = 54{,}1\%\) áp đảo
phần \(I^2_{(3)} = 8{,}4\%\) đã trình bày ở Mục 4.2.2. Kiểm định Cochran
tổng cho \(Q = 1.909{,}42\) (\(df = 287\), \(p < {,}001\)), bác bỏ giả
thuyết đồng nhất một cách rõ ràng và do đó biện minh cho lựa chọn mô
hình hiệu ứng ngẫu nhiên ba tầng.

Hàm ý phương pháp luận của phân rã này khác biệt căn bản so với một báo
cáo \(I^2\) tổng đơn thuần. Vì biến thiên hệ số chủ yếu phát sinh từ
những lựa chọn vận hành hóa bên trong cùng một bài báo, cụ thể là cách
đo lường cường độ quốc tế hóa (tỷ lệ doanh thu xuất khẩu trên tổng doanh
thu, chỉ số entropy, số thị trường nước ngoài), cách chọn thước đo hiệu
quả (dựa trên kế toán, dựa trên thị trường, dựa trên năng suất) và thiết
lập biến kiểm soát, nên việc bổ sung thêm các nghiên cứu sơ cấp mới khó
có khả năng thu hẹp dị biệt nếu các nghiên cứu đó tiếp tục báo cáo nhiều
hệ số không đồng nhất cho cùng một mẫu. Mô hình hai tầng quy ước (ước
lượng DerSimonian--Laird) cho cùng giá trị điểm \(r = 0{,}074\) với
khoảng tin cậy \([0{,}061;\ 0{,}087]\), xác nhận rằng việc bỏ qua cấu
trúc lồng ghép không làm chệch ước lượng gộp mà chỉ làm mất thông tin về
nguồn dị biệt. Đây là lý do trung tâm để luận án ưu tiên bằng chứng cấp
doanh nghiệp trên một khung phân loại thể chế nhất quán hơn là tiếp tục
tích lũy phân tích tổng hợp trên dữ liệu thứ cấp không đồng nhất.

\hypertarget{phux1ed5-hiux1ec7u-ux1ee9ng-theo-nhuxf3m-thux1ec3-chux1ebf-vuxe0-giux1edbi-hux1ea1n-cux1ee7a-bux1eb1ng-chux1ee9ng-ux111ux1ecbnh-hux1b0ux1edbng}{%
\subsection{4.2.6 Phổ hiệu ứng theo nhóm thể chế và giới hạn của bằng
chứng định
hướng}\label{phux1ed5-hiux1ec7u-ux1ee9ng-theo-nhuxf3m-thux1ec3-chux1ebf-vuxe0-giux1edbi-hux1ea1n-cux1ee7a-bux1eb1ng-chux1ee9ng-ux111ux1ecbnh-hux1b0ux1edbng}}

Bảng hệ số gộp theo từng nhóm ICRV cho thấy một bức tranh tinh tế hơn so
với kết luận ``có điều tiết'' đơn thuần. Nhóm I (Advanced đổi mới sáng
tạo) có hệ số gộp lớn nhất trong các nhóm đủ dữ liệu,
\(\bar{r} = 0{,}079\) (khoảng tin cậy \([0{,}058;\ 0{,}099]\),
\(k = 139\) hiệu ứng); nhóm II (Trung bình cao) có \(\bar{r} = 0{,}065\)
(\(k = 25\)); nhóm III (Đang nổi) có \(\bar{r} = 0{,}069\) (\(k = 91\));
và nhóm đa bối cảnh có \(\bar{r} = 0{,}053\) (\(k = 30\)). Bốn nhóm này
đều cho hệ số dương có ý nghĩa và xếp theo chiều phù hợp với kỳ vọng lý
thuyết, song khác biệt giữa chúng không đạt ý nghĩa thống kê trong các
so sánh cặp đôi: Nhóm I so với Nhóm III cho \(b = -0{,}011\)
(\(p = {,}502\)), và Nhóm I so với nhóm đa bối cảnh cho \(b = -0{,}026\)
(\(p = {,}269\)).

Do đó cần phân biệt rạch ròi giữa kết quả của kiểm định tổng thể và kết
quả của các so sánh định hướng. Kiểm định \(Q_M\) giữa các nhóm có ý
nghĩa (\(p = {,}002\)) khẳng định rằng hệ số quốc tế hóa và hiệu quả
biến thiên có hệ thống giữa các chế độ thể chế, nhưng ý nghĩa này được
dẫn dắt gần như toàn bộ bởi một nhóm cận biên: nhóm Frontier chỉ gồm ba
nghiên cứu lại cho hệ số gộp dị thường \(\bar{r} = 0{,}349\)
(\(k = 3\)), và tương phản duy nhất có ý nghĩa là Nhóm I so với Frontier
(\(b = +0{,}285\), \(p < {,}001\)). Giá trị Frontier này bị chi phối bởi
một nghiên cứu ngoại lai duy nhất (Pouresmaeili và cộng sự, 2018, với
\(r = 0{,}69\) trên 226 doanh nghiệp) và vì vậy không thể coi là một ước
lượng đáng tin cậy cho chế độ thể chế yếu nhất. Hệ quả là phổ thể chế
định hướng không được xác nhận ở cấp phân tích tổng hợp với cỡ mẫu hiện
tại; sự khác biệt giữa Nhóm I và Nhóm III chỉ là \(0{,}010\) và không có
ý nghĩa thống kê. Phát hiện này định vị lại kỳ vọng một cách quan trọng:
bằng chứng về điều tiết thể chế tồn tại, nhưng việc bóc tách phổ thể chế
cụ thể đòi hỏi bằng chứng cấp doanh nghiệp trên một khung phân loại nhất
quán, đúng như con đường mà các nghiên cứu thành phần P3--P8 và mô hình
đa quốc gia P7 của luận án theo đuổi.

Hai biến điều tiết bối cảnh khác được kiểm định trong cùng khung MARA
nhưng không đạt ý nghĩa thống kê. Chấp nhận số cấp quốc gia
(country-level Digital Adoption Index --- cDAI) cho \(Q_M = 1{,}23\)
(\(df = 2\), \(p = {,}541\)), với hệ số phân tích tổng hợp-hồi quy liên
tục \(\beta_{\text{cDAI}} = +0{,}003\) (\(p = {,}744\)) và thứ tự ba
phân nhóm không đơn điệu (Thấp \(0{,}075\), Trung bình \(0{,}063\), Cao
\(0{,}091\)). Chu kỳ nghịch lý số (Digital Paradox Lifecycle --- DPL)
cho \(Q_M = 0{,}56\) (\(df = 2\), \(p = {,}755\)), với thứ tự ba giai
đoạn ngược với dự đoán lý thuyết. Việc chỉ có một trong ba biến điều
tiết bối cảnh đạt ý nghĩa, và biến đó lại bị chi phối bởi một ô dữ liệu
cận biên, củng cố cách diễn giải rằng dị biệt quốc tế hóa và hiệu quả ở
cấp tổng hợp khó được giải thích bằng các biến bối cảnh quốc gia khi các
nghiên cứu sơ cấp hiếm khi báo cáo đủ chi tiết theo chế độ thể chế.

\hypertarget{tuxednh-vux1eefng-cux1ee7a-hiux1ec7u-ux1ee9ng-gux1ed9p-vuxe0-phux1ea1m-vi-cux1ee7a-bux1eb1ng-chux1ee9ng-lux1ec7ch-lux1ea1c-cuxf4ng-bux1ed1}{%
\subsection{4.2.7 Tính vững của hiệu ứng gộp và phạm vi của bằng chứng
lệch lạc công
bố}\label{tuxednh-vux1eefng-cux1ee7a-hiux1ec7u-ux1ee9ng-gux1ed9p-vuxe0-phux1ea1m-vi-cux1ee7a-bux1eb1ng-chux1ee9ng-lux1ec7ch-lux1ea1c-cuxf4ng-bux1ed1}}

Hiệu ứng gộp \(r = 0{,}074\) giữ ổn định qua toàn bộ các kiểm định độ
vững có sẵn. Khi chỉ giữ các hệ số tương quan được báo cáo trực tiếp, hệ
số gộp là \(0{,}077\) (\(k = 241\)); khi loại các hiệu ứng từ mẫu rất
nhỏ (\(n < 30\)), hệ số là \(0{,}073\) (\(k = 286\)); khi chỉ giữ thước
đo hiệu quả dựa trên kế toán, hệ số là \(0{,}075\) (\(k = 247\)). Kiểm
định bảo thủ nhất là giới hạn ở các nghiên cứu chỉ dùng tỷ lệ doanh thu
xuất khẩu trên tổng doanh thu làm thước đo quốc tế hóa, cho hệ số suy
giảm còn \(\bar{r} = 0{,}060\) (\(k = 138\)), vẫn dương và có ý nghĩa,
gợi ý rằng tính không đồng nhất của thước đo quốc tế hóa làm phóng đại
một phần hiệu ứng gộp khi gộp chung các thước đo rộng hơn. Phân tích
loại bỏ từng nghiên cứu (leave-one-out) cho dải hệ số
\([0{,}071;\ 0{,}076]\) và không một nghiên cứu nào trong 288 hiệu ứng
làm đổi chiều kết quả, xác nhận rằng hiệu ứng dương khiêm tốn không phải
là sản phẩm của bất kỳ quan sát ảnh hưởng đơn lẻ nào.

Về lệch lạc công bố, cần nhấn mạnh sự không thống nhất giữa hai kiểm
định bất đối xứng phễu đã nêu ở Mục 4.2.4: kiểm định Egger không vượt
ngưỡng quy ước (\(b = 0{,}487\), \(p = {,}052\)) trong khi kiểm định thứ
hạng Begg--Mazumdar lại có ý nghĩa (\(\tau = -0{,}132\),
\(p = {,}001\)), nên bằng chứng bất đối xứng dựa trên Egger đơn lẻ không
được khẳng định. Quy trình cắt-và-điền (trim-and-fill) gán thêm 57
nghiên cứu ở phía trái của phễu và đưa hệ số điều chỉnh xuống
\(\bar{r} = 0{,}035\) (khoảng tin cậy \([0{,}018;\ 0{,}051]\)); đây là
cận dưới bảo thủ, hàm ý hiệu ứng dân số thực có thể gần
\(r \approx 0{,}035\) hơn là các con số trong dải \(0{,}07\)--\(0{,}10\)
thường được trích dẫn. Cùng lúc, chỉ số an toàn fail-safe \(N = 44.782\)
vượt xa ngưỡng tới hạn \(5k + 10 = 1.200\), hàm ý rằng lệch lạc công bố
tuy hiện diện nhưng không thể đảo ngược hoàn toàn kết luận về dấu của
quan hệ. Tổng hợp lại, các kiểm định độ vững và lệch lạc công bố hội tụ
về một thông điệp duy nhất: quan hệ quốc tế hóa và hiệu quả ở cấp tổng
hợp là dương, ổn định về dấu, nhưng nhỏ hơn nhiều so với ấn tượng từ tài
liệu đã công bố.

Một hàm ý phương pháp luận trực tiếp nối kết phân tích tổng hợp P6 với
chiến lược thiết kế cấp doanh nghiệp của luận án. Vì hai biến điều tiết
thước đo cho thấy cỡ ảnh hưởng quan sát phụ thuộc mạnh vào lựa chọn vận
hành hóa, dao động từ \(0{,}059\) khi đo quốc tế hóa bằng cường độ xuất
khẩu trên doanh thu (\(K = 137\)) tới \(0{,}097\) khi đo bằng phạm vi
địa lý (\(K = 50\)) và \(0{,}111\) khi đo bằng chỉ số tổng hợp
(\(K = 31\)), và vì cỡ ảnh hưởng theo thước đo kế toán (\(0{,}078\),
\(K = 246\)) đổi dấu nhẹ khi chuyển sang năng suất lao động
(\(-0{,}019\), \(K = 12\)), một phần lớn dị biệt tầng nội bộ nghiên cứu
(\(I^2_{\text{within}} = 54{,}1\%\)) phản ánh chính sự không đồng nhất
của thước đo. Hệ quả thiết kế là rõ ràng: việc luận án dùng nhất quán
cường độ xuất khẩu (FSTS) làm thước đo quốc tế hóa và năng suất lao động
làm thước đo hiệu quả ở cấp doanh nghiệp loại bỏ chính nguồn dị biệt do
thước đo mà phân tích tổng hợp cấp tổng hợp không thể kiểm soát, cho ước
lượng điều tiết thể chế sạch hơn và biến bằng chứng cấp doanh nghiệp
trên khung phân loại thể chế nhất quán thành con đường hữu hiệu hơn để
bóc tách dị biệt so với việc tích lũy thêm phân tích tổng hợp trên dữ
liệu thứ cấp không đồng nhất.

\begin{center}\rule{0.5\linewidth}{0.5pt}\end{center}

\hypertarget{kux1ebft-quux1ea3-bux1ed1i-cux1ea3nh-thux1ec3-chux1ebf-tiuxean-tiux1ebfn-ux111ux1ed5i-mux1edbi-singapore-p4}{%
\section{4.3 Kết quả bối cảnh thể chế tiên tiến đổi mới: Singapore
(P4)}\label{kux1ebft-quux1ea3-bux1ed1i-cux1ea3nh-thux1ec3-chux1ebf-tiuxean-tiux1ebfn-ux111ux1ed5i-mux1edbi-singapore-p4}}

\hypertarget{xuxe1c-nhux1eadn-quan-hux1ec7-phi-tuyux1ebfn-chux1eef-u-ngux1b0ux1ee3c}{%
\subsection{4.3.1 Xác nhận quan hệ phi tuyến chữ U
ngược}\label{xuxe1c-nhux1eadn-quan-hux1ec7-phi-tuyux1ebfn-chux1eef-u-ngux1b0ux1ee3c}}

Phân tích trên mẫu doanh nghiệp Singapore từ dữ liệu WBES với
\(N = 623\) cơ sở (trong đó 84 cơ sở xuất khẩu dùng cho kiểm định mẫu
phụ) cho thấy quan hệ giữa FSTS và hiệu quả gần tuyến tính dương; thành
phần bậc hai chỉ gợi ý đường cong rất thoải với điểm uốn ngoài miền dữ
liệu. Mô hình bậc hai với biến FSTS và FSTS\(^2\) cho hệ số bậc hai âm,
phù hợp với dạng chữ U ngược như kỳ vọng từ H1 trong khung lý thuyết
CDCM.

Điểm điểm uốn ước lượng có giá trị trung tâm \textbf{TP ≈ 88,6\% FSTS},
nhưng khoảng tin cậy bootstrap rất rộng và \textbf{vượt ngưỡng khả thi
của FSTS} (CI: {[}53\%, 253\%{]}; trong khi FSTS theo định nghĩa không
thể vượt 100\%). Vì cận trên của khoảng tin cậy nằm ngoài miền giá trị
khả thi, điểm uốn \textbf{không được định vị một cách tin cậy trong miền
dữ liệu quan sát}: trên thực tế, trong khoảng FSTS quan sát được
(0--100\%), quan hệ FSTS--hiệu quả hoạt động kinh doanh của doanh nghiệp
Singapore là \textbf{tăng gần như đơn điệu}, và phần suy giảm bên phải
của đường chữ U ngược nằm ngoài (hoặc ở rìa) miền dữ liệu nên không được
kiểm chứng thực nghiệm. Do đó, kết quả Singapore nên được diễn giải thận
trọng là \emph{``chưa quan sát thấy ngưỡng bất lợi của quốc tế hóa trong
miền dữ liệu''} thay vì khẳng định một điểm uốn cụ thể ở 88,6\%; điều
này nhất quán với kỳ vọng lý thuyết rằng ở bối cảnh thể chế mạnh (Nhóm
I), dư địa quốc tế hóa có lợi là rất rộng (xem Hình 4.8).

\hypertarget{ux111iux1ec1u-tiux1ebft-bux1edfi-chux1ec9-sux1ed1-chux1ea5p-nhux1eadn-sux1ed1-dai}{%
\subsection{4.3.2 Điều tiết bởi chỉ số chấp nhận số
(DAI)}\label{ux111iux1ec1u-tiux1ebft-bux1edfi-chux1ec9-sux1ed1-chux1ea5p-nhux1eadn-sux1ed1-dai}}

Mô hình M5 bổ sung tương tác \(\mathrm{FSTS} \times \mathrm{DAI}\) (bao
gồm cả Tầng 1 và Tầng 2 của chỉ số DAI) cho thấy hiệu ứng khuếch đại
(amplification) tại mức FSTS cao: hệ số tương tác
\(\mathrm{FSTS}^2 \times \mathrm{DAI} = +3{,}119\) và có ý nghĩa thống
kê (\(p = {,}005\)). Cụ thể, ở mức độ quốc tế hóa cao, doanh nghiệp
Singapore có năng lực chấp nhận số (digital adoption) cao hơn duy trì
hiệu quả tốt hơn so với doanh nghiệp có mức DAI thấp hơn, cho thấy DAI
đóng vai trò nguồn lực bổ trợ (situational resource) làm chậm đà suy
giảm phía bên phải của đường phi tuyến.

\hypertarget{giux1edbi-hux1ea1n-suy-luux1eadn-vuxe0-cux1ea3nh-buxe1o-thux1ed1ng-kuxea}{%
\subsection{4.3.3 Giới hạn suy luận và cảnh báo thống
kê}\label{giux1edbi-hux1ea1n-suy-luux1eadn-vuxe0-cux1ea3nh-buxe1o-thux1ed1ng-kuxea}}

Cần nhấn mạnh rằng kết quả Singapore là \textbf{ranh giới suy luận}
(inferential bounds) do hạn chế về quy mô mẫu: \(N = 623\) quan sát toàn
mẫu và chỉ 84 cơ sở xuất khẩu cho các kiểm định mẫu phụ, dẫn đến power
thống kê ước tính chỉ đạt khoảng 16\%. Luận án diễn giải kết quả
Singapore như \emph{bằng chứng phù hợp nhưng chưa kết luận} (consistent
but underpowered evidence), không bác bỏ giả thuyết nhưng cũng không
tuyên bố xác nhận dứt khoát.

\hypertarget{phuxe2n-phux1ed1i-mux1eabu-ux111ux1ecbnh-vux1ecb-ux111iux1ec3m-uux1ed1n-bux1eb1ng-bootstrap-vuxe0-ux111a-cux1ed9ng-tuyux1ebfn}{%
\subsection{4.3.4 Phân phối mẫu, định vị điểm uốn bằng bootstrap và đa
cộng
tuyến}\label{phuxe2n-phux1ed1i-mux1eabu-ux111ux1ecbnh-vux1ecb-ux111iux1ec3m-uux1ed1n-bux1eb1ng-bootstrap-vuxe0-ux111a-cux1ed9ng-tuyux1ebfn}}

Bức tranh phân phối của mẫu Singapore củng cố cách diễn giải thận trọng
nói trên. Cường độ xuất khẩu trung bình chỉ đạt
\(\overline{\text{FSTS}} = 0{,}045\) với độ lệch chuẩn \(0{,}144\),
trong đó khoảng \(82\%\) doanh nghiệp báo cáo không xuất khẩu và chỉ
\(3\%\) vượt ngưỡng \(50\%\) doanh thu nước ngoài. Chính sự dồn nén ở
đầu thấp của miền FSTS, kết hợp với mức bão hòa số rất cao (khoảng
\(67\%\) doanh nghiệp có hiện diện website), làm cho nhánh suy giảm bên
phải của đường cong khó được nhận diện thống kê trong cỡ mẫu 623 cơ sở.
Kiểm định Lind--Mehlum do đó không bác bỏ tính đơn điệu ở mức ý nghĩa
thông thường (\(p = {,}303\)), một kết quả không có ý nghĩa nhưng giàu
thông tin (informative null), nhất quán với giả thuyết bão hòa số thay
vì một thất bại của khung chữ U ngược. Về kỹ thuật, điểm uốn trung tâm
\(88{,}6\%\) được tính từ công thức \(-\hat{\beta}_1/(2\hat{\beta}_2)\)
trên FSTS đã trừ trung bình (cho giá trị \(\approx 84{,}1\%\)) rồi cộng
lại trung bình mẫu \(0{,}045\); bootstrap 5.000 lần tái lập tái hiện
dạng chữ U ngược ở \(96{,}3\%\) số lần lặp nhưng khoảng tin cậy phân vị
\(95\%\) của điểm uốn trải rộng \([53\%; 253\%]\) (trung vị \(80\%\),
khoảng tứ phân vị \([68\%; 102\%]\)), khẳng định vị trí điểm uốn chỉ
được định vị lỏng lẻo. Đa cộng tuyến giữa số hạng bậc nhất và bậc hai
không phải là vấn đề thực chất: hệ số phóng đại phương sai (variance
inflation factor --- VIF) duy trì dưới 3 ở mọi mô hình cuối.

\hypertarget{tuxe1ch-bux1ea1ch-hai-cux1ea5u-truxfac-tci-vuxe0-dai-hiux1ec7u-ux1ee9ng-trux1ef1c-tiux1ebfp-vuxe0-tiux1ebfn-triux1ec3n-ux111ux1ed9-phuxf9-hux1ee3p-muxf4-huxecnh}{%
\subsection{4.3.5 Tách bạch hai cấu trúc TCI và DAI: hiệu ứng trực tiếp
và tiến triển độ phù hợp mô
hình}\label{tuxe1ch-bux1ea1ch-hai-cux1ea5u-truxfac-tci-vuxe0-dai-hiux1ec7u-ux1ee9ng-trux1ef1c-tiux1ebfp-vuxe0-tiux1ebfn-triux1ec3n-ux111ux1ed9-phuxf9-hux1ee3p-muxf4-huxecnh}}

Bằng chứng Singapore cho phép phân tách rạch ròi hai cấu trúc năng lực
thường bị gộp chung trong tài liệu số hóa. Năng lực công nghệ thể hiện
một hiệu ứng trực tiếp dương và bền vững trên năng suất lao động: ở mô
hình hiệu ứng trực tiếp (M5), \(\hat{\beta}(\text{TCI}) = +0{,}168\)
(\(SE = 0{,}040\); \(p < {,}001\)), đi kèm cải thiện độ phù hợp mô hình
khi \(R^2\) tăng từ \(0{,}178\) ở M2 lên \(0{,}199\) ở M5. Trong mô hình
điều tiết TCI bổ sung (M3), hệ số trực tiếp vẫn dương và có ý nghĩa
(\(\hat{\beta} = +0{,}188\), \(p < {,}001\)) nhưng các số hạng tương tác
với \(\text{FSTS}\) và \(\text{FSTS}^2\) không có ý nghĩa thống kê khi
kiểm định chung, củng cố cách đọc rằng năng lực công nghệ nâng mặt bằng
năng suất thay vì uốn hình dạng đường cong.

Ngược lại, năng lực chấp nhận số nền tảng thể hiện hành vi mang tính bối
cảnh. Ở mô hình trực tiếp (M6), \(\hat{\beta}(\text{DAI}) = +0{,}104\)
(\(SE = 0{,}038\); \(p = {,}007\)), nhưng khi đưa đồng thời TCI vào (M7)
hệ số DAI suy giảm còn \(+0{,}077\) và chỉ còn ở biên ý nghĩa
(\(p = {,}048\)), cho thấy một phần phương sai năng suất mà DAI nắm bắt
đã trùng lặp với TCI. Sự suy giảm này quan trọng về nội dung vì nó bác
bỏ cách hiểu DAI như một phần thưởng năng suất phổ quát cho mọi doanh
nghiệp.

\hypertarget{cux1a1-chux1ebf-luxe1-chux1eafn-sux1ed1-tux1b0ux1a1ng-tuxe1c-tux1ea7ng-1tux1ea7ng-2-vuxe0-hiux1ec7u-ux1ee9ng-biuxean-theo-cux1b0ux1eddng-ux111ux1ed9-xuux1ea5t-khux1ea9u}{%
\subsection{4.3.6 Cơ chế lá chắn số: tương tác Tầng 1--Tầng 2 và hiệu
ứng biên theo cường độ xuất
khẩu}\label{cux1a1-chux1ebf-luxe1-chux1eafn-sux1ed1-tux1b0ux1a1ng-tuxe1c-tux1ea7ng-1tux1ea7ng-2-vuxe0-hiux1ec7u-ux1ee9ng-biuxean-theo-cux1b0ux1eddng-ux111ux1ed9-xuux1ea5t-khux1ea9u}}

Bằng chứng mạnh nhất về vai trò của DAI xuất hiện ở các mô hình điều
tiết. Trong mô hình đầy đủ M8, số hạng trực tiếp DAI nhỏ và không có ý
nghĩa (\(\hat{\beta} = +0{,}019\); \(p = {,}705\)), tương tác bậc nhất
\(\text{FSTS} \times \text{DAI}\) âm và ở biên
(\(\hat{\beta} = -1{,}177\); \(p = {,}083\)), trong khi tương tác bậc
hai \(\text{FSTS}^2 \times \text{DAI}\) dương và có ý nghĩa thống kê
(\(\hat{\beta} = +3{,}119\); \(SE = 1{,}124\); \(p = {,}005\)). M8 cũng
đạt sức giải thích cao nhất trong các mô hình chính, với
\(R^2 = 0{,}211\) và \(R^2\) hiệu chỉnh \(= 0{,}196\). Đây chính là cơ
chế lá chắn số: tương tác bậc hai dương làm tù đi nhánh suy giảm bên
phải, tức là DAI bảo vệ doanh nghiệp khỏi đà giảm hiệu quả ở mức quốc tế
hóa cao.

Bảng hiệu ứng biên của DAI theo từng mức FSTS làm rõ tính bối cảnh này.
Tại doanh nghiệp thuần nội địa (\(\text{FSTS} = 0\)), hiệu ứng biên của
một độ lệch chuẩn DAI chỉ là \(+0{,}080\) (\(p = {,}045\)); trên dải
xuất khẩu thấp và trung bình, hiệu ứng không phân biệt được với 0 (ví dụ
tại \(\text{FSTS} = 20\%\): \(-0{,}087\), \(p = {,}444\)); nhưng tại
nhóm xuất khẩu cường độ cao, hiệu ứng trở nên dương mạnh và có ý nghĩa:
\(+0{,}660\) tại \(\text{FSTS} = 70\%\) (\(p = {,}025\)) và \(+1{,}818\)
tại \(\text{FSTS} = 100\%\) (\(p = {,}002\)). Mô thức này xác nhận DAI
là một nguồn lực mở rộng quy mô có điều kiện, phát huy giá trị chính khi
doanh nghiệp đối mặt với nhu cầu phối hợp xuyên biên giới dày đặc hơn.
Dấu dương của số hạng điều tiết bậc hai được giữ nguyên qua sáu mô hình
độ vững. Một kiểm định đảo chỗ chỉ báo (item-swap) củng cố ranh giới cấu
trúc giữa TCI và DAI: khi chuyển chỉ báo website (c22b) từ DAI sang TCI,
ý nghĩa chung của khối điều tiết DAI sụp đổ, với thống kê F chung giảm
từ \(4{,}56\) (\(p = {,}011\)) xuống \(1{,}88\) (\(p = {,}154\)). Ở mẫu
phụ chỉ gồm doanh nghiệp xuất khẩu (R5, \(N = 84\)), tương tác bậc hai
vẫn dương (\(\hat{\beta} = +2{,}821\)) với kiểm định F chung có ý nghĩa
(\(F = 6{,}32\); \(p = {,}003\)).

\hypertarget{ux111ux1ecbnh-lux1b0ux1ee3ng-cuxf4ng-suux1ea5t-thux1ed1ng-kuxea}{%
\subsection{4.3.7 Định lượng công suất thống
kê}\label{ux111ux1ecbnh-lux1b0ux1ee3ng-cuxf4ng-suux1ea5t-thux1ed1ng-kuxea}}

Cảnh báo về công suất thống kê cần được định lượng cụ thể. Hệ số Cohen's
\(f^2\) của hiệu ứng trực tiếp TCI (so sánh M2 với M5) xấp xỉ
\(0{,}036\), thuộc dải nhỏ-đến-trung bình; còn \(f^2\) của khối điều
tiết DAI (so sánh M7 với M8) chỉ khoảng \(0{,}018\), dưới ngưỡng hiệu
ứng nhỏ quy ước \(0{,}02\) của Cohen (1988). Với \(f^2 = 0{,}018\),
\(\alpha = {,}05\) và một bậc tự do tử số, để đạt công suất \(80\%\) cần
khoảng \(N > 430\): toàn mẫu phân tích (\(N = 617\)) đủ điều kiện, nhưng
mẫu phụ chỉ gồm doanh nghiệp xuất khẩu (\(N = 84\)) chỉ đạt công suất
khoảng \(16\%\). Do đó, các kết quả không có ý nghĩa về điều tiết DAI
trong mẫu phụ xuất khẩu không nên được hiểu là bằng chứng phản bác cơ
chế mở rộng quy mô có điều kiện, mà là hệ quả của cỡ mẫu chưa đủ để phát
hiện hiệu ứng nhỏ này.

\hypertarget{bux1ed9-hux1ec7-sux1ed1-ux111ux1ea7y-ux111ux1ee7-theo-chuux1ed7i-muxf4-huxecnh-phuxe2n-tux1ea7ng-vuxe0-tiux1ebfn-triux1ec3n-ux111ux1ed9-phuxf9-hux1ee3p}{%
\subsection{4.3.8 Bộ hệ số đầy đủ theo chuỗi mô hình phân tầng và tiến
triển độ phù
hợp}\label{bux1ed9-hux1ec7-sux1ed1-ux111ux1ea7y-ux111ux1ee7-theo-chuux1ed7i-muxf4-huxecnh-phuxe2n-tux1ea7ng-vuxe0-tiux1ebfn-triux1ec3n-ux111ux1ed9-phuxf9-hux1ee3p}}

Để minh bạch hóa cách hai cấu trúc năng lực tách dần khỏi đường cong
cường độ qua chuỗi mô hình lồng nhau, Bảng 4.5 hợp nhất bộ hệ số trọng
tâm của bảy mô hình chính trên mẫu Singapore. Hạng cường độ xuất khẩu
tuyến tính dương và có ý nghĩa cao ở mọi mô hình, dao động từ
\(+1{,}952\) (M7) đến \(+2{,}768\) (M4), trong khi hạng bậc hai chỉ đạt
ý nghĩa thống kê trong hai mô hình có khối tương tác DAI (M4:
\(-2{,}959\), \(p < {,}01\); M8: \(-2{,}543\), \(p < {,}05\)) và ở biên
trong mô hình nền M2 (\(-1{,}705\), \(p = {,}068\)).

\textbf{Bảng 4.5.} \emph{Tiến triển hệ số trọng tâm theo chuỗi mô hình
phân tầng, Singapore 2023 (biến phụ thuộc: ln năng suất lao động; sai số
chuẩn vững HC1).}

\begin{longtable}[]{@{}
  >{\raggedright\arraybackslash}p{(\columnwidth - 12\tabcolsep) * \real{0.1429}}
  >{\raggedleft\arraybackslash}p{(\columnwidth - 12\tabcolsep) * \real{0.1429}}
  >{\raggedleft\arraybackslash}p{(\columnwidth - 12\tabcolsep) * \real{0.1429}}
  >{\raggedleft\arraybackslash}p{(\columnwidth - 12\tabcolsep) * \real{0.1429}}
  >{\raggedleft\arraybackslash}p{(\columnwidth - 12\tabcolsep) * \real{0.1429}}
  >{\raggedleft\arraybackslash}p{(\columnwidth - 12\tabcolsep) * \real{0.1429}}
  >{\raggedleft\arraybackslash}p{(\columnwidth - 12\tabcolsep) * \real{0.1429}}@{}}
\toprule\noalign{}
\begin{minipage}[b]{\linewidth}\raggedright
Hệ số
\end{minipage} & \begin{minipage}[b]{\linewidth}\raggedleft
M2 (U ngược)
\end{minipage} & \begin{minipage}[b]{\linewidth}\raggedleft
M5 (+TCI)
\end{minipage} & \begin{minipage}[b]{\linewidth}\raggedleft
M6 (+DAI)
\end{minipage} & \begin{minipage}[b]{\linewidth}\raggedleft
M7 (T+D)
\end{minipage} & \begin{minipage}[b]{\linewidth}\raggedleft
M4 (DAI×)
\end{minipage} & \begin{minipage}[b]{\linewidth}\raggedleft
M8 (đầy đủ)
\end{minipage} \\
\midrule\noalign{}
\endhead
\bottomrule\noalign{}
\endlastfoot
FSTS (định tâm) & +2,652*** & +2,165** & +2,322*** & +1,952** &
+2,768*** & +2,409** \\
FSTS² & −1,705† & −1,168 & −1,389 & −0,965 & −2,959** & −2,543* \\
TCI (z) & --- & +0,168*** & --- & +0,153*** & --- & +0,153*** \\
DAI (z) & --- & --- & +0,104** & +0,077* & +0,046 & +0,019 \\
FSTS × DAI & --- & --- & --- & --- & −1,144 & −1,167† \\
FSTS² × DAI & --- & --- & --- & --- & +3,098** & +3,119** \\
\(R^2\) & 0,178 & 0,199 & 0,184 & 0,202 & 0,193 & 0,211 \\
\(R^2\) hiệu chỉnh & 0,168 & 0,189 & 0,174 & 0,190 & 0,180 & 0,196 \\
N & 623 & 623 & 617 & 617 & 617 & 617 \\
\end{longtable}

\emph{Nguồn: P4 Singapore, \texttt{p4/p4\_singapore\_en\_clean.md} (Bảng
2 bản thảo). † \(p < {,}10\); * \(p < {,}05\); ** \(p < {,}01\); ***
\(p < {,}001\). Sai số chuẩn vững HC1; FSTS định tâm trước khi bình
phương; mọi mô hình có quy mô, tuổi, sở hữu nước ngoài và hiệu ứng cố
định ngành. Các mô hình chứa DAI dùng \(N = 617\) do sáu cơ sở thiếu giá
trị biến liên quan DAI.}

Tiến triển độ phù hợp mô hình củng cố trật tự đóng góp của hai cấu trúc.
Việc bổ sung TCI nâng \(R^2\) từ \(0{,}178\) (M2) lên \(0{,}199\) (M5),
một mức nhảy lớn hơn so với khi bổ sung riêng DAI (lên \(0{,}184\) ở
M6); mô hình đầy đủ M8 đạt sức giải thích cao nhất trong các mô hình
chính (\(R^2 = 0{,}211\); \(R^2\) hiệu chỉnh \(= 0{,}196\)). Quỹ đạo suy
giảm của hệ số DAI từ M6 qua M7 đến M8
(\(+0{,}104 \to +0{,}077 \to +0{,}019\)) cho thấy phần lớn phương sai
năng suất mà DAI nắm bắt ở dạng hiệu ứng mức thực ra trùng lặp với TCI
hoặc được chuyển vào khối tương tác, bác bỏ cách hiểu DAI như một phần
thưởng năng suất phổ quát. Đặc trưng mẫu nền tảng cho cách đọc này:
trong \(N = 623\) cơ sở, năng suất lao động dạng log trung bình
\(11{,}42\) (độ lệch chuẩn \(1{,}18\)),
\(\overline{\mathrm{FSTS}} = 0{,}045\), và tương quan TCI--DAI chỉ ở mức
\(r = 0{,}31\), đủ thấp để biện minh cho việc giữ hai cấu trúc tách biệt
thay vì gộp thành một thước đo duy nhất.

\hypertarget{phuxe2n-xux1eed-ba-cux1a1-chux1ebf-cux1ea1nh-tranh-vuxe0-ux111ux1ed9-vux1eefng-trux1b0ux1edbc-hiux1ec7u-chux1ec9nh-chux1ecdn-lux1ef1a}{%
\subsection{4.3.9 Phân xử ba cơ chế cạnh tranh và độ vững trước hiệu
chỉnh chọn
lựa}\label{phuxe2n-xux1eed-ba-cux1a1-chux1ebf-cux1ea1nh-tranh-vuxe0-ux111ux1ed9-vux1eefng-trux1b0ux1edbc-hiux1ec7u-chux1ec9nh-chux1ecdn-lux1ef1a}}

Dấu dương của số hạng điều tiết bậc hai cho phép phân xử giữa ba cơ chế
cạnh tranh: thay thế (substitution), khuếch đại (amplification) và chọn
lựa (selection). Chữ ký thực nghiệm phân biệt khuếch đại với thay thế
nằm ở dấu của số hạng tương tác \emph{bậc hai}: hệ số dương
\(\mathrm{FSTS}^2 \times \mathrm{DAI} = +3{,}119\) (\(p = {,}005\)) hàm
ý đóng góp năng suất của DAI tăng siêu tuyến tính theo quy mô cường độ
xuất khẩu, điều mà một cơ chế thay thế dịch chuyển mức (chỉ dự đoán dịch
chuyển tuyến tính \(\mathrm{FSTS} \times \mathrm{DAI}\)) không thể tạo
ra. Tính bền vững của tương tác bậc hai dương qua các hạn chế mẫu (loại
siêu nhỏ R3, chỉ doanh nghiệp nhỏ và vừa R4, chỉ doanh nghiệp xuất khẩu
R5) lại bác bỏ một câu chuyện chọn lựa thuần túy, vì dưới chọn lựa,
tương quan đồng thời năng suất--DAI--FSTS lẽ ra phải suy giảm đáng kể
khi loại các doanh nghiệp dẫn dắt chọn lựa, song dấu dương được bảo toàn
xuyên suốt. Cụ thể, dấu dương và độ lớn của
\(\mathrm{FSTS}^2 \times \mathrm{DAI}\) giữ ổn định ở \(+3{,}521\) khi
loại doanh nghiệp siêu nhỏ (\(N = 464\)), \(+3{,}505\) khi giới hạn
doanh nghiệp nhỏ và vừa (\(N = 595\)) và \(+2{,}821\) ở mẫu chỉ doanh
nghiệp xuất khẩu với kiểm định F liên hợp có ý nghĩa (\(F = 6{,}32\),
\(p = {,}003\); \(N = 84\)).

Cơ chế khuếch đại còn được kiểm tra trước hiệu chỉnh chọn lựa biên gia
nhập. Vì thiết kế một đợt khảo sát Singapore 2023 không cung cấp công cụ
thỏa mãn ràng buộc loại trừ, một chẩn đoán dạng rút gọn ước lượng probit
giai đoạn một về tham gia xuất khẩu rồi đưa tỉ số Mills nghịch đảo vào
mô hình M8: hệ số IMR (\(\hat{\beta} = 0{,}264\), \(SE = 0{,}138\),
\(t = 1{,}92\), \(p = {,}055\)) chỉ vừa vượt ngưỡng, cho thấy chọn lựa
nhẹ ở biên gia nhập, nhưng tương tác trọng tâm
\(\mathrm{FSTS}^2 \times \mathrm{DAI}\) và hệ số TCI gần như không đổi
(\(|\Delta| < 0{,}02\) trong mỗi trường hợp). Phù hợp với cảnh báo trong
tài liệu kinh tế lượng rằng hiệu chỉnh Heckman thiếu công cụ loại trừ
khả tín có thể chệch hơn cả ước lượng OLS cần sửa, mô hình OLS được giữ
làm mô hình chính. Tổng hợp lại, bằng chứng Singapore ủng hộ khuếch đại
là cơ chế trọng tâm, giữ chọn lựa như một quan ngại dư có giới hạn, và
đặt Singapore như một trường hợp biên ở trần thể chế nơi bão hòa số nội
địa xóa hiệu ứng phần thưởng phổ quát còn tầng giao dịch Tầng 2 sẵn có
về mặt thể chế làm cơ chế mở rộng quy mô có điều kiện trở nên nhận diện
được về thực nghiệm.

\begin{center}\rule{0.5\linewidth}{0.5pt}\end{center}

\hypertarget{kux1ebft-quux1ea3-bux1ed1i-cux1ea3nh-thux1ec3-chux1ebf-trung-buxecnh-cao-trung-quux1ed1c-20122024-p5}{%
\section{4.4 Kết quả bối cảnh thể chế trung bình cao: Trung Quốc,
2012--2024
(P5)}\label{kux1ebft-quux1ea3-bux1ed1i-cux1ea3nh-thux1ec3-chux1ebf-trung-buxecnh-cao-trung-quux1ed1c-20122024-p5}}

\hypertarget{xuxe1c-nhux1eadn-quan-hux1ec7-phi-tuyux1ebfn-vuxe0-ux111iux1ec3m-uux1ed1n}{%
\subsection{4.4.1 Xác nhận quan hệ phi tuyến và điểm
uốn}\label{xuxe1c-nhux1eadn-quan-hux1ec7-phi-tuyux1ebfn-vuxe0-ux111iux1ec3m-uux1ed1n}}

Phân tích hai giai đoạn trên dữ liệu WBES Trung Quốc (Nhóm III, trung
bình cao) giai đoạn 2012--2024 xác nhận dạng quan hệ phi tuyến chữ U
ngược giữa FSTS và hiệu quả hoạt động kinh doanh của doanh nghiệp. Điểm
điểm uốn được ước lượng tại:

\begin{itemize}
\tightlist
\item
  \textbf{Năm 2012}: TP = 49,37\% FSTS
\item
  \textbf{Năm 2024}: TP = 47,19\% FSTS
\item
  \textbf{gộp (toàn mẫu)}: TP = 48,78\% FSTS
\end{itemize}

Các con số này cho thấy doanh nghiệp Trung Quốc đạt hiệu quả tối ưu khi
cường độ quốc tế hóa (FSTS) ở mức gần 49\%, một mức độ tập trung quốc tế
vừa phải trong bối cảnh nền kinh tế nội địa khổng lồ của Trung Quốc cũng
cung cấp cơ hội tăng trưởng song song.

\hypertarget{kiux1ec3m-ux111ux1ecbnh-tuxednh-ux1ed5n-ux111ux1ecbnh-cux1ee7a-ux111iux1ec3m-uux1ed1n-theo-thux1eddi-gian}{%
\subsection{4.4.2 Kiểm định tính ổn định của điểm uốn theo thời
gian}\label{kiux1ec3m-ux111ux1ecbnh-tuxednh-ux1ed5n-ux111ux1ecbnh-cux1ee7a-ux111iux1ec3m-uux1ed1n-theo-thux1eddi-gian}}

Luận án đặt câu hỏi thực nghiệm quan trọng: liệu hình dạng phi tuyến và
điểm uốn có thay đổi qua thập kỷ 2012--2024, giai đoạn Trung Quốc chứng
kiến nhiều thay đổi chính sách thương mại, căng thẳng địa chính trị và
chuyển đổi số mạnh mẽ? Kiểm định Paternoster cho kết quả
\(z = +0{,}82\), \(p = {,}412\) (vế FSTS) và \(z = -0{,}61\),
\(p = {,}545\) (vế FSTS²), không có ý nghĩa thống kê. Điều này xác nhận
rằng \textbf{điểm uốn ổn định qua thời gian}: sự dịch chuyển từ TP =
49,37\% (2012) xuống TP = 47,19\% (2024) là không đáng kể về mặt thống
kê.

Kiểm định tương tác thời gian (Temporal H6): \(F = 1{,}83\),
\(p = {,}176\), không có ý nghĩa thống kê, nghĩa là hình dạng tổng thể
của quan hệ phi tuyến không biến đổi đáng kể qua giai đoạn phân tích.
Phát hiện này cho thấy trong bối cảnh trung bình cao như Trung Quốc, cơ
chế căn bản của quan hệ quốc tế hóa và hiệu quả duy trì tính ổn định
đáng kể, có thể phản ánh sự bù trừ giữa chi phí thích ứng tăng lên do
chính sách thương mại khó đoán và năng lực tổ chức ngày càng trưởng
thành của các doanh nghiệp xuất khẩu Trung Quốc.

\hypertarget{vai-truxf2-cux1ee7a-tci-vuxe0-dai}{%
\subsection{4.4.3 Vai trò của TCI và
DAI}\label{vai-truxf2-cux1ee7a-tci-vuxe0-dai}}

Tương tác \(\mathrm{FSTS} \times \mathrm{TCI}\) không đạt ý nghĩa thống
kê (null moderation, \(p > {,}10\)), cho thấy chỉ số năng lực công nghệ
(TCI) không phát huy vai trò điều tiết đáng kể ở bối cảnh Trung Quốc
theo mô hình P5. Tương tự, DAI cũng không cho thấy điều tiết có ý nghĩa
trong bối cảnh này. kết quả không có ý nghĩa về điều tiết tại Trung Quốc
có thể phản ánh thực tế rằng năng lực công nghệ đã được phân phối tương
đối đồng đều hơn trong mẫu doanh nghiệp xuất khẩu Trung Quốc, làm giảm
phương sai đủ để phát hiện hiệu ứng điều tiết.

\hypertarget{bux1ed9-hux1ec7-sux1ed1-muxf4-huxecnh-ngux1b0ux1ee1ng-chuxednh-m2-vuxe0-ux111ux1ed9-lux1edbn-kinh-tux1ebf}{%
\subsection{4.4.4 Bộ hệ số mô hình ngưỡng chính (M2) và độ lớn kinh
tế}\label{bux1ed9-hux1ec7-sux1ed1-muxf4-huxecnh-ngux1b0ux1ee1ng-chuxednh-m2-vuxe0-ux111ux1ed9-lux1edbn-kinh-tux1ebf}}

Mô hình ngưỡng chính M2 (mô hình
\(\ln(\mathrm{LP}) \sim \mathrm{FSTS} + \mathrm{FSTS}^2 + \ln(\mathrm{Emp}) + \text{firmage} + \text{foreigndummy}\))
được ước lượng riêng cho từng đợt khảo sát và mẫu gộp, củng cố cho hình
dạng chữ U ngược bằng một bộ hệ số đầy đủ. Hạng cường độ xuất khẩu tuyến
tính dương trong mọi mẫu: \(\hat{\beta}_1 = +2{,}07\) năm 2012
(\(p < {,}001\)), \(+1{,}50\) năm 2024 (\(p = {,}010\)) và \(+1{,}78\)
trong mẫu gộp (\(p < {,}001\)). Hạng bậc hai âm tương ứng:
\(\hat{\beta}_2 = -2{,}09\) năm 2012 (\(p < {,}001\)), \(-1{,}59\) năm
2024 (\(p = {,}026\)) và \(-1{,}83\) gộp (\(p < {,}001\)). Kiểm định U
Lind--Mehlum xác nhận chữ U ngược một cách chính thức với \(p < {,}001\)
trong cả mẫu 2012 và mẫu gộp, và \(p = {,}037\) trong mẫu 2024.

Khoảng tin cậy 95\% theo phương pháp delta của điểm uốn là
\([43{,}2\%; 55{,}6\%]\) năm 2012, \([34{,}5\%; 59{,}9\%]\) năm 2024 và
\([42{,}7\%; 54{,}9\%]\) trong mẫu gộp. Sự chồng lấp chặt chẽ của ba
khoảng tin cậy này là cơ sở định lượng trực tiếp cho kết luận về tính ổn
định theo thời gian đã trình bày ở Mục 4.4.2. Về độ lớn kinh tế, tại
điểm uốn đặc thù đợt khảo sát, \(\ln(\mathrm{LP})\) dự đoán cao hơn
khoảng \(0{,}51\) điểm log năm 2012 và \(0{,}35\) điểm log năm 2024 so
với mức cơ sở \(\mathrm{FSTS} = 0\), tức tương ứng khoảng \(67\%\) và
\(42\%\) chênh lệch năng suất ở mức kiểm soát trung bình hình học. Ngược
lại, tại \(\mathrm{FSTS} = 1\) (một nhà xuất khẩu 100\% giả định), năng
suất dự đoán gần như ngang bằng mức cơ sở trong cả hai đợt, minh hoạ
trực quan logic ``tối ưu có giới hạn'' nằm ở trung tâm của giả thuyết
H1.

\hypertarget{quy-muxf4-mux1eabu-ux111ux1eb7c-trux1b0ng-nhuxf3m-doanh-nghiux1ec7p-vuxe0-phuxe2n-tuxe1ch-tux103ng-trux1b0ux1edfng-nux103ng-suux1ea5t}{%
\subsection{4.4.5 Quy mô mẫu, đặc trưng nhóm doanh nghiệp và phân tách
tăng trưởng năng
suất}\label{quy-muxf4-mux1eabu-ux111ux1eb7c-trux1b0ng-nhuxf3m-doanh-nghiux1ec7p-vuxe0-phuxe2n-tuxe1ch-tux103ng-trux1b0ux1edfng-nux103ng-suux1ea5t}}

Mẫu phân tích cơ sở gồm \(N = 2.610\) quan sát năm 2012 và \(N = 1.934\)
năm 2024, với mẫu gộp \(N = 4.544\) trong mô hình ngưỡng chính. Cường độ
xuất khẩu trung bình thấp và lệch mạnh: \(6{,}9\%\) năm 2012 và
\(5{,}2\%\) năm 2024, với chỉ \(15{,}4\%\) doanh nghiệp năm 2012 và
\(12{,}6\%\) năm 2024 báo cáo cường độ xuất khẩu dương. Năng suất lao
động trung bình dạng log tăng từ \(12{,}52\) (2012) lên \(13{,}01\)
(2024), một dịch chuyển mức \(0{,}49\) điểm log tương đương năng suất
cao hơn khoảng \(64\%\) về mức (\(\exp(0{,}49) \approx 1{,}64\)). Một
phân tách Oaxaca--Blinder không chính thức quy khoảng \(70\%\) dịch
chuyển này cho tăng trưởng năng suất trong cùng ngành và khoảng \(30\%\)
cho thay đổi thành phần ngành.

Năng lực công nghệ vận hành bền vững như một yếu tố dịch chuyển mức trực
tiếp chứ không phải yếu tố điều tiết độ cong. Hệ số TCI chuẩn hoá z
trong cùng đợt là \(\hat{\beta}_z = +0{,}260\) năm 2012, \(+0{,}426\)
năm 2024 và \(+0{,}380\) trong mẫu gộp (đều \(p < {,}001\)), với kiểm
định Paternoster về thay đổi chéo đợt khảo sát cho \(z = -2{,}55\),
\(p = {,}011\), cho thấy liên hệ giữa năng lực công nghệ và năng suất
được củng cố có ý nghĩa giữa hai đợt.

\hypertarget{cux1ea5u-truxfac-kiux1ec3m-ux111ux1ecbnh-ux111iux1ec1u-tiux1ebft-ba-chiux1ec1u-vuxe0-cux1a1-chux1ebf-vux1ed1n-lux1b0u-ux111ux1ed9ng}{%
\subsection{4.4.6 Cấu trúc kiểm định điều tiết ba chiều và cơ chế vốn
lưu
động}\label{cux1ea5u-truxfac-kiux1ec3m-ux111ux1ecbnh-ux111iux1ec1u-tiux1ebft-ba-chiux1ec1u-vuxe0-cux1a1-chux1ebf-vux1ed1n-lux1b0u-ux111ux1ed9ng}}

Kết quả điều tiết không có ý nghĩa trình bày trong Mục 4.4.3 được nâng
đỡ bởi một bộ ba kiểm định F liên hợp trên mô hình điều tiết ba chiều
gộp (\(N = 3.559\)). Kiểm định dịch chuyển chéo đợt khảo sát cho
\(F(2, 3.558) = 2{,}24\), \(p = {,}107\) (không bác bỏ tính bằng nhau);
kiểm định điều tiết độ cong theo năng lực cho \(F(2, 3.558) = 3{,}26\),
\(p = {,}039\) (chỉ ở mức biên); và kiểm định điều tiết động phụ thuộc
năng lực cho \(F(2, 3.558) = 0{,}27\), \(p = {,}760\) (không hỗ trợ).
Đáng chú ý, dù kiểm định F liên hợp về điều tiết độ cong đạt ý nghĩa
biên, không một hệ số tương tác riêng lẻ nào phân biệt được với 0:
\(\text{TCI} \times \mathrm{FSTS}\) cho \(\beta = -0{,}414\),
\(p = {,}443\) và \(\text{TCI} \times \mathrm{FSTS}^2\) cho
\(\beta = +0{,}051\), \(p = {,}934\), củng cố cách diễn giải rằng năng
lực công nghệ là yếu tố dịch chuyển mức chứ không phải yếu tố điều tiết
ngưỡng.

Việc khảo sát cơ chế vốn lưu động (working capital) cũng không nhận dạng
được một yếu tố điều tiết bền vững. Tương tác thấu chi
\(\times \mathrm{FSTS}^2\) đảo dấu giữa hai đợt (\(\beta = +0{,}27\),
\(p = {,}050\) năm 2012; \(\beta = -0{,}38\), \(p = {,}081\) năm 2024);
tương tác tín dụng thương mại \(\times \mathrm{FSTS}^2\) âm có ý nghĩa
năm 2012 (\(\beta = -0{,}017\), \(p = {,}009\)) nhưng không có ý nghĩa
năm 2024 (\(p = {,}660\)). Như vậy, cơ chế bẫy vốn lưu động vẫn là cơ sở
lý thuyết nhất quán nhất cho đoạn suy giảm sau ngưỡng nhưng chưa được
nhận dạng trực tiếp với công cụ WBES hiện có, đòi hỏi dữ liệu vi mô tài
chính như thời gian chu kỳ chuyển đổi tiền mặt hoặc số ngày khoản phải
thu chưa thanh toán cho nghiên cứu tương lai.

\hypertarget{kiux1ec3m-ux111ux1ecbnh-ux111ux1ed9-vux1eefng-mux1eabu-phux1ee5-chux1ec9-nguxe0nh-sux1ea3n-xuux1ea5t}{%
\subsection{4.4.7 Kiểm định độ vững: mẫu phụ chỉ ngành sản
xuất}\label{kiux1ec3m-ux111ux1ecbnh-ux111ux1ed9-vux1eefng-mux1eabu-phux1ee5-chux1ec9-nguxe0nh-sux1ea3n-xuux1ea5t}}

Để loại trừ quan ngại trộn lẫn ngành, mô hình M2 được tái ước lượng trên
mẫu phụ chỉ ngành sản xuất (\(N = 1.656\) năm 2012, \(N = 1.062\) năm
2024, gộp \(N = 2.718\)). Dưới ràng buộc này, điểm uốn dịch về khoảng
\(42\%\) năm 2012 (khoảng tin cậy \([37{,}8; 46{,}8]\)) và \(30\%\) năm
2024 (khoảng tin cậy \([15{,}1; 44{,}2]\)), khoảng cuối rộng hơn nhiều
do sức mạnh thống kê giảm, song kiểm định Paternoster vẫn không bác bỏ
tính bằng nhau chéo đợt khảo sát (\(z = +1{,}51\), \(p = {,}130\) cho
FSTS; \(z = -0{,}96\), \(p = {,}337\) cho \(\mathrm{FSTS}^2\)). Khi bổ
sung tầng ngành như hiệu ứng cố định vào M2 gộp, hệ số FSTS giảm từ
\(+1{,}78\) xuống \(+1{,}56\) (giảm \(12{,}8\%\)) và hệ số
\(\mathrm{FSTS}^2\) từ \(-1{,}83\) xuống \(-1{,}55\) (giảm
\(15{,}2\%\)), nhưng chữ U ngược được bảo toàn và điểm uốn vẫn nằm trong
dải vận hành an toàn 30--60\%.

\hypertarget{bux1eb1ng-chux1ee9ng-bux1ed5-trux1ee3-cux1ea5p-doanh-nghiux1ec7p-nhux1ecf-vuxe0-vux1eeba-chux1eef-u-ngux1b0ux1ee3c-vuxe0-cux1a1-chux1ebf-bux1eaby-vux1ed1n-lux1b0u-ux111ux1ed9ng-p2}{%
\subsection{4.4.8 Bằng chứng bổ trợ cấp doanh nghiệp nhỏ và vừa: chữ U
ngược và cơ chế bẫy vốn lưu động
(P2)}\label{bux1eb1ng-chux1ee9ng-bux1ed5-trux1ee3-cux1ea5p-doanh-nghiux1ec7p-nhux1ecf-vuxe0-vux1eeba-chux1eef-u-ngux1b0ux1ee3c-vuxe0-cux1a1-chux1ebf-bux1eaby-vux1ed1n-lux1b0u-ux111ux1ed9ng-p2}}

Nghiên cứu thành phần P2 cung cấp một kiểm định độc lập về quan hệ chữ U
ngược trong bối cảnh Trung Quốc, trên một lát cắt khác của tổng thể:
\textbf{4.290 quan sát doanh nghiệp và năm} thuộc các doanh nghiệp sản
xuất nhỏ và vừa (mã ngành ISIC 15--37). Cường độ xuất khẩu trung bình
của mẫu là 23,7\% (độ lệch chuẩn 28,4\%), trong đó 18\% doanh nghiệp có
cường độ xuất khẩu vượt 50\%, tạo nên một phân bố trải rộng từ thuần nội
địa đến định hướng xuất khẩu cao, phù hợp để định vị điểm uốn.

Ước lượng mô hình bậc hai với hiệu ứng cố định ngành và năm xác nhận chữ
U ngược: hệ số nhánh đi lên dương và có ý nghĩa cao
(\(\hat{\beta}_1 = +1{,}704\), \(p < {,}001\)), hệ số nhánh đi xuống âm
và có ý nghĩa, định vị điểm uốn ở mức \textbf{47,8\% cường độ xuất
khẩu}. Giá trị này nằm sát ngưỡng của nghiên cứu Trung Quốc theo chuỗi
thời gian (P5: 48,78\%) và trong dải vận hành an toàn 30--60\% được xác
nhận xuyên suốt luận án. Đóng góp đặc thù của P2 là cơ chế giải thích
nhánh đi xuống: \textbf{bẫy vốn lưu động}. Khi cường độ xuất khẩu vượt
47,8\%, chi phí biên của việc tài trợ chu kỳ chuyển đổi tiền mặt kéo
dài, do doanh nghiệp xuất khẩu phải ứng vốn cho sản xuất, đóng gói và
vận chuyển từ rất lâu trước khi thu được khoản phải thu nước ngoài, bắt
đầu lấn át lợi ích từ quy mô và học hỏi qua xuất khẩu, kéo năng suất đi
xuống. Cơ chế này bổ sung cho cách đọc dựa trên chi phí điều phối và
năng lực, làm rõ một kênh tài chính cụ thể đứng sau độ cong ở các doanh
nghiệp nhỏ và vừa thâm dụng vốn lưu động. Nhất quán với phần còn lại của
luận án, năng lực số và đặc điểm giới tính của nhà quản trị trong P2 vận
hành như yếu tố nâng mặt bằng năng suất chứ không uốn lại hình dạng
đường cong.

\begin{center}\rule{0.5\linewidth}{0.5pt}\end{center}

\hypertarget{kux1ebft-quux1ea3-bux1ed1i-cux1ea3nh-chuyux1ec3n-ux111ux1ed5i-bux1eadc-thux1ea5ptrung-viux1ec7t-nam-p3}{%
\section{4.5 Kết quả bối cảnh chuyển đổi bậc thấp--trung: Việt Nam
(P3)}\label{kux1ebft-quux1ea3-bux1ed1i-cux1ea3nh-chuyux1ec3n-ux111ux1ed5i-bux1eadc-thux1ea5ptrung-viux1ec7t-nam-p3}}

\hypertarget{xuxe1c-nhux1eadn-quan-hux1ec7-phi-tuyux1ebfn-vuxe0-hux1ec7-sux1ed1-hux1ed3i-quy}{%
\subsection{4.5.1 Xác nhận quan hệ phi tuyến và hệ số hồi
quy}\label{xuxe1c-nhux1eadn-quan-hux1ec7-phi-tuyux1ebfn-vuxe0-hux1ec7-sux1ed1-hux1ed3i-quy}}

Phân tích hồi quy bảng (panel regression) trên ba làn sóng khảo sát WBES
Việt Nam (2009, 2015, 2023) xác nhận dạng phi tuyến chữ U ngược. Mô hình
M2, mô hình phi tuyến chính (mẫu gộp, \(N = 2.958\)), cho kết quả:

\[\hat{\beta}(\text{FSTS}_c) = +0{,}984, \quad \hat{\beta}(\text{FSTS}_c^2) = -1{,}909\]

Cả hai hệ số đều có ý nghĩa thống kê, với hệ số bậc nhất dương và bậc
hai âm xác nhận đường cong chữ U ngược. Kiểm định hình thức bằng phương
pháp Lind--Mehlum (U-test) cho kết quả \(p < {,}001\), bằng chứng thống
kê mạnh nhất trong tất cả các mẫu quốc gia của luận án.

\hypertarget{ux111iux1ec3m-uux1ed1n-theo-luxe0n-suxf3ng}{%
\subsection{4.5.2 Điểm uốn theo làn
sóng}\label{ux111iux1ec3m-uux1ed1n-theo-luxe0n-suxf3ng}}

Điểm uốn được ước lượng riêng cho từng làn sóng và cho gộp:

\begin{longtable}[]{@{}ll@{}}
\toprule\noalign{}
Làn sóng & điểm uốn \\
\midrule\noalign{}
\endhead
\bottomrule\noalign{}
\endlastfoot
2009 & 46,2\% FSTS \\
2015 & 39,3\% FSTS \\
2023 & 41,6\% FSTS \\
gộp & 39,7\% FSTS \\
\end{longtable}

Điểm điểm uốn trong khoảng \textbf{39--46\% FSTS} đặt Việt Nam ở mức
``thấp hơn đáng kể'' so với Singapore (TP ≈ 88,6\%) và tương đương hoặc
thấp hơn Trung Quốc (TP ≈ 48,78\%). Kiểm định Paternoster so sánh điểm
uốn giữa 2009 và 2015 cho kết quả \(z = 3{,}353\), \(p < {,}001\), có ý
nghĩa thống kê, cho thấy điểm uốn đã dịch chuyển đáng kể từ 46,2\% xuống
39,3\% giữa hai làn sóng.

\hypertarget{ux111iux1ec1u-tiux1ebft-bux1edfi-tci-tuxedch-cux1ef1c-vuxe0-cuxf3-uxfd-nghux129a}{%
\subsection{4.5.3 Điều tiết bởi TCI: Tích cực và có ý
nghĩa}\label{ux111iux1ec1u-tiux1ebft-bux1edfi-tci-tuxedch-cux1ef1c-vuxe0-cuxf3-uxfd-nghux129a}}

Mô hình M3 với tương tác \(\mathrm{FSTS} \times \mathrm{TCI}\) cho kết
quả \textbf{dương và có ý nghĩa thống kê}: TCI làm nâng cao mặt bằng
hiệu quả và khuếch đại tác động tích cực của quốc tế hóa. Doanh nghiệp
Việt Nam có năng lực công nghệ tốt hơn (đo bằng TCI) đạt hiệu quả cao
hơn từ quá trình quốc tế hóa, nhất quán với H2 trong khung lý thuyết
CDCM.

\hypertarget{ux111iux1ec1u-tiux1ebft-bux1edfi-dai-khuxf4ng-cuxf3-uxfd-nghux129a}{%
\subsection{4.5.4 Điều tiết bởi DAI: Không có ý
nghĩa}\label{ux111iux1ec1u-tiux1ebft-bux1edfi-dai-khuxf4ng-cuxf3-uxfd-nghux129a}}

Mô hình M4 với tương tác \(\mathrm{FSTS} \times \mathrm{DAI}\) (chỉ bao
gồm Tầng 1 DAI do giới hạn đo lường WBES trong các làn sóng này) cho kết
quả yếu dương nhưng không có ý nghĩa thống kê. Kết quả này có thể phản
ánh hạn chế về đo lường: Tầng 1 DAI chỉ nắm bắt mức độ chấp nhận số bề
mặt (website, email, online sales) chứ không đo lường được chiều sâu
năng lực số. Với dữ liệu WBES Việt Nam chỉ cung cấp Tầng 1, không đủ
phân biệt để phát hiện hiệu ứng điều tiết tinh tế hơn.

\hypertarget{bux1ed9-hux1ec7-sux1ed1-ux111ux1ea7y-ux111ux1ee7-theo-luxe0n-suxf3ng-vuxe0-quux1ef9-ux111ux1ea1o-ux111iux1ec1u-tiux1ebft}{%
\subsection{4.5.5 Bộ hệ số đầy đủ theo làn sóng và quỹ đạo điều
tiết}\label{bux1ed9-hux1ec7-sux1ed1-ux111ux1ea7y-ux111ux1ee7-theo-luxe0n-suxf3ng-vuxe0-quux1ef9-ux111ux1ea1o-ux111iux1ec1u-tiux1ebft}}

Bộ hệ số chữ U ngược M2 nhất quán qua cả ba làn sóng. Hạng tuyến tính
\(\mathrm{FSTS}^{c}\) dương trong mọi đợt (\(\hat{\beta}_1 = +1{,}045\)
năm 2009, \(p = {,}015\); \(+1{,}159\) năm 2015, \(p = {,}029\);
\(+0{,}962\) năm 2023, \(p = {,}039\); \(+0{,}984\) gộp,
\(p < {,}001\)), trong khi hạng bậc hai âm (\(\hat{\beta}_2 = -1{,}823\)
năm 2009, \(p = {,}005\); \(-2{,}115\) năm 2015, \(p = {,}004\);
\(-1{,}686\) năm 2023, \(p = {,}008\); \(-1{,}909\) gộp,
\(p < {,}001\)). Độ cong sắc nét nhất quan sát được trong làn sóng 2015
(\(\hat{\beta}_2 = -2{,}115\)), nhất quán với một giai đoạn chuyển đổi
trong đó kênh số bị nén hoàn toàn. Khoảng tin cậy 95\% theo phương pháp
delta cho từng điểm uốn lần lượt là \([37{,}4\%; 55{,}1\%]\) năm 2009,
\([30{,}3\%; 48{,}4\%]\) năm 2015, \([31{,}7\%; 51{,}5\%]\) năm 2023 và
\([34{,}0\%; 45{,}5\%]\) trong mẫu gộp.

Năng lực công nghệ \(\tilde{\mathrm{TCI}}\) có liên hệ trực tiếp thuận
chiều và bền vững qua cả ba đợt (\(\hat{\beta} = +0{,}215\) năm 2009,
\(+0{,}128\) năm 2015, \(+0{,}123\) năm 2023, đều có ý nghĩa) và trong
mẫu gộp (\(\hat{\beta} = +0{,}179\), \(p < {,}001\)), với vai trò điều
tiết phân biệt được trong ba trong bốn bảng (kiểm định liên hợp M3:
\(p = {,}040\) năm 2009, \(p = {,}027\) năm 2023, \(p = {,}003\) gộp,
chỉ \(p = {,}713\) năm 2015). Trong mẫu gộp, tương tác tuyến tính
\(\mathrm{FSTS}^{c} \times \tilde{\mathrm{TCI}} = -0{,}587\)
(\(p = {,}003\)) và bậc hai
\((\mathrm{FSTS}^{c})^2 \times \tilde{\mathrm{TCI}} = +0{,}640\)
(\(p = {,}031\)), cho thấy chữ U ngược phẳng lại đối với doanh nghiệp có
năng lực cao chứ không dịch chuyển về mức.

\hypertarget{tuxednh-phux1ee5-thuux1ed9c-giai-ux111oux1ea1n-cux1ee7a-dai-vuxe0-kiux1ec3m-ux111ux1ecbnh-paternoster-chuxe9o-luxe0n-suxf3ng}{%
\subsection{4.5.6 Tính phụ thuộc giai đoạn của DAI và kiểm định
Paternoster chéo làn
sóng}\label{tuxednh-phux1ee5-thuux1ed9c-giai-ux111oux1ea1n-cux1ee7a-dai-vuxe0-kiux1ec3m-ux111ux1ecbnh-paternoster-chuxe9o-luxe0n-suxf3ng}}

Trái với năng lực công nghệ, áp dụng công nghệ số nền tảng
\(\tilde{\mathrm{DAI}}\) theo một quỹ đạo không đơn điệu: mạnh nhất năm
2009 (\(\hat{\beta} = +0{,}175\), \(p < {,}001\)), mất ý nghĩa năm 2015
(\(\hat{\beta} = -0{,}044\), \(p = {,}377\)), tái xuất hiện năm 2023
(\(\hat{\beta} = +0{,}095\), \(p = {,}038\)) và dương trong mẫu gộp
(\(\hat{\beta} = +0{,}078\), \(p = {,}004\)). Kiểm định Paternoster xác
nhận các dịch chuyển này phân biệt được về mặt thống kê: mức giảm từ
2009 đến 2015 cho \(z = 3{,}35\) (\(p = {,}001\)) và sự phục hồi từ 2015
đến 2023 cho \(z = -2{,}05\) (\(p = {,}040\)), trong khi các tham số độ
cong \(\mathrm{FSTS}^{c}\) và \((\mathrm{FSTS}^{c})^2\) không phân biệt
được giữa các đợt (mọi \(p > {,}33\)). Điều tiết DAI tập trung hoàn toàn
trong làn sóng 2023, nơi
\(\mathrm{FSTS}^{c} \times \tilde{\mathrm{DAI}} = -0{,}912\)
(\(p = {,}043\)) và kiểm định liên hợp M8 ở mức cận ngưỡng
(\(p = {,}062\)). Theo phân loại Haans, Pieters và He (2016), mẫu hình
2023 phù hợp với điều tiết Loại I (làm phẳng độ dốc, giữ nguyên hình
dạng chữ U ngược) chứ không phải Loại II.

\hypertarget{phuxe2n-ruxe3-tham-giacux1b0ux1eddng-ux111ux1ed9-vuxe0-bux1eb1ng-chux1ee9ng-nux1ed9i-sinh}{%
\subsection{4.5.7 Phân rã tham gia--cường độ và bằng chứng nội
sinh}\label{phuxe2n-ruxe3-tham-giacux1b0ux1eddng-ux111ux1ed9-vuxe0-bux1eb1ng-chux1ee9ng-nux1ed9i-sinh}}

Một phát hiện thực nghiệm cốt lõi của P3 là chữ U ngược toàn mẫu được
nhận diện chủ yếu thông qua biên tham gia chứ không qua độ cong trong
nhóm doanh nghiệp xuất khẩu. Tỉ lệ doanh nghiệp báo cáo xuất khẩu trực
tiếp dương giảm từ \(28{,}4\%\) năm 2009 xuống \(20{,}7\%\) năm 2015 và
\(18{,}8\%\) năm 2023, và chỉ khoảng \(1{,}0\%\) doanh nghiệp gộp (29
trong 2.958) nằm trong dải \(\pm 5\) điểm phần trăm quanh điểm uốn. Khi
tái ước lượng trên nhóm phụ chỉ doanh nghiệp xuất khẩu
(\(\mathrm{FSTS} > 0\); gộp \(N = 669\)), hạng tuyến tính trở nên âm
(\(\mathrm{FSTS}^{c}\) \(\hat{\beta} = -0{,}861\), \(p < {,}001\)) còn
hạng bậc hai mất ý nghĩa (\((\mathrm{FSTS}^{c})^2\)
\(\hat{\beta} = -0{,}200\), \(p = {,}660\); kiểm định liên hợp
\(p = {,}462\)). Do đó độ cong toàn mẫu phản ánh chủ yếu bước nhảy năng
suất khi doanh nghiệp chuyển từ không xuất khẩu sang xuất khẩu (biên
tham gia), chứ không phải độ cong cường độ trong nhóm doanh nghiệp xuất
khẩu (biên cường độ).

Bằng chứng nội sinh và chọn lựa làm rõ thêm sự phân biệt giữa hai kênh
năng lực. Hiệu chỉnh Heckman hai bước cho tỉ số Mills nghịch đảo
(inverse Mills ratio) không có ý nghĩa qua mọi đợt (mọi
\(|\lambda| < 0{,}84\), \(p > {,}25\)), cho thấy không phát hiện thiên
lệch chọn lựa trên nhóm phụ doanh nghiệp xuất khẩu. Ước lượng biến công
cụ hai giai đoạn (2SLS) với công cụ là tỉ lệ áp dụng đồng cấp ngành ×
vùng × đợt (thống kê F giai đoạn một 22--35, vượt xa ngưỡng
Staiger--Stock) cho hệ số DAI được công cụ hoá bằng không
(\(\hat{\beta} = 0{,}018\), \(p = {,}942\)) nhưng hệ số TCI được công cụ
hoá dương mạnh (\(\hat{\beta} = 1{,}639\), \(p < {,}001\)). Hệ số DAI
2SLS bằng không xác nhận rằng hiện diện website không gây ra cải thiện
năng suất một cách khả tín ở Việt Nam 2023, củng cố cách diễn giải
``tính lỗi thời của biến đại diện Tầng 1'': đến năm 2023 tỉ lệ doanh
nghiệp có website đạt \(49{,}8\%\) so với \(42{,}5\%\) năm 2009, khiến
chỉ báo này gần như phổ quát và mất sức mạnh phân biệt. Đối sánh điểm xu
hướng (propensity score matching) xác nhận thêm: tác động bình quân lên
nhóm được xử lý cho sở hữu website là \(0{,}298\)--\(0{,}321\)
(\(p < {,}001\)) và cho công nghệ nước ngoài hoặc chứng nhận là
\(0{,}609\)--\(0{,}637\) (\(p < {,}001\)).

\hypertarget{cux1ea5u-truxfac-lux1ea1m-phuxe1t-sux1ed1-khuxf4ng-cux1ee7a-fsts-vuxe0-hux1ec7-quux1ea3-nhux1eadn-dux1ea1ng-hai-biuxean}{%
\subsection{4.5.8 Cấu trúc lạm phát số không của FSTS và hệ quả nhận
dạng hai
biên}\label{cux1ea5u-truxfac-lux1ea1m-phuxe1t-sux1ed1-khuxf4ng-cux1ee7a-fsts-vuxe0-hux1ec7-quux1ea3-nhux1eadn-dux1ea1ng-hai-biuxean}}

Bản chất nhận dạng của chữ U ngược Việt Nam chỉ trở nên minh bạch khi
đặt vào cấu trúc phân phối của biến cường độ xuất khẩu. Biến FSTS bị
chặn dưới ở 0 và lạm phát số không (zero-inflated) nghiêm trọng: tỉ
trọng doanh nghiệp không xuất khẩu là \(71{,}6\%\) năm 2009,
\(79{,}3\%\) năm 2015 và \(81{,}2\%\) năm 2023, mức gộp \(77{,}4\%\). Vì
khối lượng quan sát dồn ở \(\mathrm{FSTS} = 0\) trong khi điểm uốn nằm
quanh 40\%, quan hệ quốc tế hóa và hiệu quả phân tách về mặt khái niệm
thành hai biên độc lập: biên tham gia, tức bước nhảy từ
\(\mathrm{FSTS} = 0\) sang \(\mathrm{FSTS} > 0\) nơi doanh nghiệp vượt
ngưỡng gia nhập và hấp thụ chi phí gia nhập cố định; và biên cường độ,
tức biến thiên FSTS trong nội bộ nhóm doanh nghiệp xuất khẩu. Lý thuyết
dự đoán phần thưởng năng suất có thể phát sinh ở cả hai biên, nhưng
trong các nền kinh tế chuyển đổi với hạ tầng xuất khẩu mỏng, biên tham
gia là biên nâng năng suất chủ đạo, còn biên cường độ thể hiện lợi suất
giảm dần hoặc bằng không một khi đã vượt ngưỡng tham gia. Cơ chế lý
thuyết đứng sau nhánh đi xuống là chi phí giao dịch thể chế (Williamson,
1985; Hennart, 2007): trong nền kinh tế chuyển đổi, khe hở thực thi hợp
đồng, thiếu hụt logistics và bất cân xứng thông tin làm tăng chi phí
biên quản lý mỗi thị trường nước ngoài bổ sung, khiến nhánh suy giảm sau
ngưỡng bám ở mức cường độ xuất khẩu thấp hơn so với các bối cảnh nơi hạ
tầng số và chất lượng thực thi hợp đồng đã hấp thụ phần lớn ma sát này.
Đây chính là cơ sở lý thuyết giải thích vì sao điểm uốn Việt Nam
(39--46\%) nằm thấp hơn đáng kể so với điểm uốn Singapore
(\textasciitilde88,6\%) ở biên thể chế mạnh nhất.

\hypertarget{bux1ed9-hux1ec7-sux1ed1-biuxean-ux111ux1ed9-kinh-tux1ebf-tux1eebng-mux1ee9c-fsts-vuxe0-ux111ux1ecbnh-lux1b0ux1ee3ng-bux1b0ux1edbc-nhux1ea3y-tham-gia}{%
\subsection{4.5.9 Bộ hệ số biên độ kinh tế từng mức FSTS và định lượng
bước nhảy tham
gia}\label{bux1ed9-hux1ec7-sux1ed1-biuxean-ux111ux1ed9-kinh-tux1ebf-tux1eebng-mux1ee9c-fsts-vuxe0-ux111ux1ecbnh-lux1b0ux1ee3ng-bux1b0ux1edbc-nhux1ea3y-tham-gia}}

Bộ hiệu ứng biên theo từng mức cường độ xuất khẩu trong mô hình gộp M2
(\(N = 2.958\)) cho phép định lượng độ lớn kinh tế của đường cong thay
vì chỉ đọc dấu hệ số. Đạo hàm biên
\(\mathrm{d}y/\mathrm{d}\mathrm{FSTS} = \beta_1 + 2\beta_2(\mathrm{FSTS} - \overline{\mathrm{FSTS}})\)
với \(\beta_1 = 0{,}984\), \(\beta_2 = -1{,}909\) và
\(\overline{\mathrm{FSTS}} \approx 13{,}9\%\) cho độ dốc \(+1{,}13\) tại
mức 10\% FSTS, \(+0{,}75\) tại 20\%, \(+0{,}37\) tại 30\%, triệt tiêu
tại điểm uốn \(39{,}7\%\), rồi đảo âm \(-0{,}39\) tại 50\% và
\(-0{,}78\) tại 60\%. Tại hai biên dữ liệu, độ dốc là \(+1{,}52\) ở
\(\mathrm{FSTS} = 0\) và \(-2{,}30\) ở \(\mathrm{FSTS} = 100\%\), hai
dấu trái ngược mà kiểm định Lind--Mehlum xác nhận là không tương thích
với quan hệ đơn điệu (\(p < {,}001\)). Bốn điều kiện nhận diện chữ U
ngược của Haans và cộng sự (2016) đều thỏa trong mẫu gộp: C1 độ dốc
dương ở biên dưới (\(\beta_1 = +0{,}984\), \(p < {,}001\)); C2 độ cong
âm (\(\beta_2 = -1{,}909\), \(p < {,}001\)); C3 điểm uốn \(39{,}7\%\)
nằm trong miền dữ liệu \([0\%, 100\%]\); và C4 dấu độ dốc đảo chiều giữa
hai biên. Độ lớn của bước nhảy biên tham gia được thước đo đối sánh điểm
xu hướng xác nhận độc lập: tác động bình quân lên nhóm được xử lý của
trạng thái có công nghệ nước ngoài hoặc chứng nhận đạt
\(0{,}609\)--\(0{,}637\) điểm log (\(p < {,}001\)), tức khoảng 84--89\%
chênh lệch năng suất, lớn hơn nhiều so với tác động của riêng sở hữu
website (\(0{,}298\)--\(0{,}321\) điểm log).

\hypertarget{pin-chux1eb7t-ux111ux1ed9-vux1eefng-tuxe1ch-nguxe0nh-muxf4-huxecnh-bux1eadc-ba-vuxe0-batterie-kiux1ec3m-ux111ux1ecbnh-nux1ed9i-sinh}{%
\subsection{4.5.10 Pin chặt độ vững: tách ngành, mô hình bậc ba và
batterie kiểm định nội
sinh}\label{pin-chux1eb7t-ux111ux1ed9-vux1eefng-tuxe1ch-nguxe0nh-muxf4-huxecnh-bux1eadc-ba-vuxe0-batterie-kiux1ec3m-ux111ux1ecbnh-nux1ed9i-sinh}}

Diễn giải chữ U ngược và sự phân biệt hai cấu trúc năng lực được nâng đỡ
bởi một loạt kiểm định độ vững có hệ thống. Thứ nhất, tách theo ngành
cho thấy chữ U ngược tồn tại ở cả hai mảng nhưng sắc nét hơn ở chế biến
chế tạo (\(\hat{\beta}_1 = +0{,}971\), \(p = {,}001\);
\(\hat{\beta}_2 = -1{,}883\), \(p < {,}001\); \(N = 1.854\)) so với phi
chế tạo (\(\hat{\beta}_1 = +1{,}615\), \(p = {,}064\);
\(\hat{\beta}_2 = -2{,}479\), \(p = {,}046\); \(N = 1.104\)); quan trọng
hơn, các kênh năng lực và số hóa vận hành chủ yếu trong chế biến chế
tạo, nơi \(\tilde{\mathrm{TCI}}\) trực tiếp \(= +0{,}223\)
(\(p < {,}001\)) và \(\tilde{\mathrm{DAI}}\) trực tiếp \(= +0{,}087\)
(\(p = {,}009\)), trong khi ở phi chế tạo \(\tilde{\mathrm{TCI}}\) chỉ ở
biên (\(+0{,}090\), \(p = {,}096\)) và \(\tilde{\mathrm{DAI}}\) không có
ý nghĩa (\(+0{,}068\), \(p = {,}133\)). Thứ hai, kiểm định dạng hàm bằng
mô hình bậc ba bác bỏ đường cong chữ S: số hạng bậc ba không có ý nghĩa
(\(\hat{\beta} = -1{,}763\), \(p = {,}287\)), và cả AIC lẫn BIC đều ưu
tiên mô hình bậc hai (AIC giảm \(0{,}99\), BIC giảm \(7{,}0\) khi quay
về mô hình bậc hai). Thứ ba, loại doanh nghiệp siêu nhỏ (giữ
\(\ell_1 \geq 10\); \(N = 2.473\)) bảo toàn chữ U ngược
(\(\hat{\beta}_1 = +0{,}794\), \(p = {,}008\);
\(\hat{\beta}_2 = -1{,}647\), \(p < {,}001\)) và hiệu ứng TCI
(\(+0{,}188\), \(p < {,}001\)). Thứ tư, kiểm định giới hạn ổn định hệ số
của Oster (2019) với giả định
\(R^2_{\max} = 1{,}3 \times R^2_{\text{controlled}}\) cho thấy không hệ
số trọng tâm nào đổi dấu hay sụp về 0 dưới các độ lớn hợp lý của chọn
lựa không quan sát được, loại trừ khả năng thiên lệch biến bỏ sót giải
thích phát hiện chính. Bộ kiểm định này cộng hưởng với hệ số TCI 2SLS
dương mạnh (\(1{,}639\), lớn hơn ước lượng OLS, nhất quán với hiệu chỉnh
thiên lệch suy giảm do sai số đo) để xác lập rằng năng lực công nghệ là
cơ chế năng suất bền vững nhất trong bối cảnh thể chế Việt Nam, trong
khi hiện diện website mất sức phân biệt nhân quả.

\hypertarget{bux1ed1i-cux1ea3nh-vux129-muxf4-ux111ux1ed3ng-huxe0nh-vuxe0-ux111ux1ecbnh-vux1ecb-thiux1ebft-kux1ebf-ux111ux1a1n-quux1ed1c-gia-theo-thux1eddi-gian}{%
\subsection{4.5.11 Bối cảnh vĩ mô đồng hành và định vị thiết kế đơn quốc
gia theo thời
gian}\label{bux1ed1i-cux1ea3nh-vux129-muxf4-ux111ux1ed3ng-huxe0nh-vuxe0-ux111ux1ecbnh-vux1ecb-thiux1ebft-kux1ebf-ux111ux1a1n-quux1ed1c-gia-theo-thux1eddi-gian}}

Tính bền vững của dải điểm uốn 39--46\% qua 14 năm trở nên thuyết phục
hơn khi đối chiếu với bức tranh vĩ mô đồng hành ba làn sóng. Tại thời
điểm cuối cửa sổ quan sát, tăng trưởng GDP Việt Nam đạt \(7{,}09\%\) năm
2024, dòng vốn đầu tư trực tiếp nước ngoài đạt \(20{,}17\) tỷ đô-la Mỹ
(tương đương \(4{,}23\%\) GDP), và cấu trúc xuất khẩu dịch chuyển mạnh
sang máy móc và điện tử ở mức \(45{,}8\%\) tổng giá trị thương mại, soán
ngôi cụm dệt may từng thống trị (World Bank, 2025b). Sự nâng cấp cấu
trúc sản xuất này diễn ra song hành với ba đợt khảo sát, củng cố cách
đọc rằng dải điểm uốn ổn định không phải do nền kinh tế tĩnh mà do các
ràng buộc cấp doanh nghiệp (năng lực điều phối, thực thi hợp đồng, tiếp
cận tài trợ thương mại) thay đổi chậm hơn môi trường vận hành rộng. Về
định vị thiết kế, tài liệu phân tích tổng hợp gần đây cho thấy quan hệ
chữ U ngược giữa quốc tế hóa và hiệu quả suy yếu đáng kể, đôi khi biến
mất, dưới nhận dạng nghiêm ngặt hơn trong các mẫu gộp xuyên quốc gia
(Pisani và cộng sự, 2020), gợi ý rằng việc gộp các bối cảnh thể chế
không đồng nhất thổi phồng tính quy luật biểu kiến của dạng hàm. Thiết
kế dọc theo thời gian trong nội bộ một quốc gia, giữ cố định bối cảnh
thể chế trong khi thay đổi thời gian, do đó cung cấp kiểm định khả tín
hơn về việc liệu chữ U ngược là tạo tác của gộp xuyên quốc gia hay một
quan hệ thực bền vững khi môi trường thể chế tiến hóa, và bảng ba đợt
WBES Việt Nam (2009, 2015, 2023) đáp ứng đúng yêu cầu thiết kế này. Cách
đọc này định vị đóng góp của P3 vượt khỏi một báo cáo điểm uốn đơn lẻ:
nó chuyển trọng tâm tranh luận từ ``quốc tế hóa có giúp ích không?''
sang ``dưới điều kiện thể chế nào thì ngưỡng dịch chuyển?''.

\begin{center}\rule{0.5\linewidth}{0.5pt}\end{center}

\hypertarget{kux1ebft-quux1ea3-kiux1ec3m-ux111ux1ecbnh-touxe0n-mux1eabu-ux111a-bux1ed1i-cux1ea3nh-chuxe2u-uxe1-p7}{%
\section{4.6 Kết quả kiểm định toàn mẫu đa bối cảnh châu Á
(P7)}\label{kux1ebft-quux1ea3-kiux1ec3m-ux111ux1ecbnh-touxe0n-mux1eabu-ux111a-bux1ed1i-cux1ea3nh-chuxe2u-uxe1-p7}}

\hypertarget{mux1eabu-vuxe0-muxf4-huxecnh-chuxednh}{%
\subsection{4.6.1 Mẫu và mô hình
chính}\label{mux1eabu-vuxe0-muxf4-huxecnh-chuxednh}}

Phân tích P7 là kiểm định quy mô lớn nhất của luận án, kết hợp tất cả
các bối cảnh ICRV trong khu vực châu Á và Thái Bình Dương trên khung
\textbf{50 nền kinh tế} (gồm Nhật Bản đợt khảo sát đầu tiên năm 2025),
103 cặp nền kinh tế và năm, với \textbf{\(N = 81.022\) quan sát} ở mô
hình bậc hai cơ sở (M2) và \(N = 79.080\) khi bổ sung đầy đủ biến kiểm
soát (M5). Biến phụ thuộc là năng suất lao động chuẩn hóa z trong từng
cặp nền kinh tế và năm (trung hòa tiền tệ, winsorize 1/99 trong đợt);
mọi mô hình có hiệu ứng cố định nền kinh tế và năm, sai số chuẩn cụm
theo nền kinh tế. Mô hình M5, cấu hình đầy đủ với các biến kiểm soát sở
hữu nước ngoài, tuổi doanh nghiệp, TCI và DAI, ước lượng điểm uốn gộp:

\[\text{TP(M5)} = 43{,}6\% \text{ FSTS} \quad (p_{\text{LM}} < {,}001)\]

Điểm uốn dao động nhất quán trong dải \textbf{43,6--51,5\%} FSTS qua các
mô hình M2 đến M10, với kiểm định Lind và Mehlum có ý nghĩa thống kê cao
ở mọi cấu hình. Tính nhất quán này bác bỏ giả thuyết rằng chữ U ngược
chỉ là sản phẩm phụ của mẫu nhỏ hay thiếu biến kiểm soát.

Bảng 4.6 trình bày tiến triển các mô hình lồng nhau, từ mô hình bậc hai
cơ sở đến các mô hình điều tiết, cho thấy điểm uốn ổn định và quan hệ
chữ U ngược được xác nhận qua mọi cấu hình.

\textbf{Bảng 4.6.} \emph{Tiến triển mô hình của P7 trên khung 50 nền
kinh tế: điểm uốn và kiểm định Lind và Mehlum.}

\begin{longtable}[]{@{}
  >{\raggedright\arraybackslash}p{(\columnwidth - 12\tabcolsep) * \real{0.1429}}
  >{\raggedleft\arraybackslash}p{(\columnwidth - 12\tabcolsep) * \real{0.1429}}
  >{\raggedright\arraybackslash}p{(\columnwidth - 12\tabcolsep) * \real{0.1429}}
  >{\raggedright\arraybackslash}p{(\columnwidth - 12\tabcolsep) * \real{0.1429}}
  >{\raggedleft\arraybackslash}p{(\columnwidth - 12\tabcolsep) * \real{0.1429}}
  >{\raggedleft\arraybackslash}p{(\columnwidth - 12\tabcolsep) * \real{0.1429}}
  >{\raggedright\arraybackslash}p{(\columnwidth - 12\tabcolsep) * \real{0.1429}}@{}}
\toprule\noalign{}
\begin{minipage}[b]{\linewidth}\raggedright
Mô hình
\end{minipage} & \begin{minipage}[b]{\linewidth}\raggedleft
N
\end{minipage} & \begin{minipage}[b]{\linewidth}\raggedright
β₁ (p)
\end{minipage} & \begin{minipage}[b]{\linewidth}\raggedright
β₂ (p)
\end{minipage} & \begin{minipage}[b]{\linewidth}\raggedleft
Điểm uốn
\end{minipage} & \begin{minipage}[b]{\linewidth}\raggedleft
\(p_{\text{LM}}\)
\end{minipage} & \begin{minipage}[b]{\linewidth}\raggedright
Mô hình bổ sung
\end{minipage} \\
\midrule\noalign{}
\endhead
\bottomrule\noalign{}
\endlastfoot
M2 & 81.022 & +1,189 (\textless,001) & −1,398 (\textless,001) & 51,5\% &
\textless,001 & FSTS + FSTS², FE hai chiều \\
\textbf{M5} & \textbf{79.080} & \textbf{+0,500 (\textless,001)} &
\textbf{−0,721 (\textless,001)} & \textbf{43,6\%} &
\textbf{\textless,001} & \textbf{+ kiểm soát (mô hình chính)} \\
M7 & 79.080 & +0,627 (\textless,001) & −0,836 (\textless,001) & 46,4\% &
\textless,001 & + TCI và tương tác TCI \\
M8 & 79.080 & +0,704 (\textless,001) & −0,924 (\textless,001) & 47,1\% &
\textless,001 & + DAI và tương tác DAI \\
M10 & 79.080 & +0,615 (\textless,001) & −0,833 (\textless,001) & 45,9\%
& \textless,001 & + tương tác chế độ thể chế yếu \\
\end{longtable}

\emph{Nguồn: ước lượng của tác giả trên khung 50 nền (tái lập:
\texttt{scripts/p7\_run\_50econ.py}). Biến kiểm soát: sở hữu nước ngoài
≥10\%, ln(tuổi doanh nghiệp), TCI chuẩn hóa, DAI. Điểm uốn xác nhận bằng
kiểm định Lind và Mehlum (2010), một phía. Sai số chuẩn cụm theo nền
kinh tế.}

\textbf{Lưu ý xác thực.} Bộ kết quả này được ước lượng bằng công cụ hồi
quy hiệu ứng cố định tương đương reghdfe (pyfixest) với spec được tài
liệu hóa đầy đủ và tái lập được; bộ do-file Stata tương ứng
(\texttt{scripts/stata/}) cho phép xác thực lại từng hệ số trên máy có
Stata trước khi nộp tạp chí.

\hypertarget{ux111iux1ec3m-uux1ed1n-theo-chux1ebf-ux111ux1ed9-thux1ec3-chux1ebf-tuxe1i-ux111ux1ecbnh-huxecnh-h5}{%
\subsection{4.6.2 Điểm uốn theo chế độ thể chế: tái định hình
H5}\label{ux111iux1ec3m-uux1ed1n-theo-chux1ebf-ux111ux1ed9-thux1ec3-chux1ebf-tuxe1i-ux111ux1ecbnh-huxecnh-h5}}

Ước lượng dạng M5 trong từng nhóm ICRV cho thấy một bức tranh điều tiết
thể chế \textbf{tinh tế hơn} giả thuyết ban đầu:

\textbf{Bảng 4.7.} \emph{Điểm uốn theo nhóm ICRV (khung 50 nền, dạng M5
trong từng nhóm).}

\begin{longtable}[]{@{}
  >{\raggedright\arraybackslash}p{(\columnwidth - 12\tabcolsep) * \real{0.1429}}
  >{\raggedleft\arraybackslash}p{(\columnwidth - 12\tabcolsep) * \real{0.1429}}
  >{\raggedright\arraybackslash}p{(\columnwidth - 12\tabcolsep) * \real{0.1429}}
  >{\raggedright\arraybackslash}p{(\columnwidth - 12\tabcolsep) * \real{0.1429}}
  >{\raggedleft\arraybackslash}p{(\columnwidth - 12\tabcolsep) * \real{0.1429}}
  >{\raggedleft\arraybackslash}p{(\columnwidth - 12\tabcolsep) * \real{0.1429}}
  >{\raggedright\arraybackslash}p{(\columnwidth - 12\tabcolsep) * \real{0.1429}}@{}}
\toprule\noalign{}
\begin{minipage}[b]{\linewidth}\raggedright
Nhóm ICRV
\end{minipage} & \begin{minipage}[b]{\linewidth}\raggedleft
N
\end{minipage} & \begin{minipage}[b]{\linewidth}\raggedright
β₁ (p)
\end{minipage} & \begin{minipage}[b]{\linewidth}\raggedright
β₂ (p)
\end{minipage} & \begin{minipage}[b]{\linewidth}\raggedleft
Điểm uốn
\end{minipage} & \begin{minipage}[b]{\linewidth}\raggedleft
\(p_{\text{LM}}\)
\end{minipage} & \begin{minipage}[b]{\linewidth}\raggedright
Kết luận
\end{minipage} \\
\midrule\noalign{}
\endhead
\bottomrule\noalign{}
\endlastfoot
I, Adv. đổi mới (gồm Nhật Bản) & 5.581 & +0,698 (,17) & −0,504 (,31) &
{[}79,1\%{]} & ,30 & gần tuyến tính dương; điểm uốn ngoài miền, không
định vị \\
II, Adv. tài nguyên & 2.075 & +0,854 (,045) & −0,730 (,082) & 62,0\% &
,081 & chữ U ngược ở biên ý nghĩa \\
III, Trung bình cao & 12.055 & +0,168 (,42) & −0,189 (,48) &
{[}55,0\%{]} & ,26 & không xác định ở cấp nhóm \\
\textbf{IV, Chuyển đổi TB-thấp} & \textbf{42.094} & \textbf{+0,709
(\textless,001)} & \textbf{−1,012 (\textless,001)} & \textbf{43,0\%} &
\textbf{\textless,001} & \textbf{chữ U ngược rõ nét nhất} \\
V, Đang nổi & 15.457 & +0,218 (,44) & −0,455 (,21) & {[}34,8\%{]} & ,18
& quan hệ mất cấu trúc \\
VI, SIDS & 1.818 & −0,728 (,010) & +0,870 (,013) & n/a & n/a & không có
chữ U ngược \\
\end{longtable}

\emph{Ghi chú: giá trị trong ngoặc vuông là điểm uốn đại số khi đường
cong không được xác nhận thống kê, chỉ mang tính tham khảo. Nguồn: như
Bảng 4.6.}

Kết quả này \textbf{tái định hình H5} theo hướng giàu thông tin hơn: thể
chế không chỉ dịch chuyển vị trí điểm uốn mà quyết định \textbf{chính sự
tồn tại của dạng hàm chữ U ngược}. Ba vùng được nhận diện rõ. Ở
\textbf{vùng giữa của phổ thể chế} (Nhóm IV), chữ U ngược định hình rõ
nét và tin cậy nhất, với điểm uốn 43,0\% trùng khớp với các ước lượng
đơn quốc gia độc lập trong cùng nhóm (P3 Việt Nam 39--46\%; P9 Ấn Độ
2022 ở mức 40,7\%). Ở \textbf{biên trên của phổ thể chế} (Nhóm I, nay
gồm cả Nhật Bản), đường cong duỗi thẳng thành quan hệ gần tuyến tính
dương với điểm uốn nằm ngoài miền dữ liệu và không định vị được, trùng
khớp với phát hiện đơn quốc gia tại Singapore (P4: gần tuyến tính, điểm
uốn ước lượng \textasciitilde88,6\% không định vị chắc chắn). Ở
\textbf{biên dưới của phổ thể chế} (Nhóm V và VI), quan hệ mất cấu trúc:
các hệ số không còn ý nghĩa thống kê ở Nhóm V và Nhóm VI không có dạng
chữ U ngược, nhất quán với việc tương quan FSTS và năng suất mất ý nghĩa
ở SIDS trong phân tích mô tả (Chuyên đề 1, Bảng 2.3.8.1) và với phát
hiện FIP của P8.

Tương tác chế độ thể chế yếu trong mô hình gộp củng cố cách đọc này: hệ
số FSTS×(chế độ yếu) = −0,523 (\(p = {,}087\)) cho thấy độ dốc nhánh đi
lên của đường cong yếu hơn rõ rệt ở các chế độ thể chế yếu, nghĩa là
cùng một mức tăng cường độ xuất khẩu mang lại phần thưởng năng suất nhỏ
hơn khi đệm thể chế mỏng. Kiểm định điều tiết tổng thể về ICRV của phân
tích tổng hợp (\(Q_M = 17{,}35\), \(p = {,}002\), P6) cung cấp nguồn xác
nhận độc lập thứ hai rằng chế độ thể chế là biến điều tiết có ý nghĩa
của quan hệ quốc tế hóa và hiệu quả.

\begin{figure}
\centering
\includegraphics{figures/fig_4_x_ip_curves_icrv_gradient.png}
\caption{Hình 4.8. Họ đường cong quốc tế hóa và hiệu quả phụ thuộc chế
độ thể chế ICRV}
\end{figure}

\textbf{Hình 4.8.} \emph{Họ đường cong quốc tế hóa và hiệu quả phụ thuộc
chế độ thể chế ICRV.} Chữ U ngược định hình rõ nét nhất ở vùng giữa phổ
thể chế (Nhóm IV, điểm uốn 43,0\%); ở biên trên (Nhóm I) đường cong duỗi
thẳng gần tuyến tính dương; ở biên dưới (Nhóm V và VI) quan hệ mất cấu
trúc, riêng SIDS chuyển sang quan hệ âm theo mức (FIP, H1b, theo spec
mức của P8). \emph{Nguồn: minh họa của tác giả từ kết quả ước lượng 50
nền và P4, P8.}

\hypertarget{ux111iux1ec1u-tiux1ebft-bux1edfi-tci-vuxe0-dai-trong-touxe0n-mux1eabu}{%
\subsection{4.6.3 Điều tiết bởi TCI và DAI trong toàn
mẫu}\label{ux111iux1ec1u-tiux1ebft-bux1edfi-tci-vuxe0-dai-trong-touxe0n-mux1eabu}}

\begin{itemize}
\tightlist
\item
  \textbf{TCI}: hiệu ứng chính dương và có ý nghĩa cao,
  \(\hat{\beta}(\text{TCI}) = +0{,}141\) (\(p < {,}001\)): một độ lệch
  chuẩn TCI nâng năng suất chuẩn hóa thêm 0,141 độ lệch chuẩn. Các tương
  tác FSTS×TCI và FSTS²×TCI không có ý nghĩa thống kê, xác nhận TCI là
  \textbf{năng lực nâng mặt bằng} chứ không uốn hình dạng đường cong,
  nhất quán với kết luận từ các nghiên cứu đơn quốc gia.
\item
  \textbf{DAI}: hiệu ứng chính dương và có ý nghĩa cao,
  \(\hat{\beta}(\text{DAI}) = +0{,}201\) (\(p < {,}001\)): doanh nghiệp
  có hiện diện số cơ bản đạt năng suất chuẩn hóa cao hơn 0,2 độ lệch
  chuẩn. Các tương tác với đường cong (FSTS×DAI = −0,271,
  \(p = {,}149\); FSTS²×DAI = +0,252, \(p = {,}367\)) mang \textbf{cùng
  chiều nén đường cong} như dự đoán của khung CDCM nhưng không đạt ý
  nghĩa thống kê trên toàn mẫu 50 nền; bằng chứng điều tiết DAI rõ nét
  nhất vẫn là ở cấp bối cảnh đơn quốc gia (P4 Singapore: FSTS²×DAI =
  +3,119, \(p = {,}005\)). H3 do đó được hỗ trợ ở hiệu ứng mức và ở bối
  cảnh đặc thù, chưa được xác nhận như hiệu ứng uốn đường cong phổ quát.
  Khoảng cách độ lớn giữa hai ước lượng (3,119 ở Singapore so với 0,252
  ở mẫu gộp) không phải mâu thuẫn mà phản ánh bản chất trung bình hóa
  của mẫu gộp: phần lớn 50 nền chỉ đo được DAI ở Tầng 1 (hiện diện web),
  nơi công cụ số chưa đủ chiều sâu để uốn đường cong, trong khi hiệu ứng
  uốn mạnh ở Singapore gắn với năng lực số Tầng 2 hiếm có trong mẫu; gộp
  hai loại bối cảnh này lại làm pha loãng hiệu ứng uốn về một trung bình
  không có ý nghĩa, đúng như logic phụ thuộc bối cảnh của khung CDCM dự
  báo chứ không phủ định cơ chế lá chắn số.
\end{itemize}

\hypertarget{diux1ec5n-tiux1ebfn-ux111iux1ec3m-uux1ed1n-vuxe0-bux1ed1n-ux111iux1ec1u-kiux1ec7n-nhux1eadn-diux1ec7n-chux1eef-u-ngux1b0ux1ee3c}{%
\subsection{4.6.4 Diễn tiến điểm uốn và bốn điều kiện nhận diện chữ U
ngược}\label{diux1ec5n-tiux1ebfn-ux111iux1ec3m-uux1ed1n-vuxe0-bux1ed1n-ux111iux1ec1u-kiux1ec7n-nhux1eadn-diux1ec7n-chux1eef-u-ngux1b0ux1ee3c}}

Khi đọc theo chuỗi mô hình lồng nhau, điểm uốn dịch chuyển có quy luật
chứ không ngẫu nhiên. Mô hình tuyến tính cơ sở cho hệ số FSTS dương và
có ý nghĩa cao (\(\hat{\beta} = +0{,}219\), \(p < {,}001\)), xác nhận
rằng ở mức trung bình toàn mẫu quan hệ quốc tế hóa và năng suất là
dương. Khi bổ sung số hạng bậc hai (mô hình M2), đường cong định hình
chữ U ngược với điểm uốn ở mức 51,5\% FSTS. Việc đưa vào lần lượt các
biến kiểm soát nhân khẩu doanh nghiệp rồi khối năng lực (TCI và DAI) kéo
điểm uốn xuống 43,8\% rồi 43,6\%, chính là điểm uốn của mô hình chính
(M5). Sự dịch chuyển xuống khoảng bảy điểm phần trăm này có ý nghĩa thực
chất: một phần độ cong biểu kiến ở mô hình thô phản ánh chênh lệch năng
lực giữa các doanh nghiệp xuất khẩu nhiều và ít; sau khi kiểm soát năng
lực, ngưỡng tối ưu thực của cường độ xuất khẩu thấp hơn và nằm sát các
ước lượng đơn quốc gia trong vùng chuyển đổi (Việt Nam 39--46\%, Ấn Độ
2022 ở mức 40,7\%).

Chữ U ngược được xác nhận chính thức theo bốn điều kiện của Haans và
cộng sự (2016) chứ không chỉ dựa vào dấu của hệ số bậc hai. Thứ nhất, hệ
số nhánh đi lên dương (\(\beta_1 > 0\)). Thứ hai, hệ số nhánh đi xuống
âm (\(\beta_2 < 0\)). Thứ ba, điểm uốn nằm hẳn trong miền dữ liệu quan
sát của FSTS. Thứ tư, độ dốc ở hai biên thỏa mãn dự đoán hướng, dương ở
biên dưới và âm ở biên trên. Kiểm định Lind và Mehlum (2010), vốn kiểm
tra đồng thời các điều kiện này, cho giá trị \(p_{\text{LM}} < {,}001\)
ở mô hình chính, nghĩa là chữ U ngược được nhận diện chắc chắn chứ không
phải tạo tác của một quan hệ đơn điệu bị ép vào dạng bậc hai.

\hypertarget{tuxednh-ux1ed5n-ux111ux1ecbnh-cux1ee7a-ux111iux1ec3m-uux1ed1n-trux1b0ux1edbc-ux111ux1eb7c-trux1b0ng-quux1ea3n-luxfd-vuxe0-tux1b0ux1a1ng-tuxe1c-thux1ec3-chux1ebf}{%
\subsection{4.6.5 Tính ổn định của điểm uốn trước đặc trưng quản lý và
tương tác thể
chế}\label{tuxednh-ux1ed5n-ux111ux1ecbnh-cux1ee7a-ux111iux1ec3m-uux1ed1n-trux1b0ux1edbc-ux111ux1eb7c-trux1b0ng-quux1ea3n-luxfd-vuxe0-tux1b0ux1a1ng-tuxe1c-thux1ec3-chux1ebf}}

Điểm uốn gộp tỏ ra trơ trước việc bổ sung các đặc trưng quản lý. Trong
mô hình có biến quản lý (\(N = 75.029\)), kinh nghiệm ngành của người
quản lý không có ý nghĩa thống kê (\(\hat{\beta} = +0{,}002\),
\(p > {,}10\)), còn hệ số người quản lý cấp cao là nữ mang dấu âm và có
ý nghĩa cao (\(\hat{\beta} = -0{,}104\), \(p < {,}001\)); điểm uốn gần
như không đổi ở mức 44,0\%. Dấu âm của hệ số quản lý nữ nhất quán với
phần còn lại của luận án (nghiên cứu đơn quốc gia Nhật Bản P10 và phân
tích gộp đều ước lượng hệ số âm) và được đọc thận trọng như một quy luật
mô tả phản ánh chọn lọc cấu thành, theo đó các cơ sở do nữ lãnh đạo tập
trung nhiều hơn vào các phân khúc nhỏ, biên lợi nhuận thấp ở những nền
kinh tế mà lãnh đạo nữ cấp cao còn hiếm, chứ không phải bằng chứng nhân
quả về hiệu quả của lãnh đạo nữ. Vì các tương tác giữa FSTS và đặc trưng
quản lý không có ý nghĩa, đặc trưng quản lý dịch chuyển mặt bằng năng
suất chứ không giúp doanh nghiệp đi qua đánh đổi quốc tế hóa hiệu quả
hơn.

Tương tự, chế độ thể chế ICRV không biểu hiện như một dịch chuyển tuyến
tính đơn của đường cong gộp. Vì biến ICRV bất biến theo thời gian trong
mỗi nền kinh tế, hiệu ứng chính của nó bị hấp thụ bởi hiệu ứng cố định
nền kinh tế; chỉ các số hạng tương tác được nhận diện. Trong mô hình
tương tác thể chế gộp, cả FSTS×ICRV (\(\hat{\beta} = -0{,}086\)) lẫn
FSTS²×ICRV (\(\hat{\beta} = -0{,}013\)) đều không có ý nghĩa, khiến điểm
uốn gộp không định vị được. Đây không phải bằng chứng phủ định H5 mà là
chỉ dấu phương pháp luận: cấu trúc thể chế của quan hệ chỉ hiện ra khi
ước lượng điểm uốn trong từng nhóm ICRV (Bảng 4.7), nơi nó biểu hiện
thành ba vùng rời rạc chứ không phải một phổ liên tục. Nói cách khác,
điều tiết thể chế là điều tiết về sự tồn tại của dạng hàm, không phải
điều tiết về vị trí của một điểm uốn phổ quát.

\hypertarget{ux111ux1eb7c-trux1b0ng-mux1eabu-ux111a-cux1ed9ng-tuyux1ebfn-vuxe0-uxfd-nghux129a-tux1ed5ng-hux1ee3p-ba-vuxf9ng}{%
\subsection{4.6.6 Đặc trưng mẫu, đa cộng tuyến và ý nghĩa tổng hợp ba
vùng}\label{ux111ux1eb7c-trux1b0ng-mux1eabu-ux111a-cux1ed9ng-tuyux1ebfn-vuxe0-uxfd-nghux129a-tux1ed5ng-hux1ee3p-ba-vuxf9ng}}

Cường độ xuất khẩu trung bình toàn mẫu là 9,0\% FSTS, phản ánh một quần
thể doanh nghiệp thiên về thị trường nội địa với một đuôi xuất khẩu
mạnh, đặc điểm điển hình của dữ liệu doanh nghiệp WBES ở các nền kinh tế
đang phát triển và mới nổi. Để giảm đa cộng tuyến giữa số hạng tuyến
tính và bậc hai, FSTS được căn giữa quanh trung bình mẫu trước khi lập
số hạng bình phương và các số hạng tương tác; nhờ vậy hệ số nhánh đi lên
và nhánh đi xuống được ước lượng ổn định qua toàn bộ chuỗi mô hình. Mọi
mô hình mang hiệu ứng cố định nền kinh tế và năm cùng sai số chuẩn cụm
theo nền kinh tế (CRV1), kiểm soát đồng thời chênh lệch mặt bằng giá và
năng suất giữa các nền kinh tế và xu hướng theo thời gian.

Tổng hợp lại, ba mức bằng chứng, gồm diễn tiến điểm uốn ổn định trong
dải 43,6--51,5\%, tính trơ của điểm uốn trước năng lực và đặc trưng quản
lý, và cấu trúc ba vùng theo chế độ thể chế, hội tụ về một cách đọc
thống nhất: năng lực công nghệ và số hóa cùng đặc trưng quản lý nâng
\emph{mặt bằng} năng suất nhưng không uốn lại hình dạng đường cong gộp;
cái uốn lại đường cong là chế độ thể chế, và nó tác động không như một
phổ liên tục mà như ba vùng rời rạc. Ở giữa phổ thể chế (vùng chuyển đổi
thu nhập trung bình thấp), chữ U ngược định hình sắc nét quanh điểm uốn
43\%; ở biên trên (thể chế mạnh nhất, gồm Nhật Bản), nhánh đi xuống bị
đẩy ra ngoài miền quan sát và quan hệ gần như tuyến tính dương; ở biên
dưới (đang nổi và SIDS Thái Bình Dương), chữ U ngược mất cấu trúc và đảo
dấu ở nhóm SIDS thành hình phạt quốc tế hóa cưỡng bức. Cách đọc ba vùng
này dung hòa các kết quả đơn quốc gia tưởng như mâu thuẫn của luận án
trong một khung ước lượng hài hòa trên 50 nền kinh tế.

\hypertarget{hux1ed9i-tux1ee5-chuxe9o-nghiuxean-cux1ee9u-cux1ee7a-ux111iux1ec3m-uux1ed1n-vuxf9ng-giux1eefa-vuxe0-logic-chux1ecdn-lux1ecdc-melitz-ux1edf-biuxean-truxean}{%
\subsection{4.6.7 Hội tụ chéo nghiên cứu của điểm uốn vùng giữa và logic
chọn lọc Melitz ở biên
trên}\label{hux1ed9i-tux1ee5-chuxe9o-nghiuxean-cux1ee9u-cux1ee7a-ux111iux1ec3m-uux1ed1n-vuxf9ng-giux1eefa-vuxe0-logic-chux1ecdn-lux1ecdc-melitz-ux1edf-biuxean-truxean}}

Sức thuyết phục của cấu trúc ba vùng nằm ở chỗ ước lượng gộp 50 nền và
các ước lượng đơn quốc gia độc lập hội tụ chặt chẽ ở từng vùng dù được
ước lượng trên các mẫu, công cụ khảo sát và mô hình khác nhau. Ở vùng
giữa, điểm uốn cấp nhóm của Nhóm IV (43,0\%) trùng khớp đáng kể với điểm
uốn đơn quốc gia của Việt Nam, vốn được kiểm định Lind--Mehlum xác nhận
trong cả ba đợt (2009 \(p = {,}006\), 2015 \(p = {,}009\), 2023
\(p = {,}013\)) và mẫu gộp (\(p < {,}001\)) với dải điểm uốn
\(39{,}3\)--\(46{,}2\%\), và với điểm uốn Ấn Độ 2022 ở mức \(40{,}7\%\).
Sự trùng khớp ba ước lượng độc lập quanh mốc 40--43\% ở cùng chế độ thể
chế ICRV IV cung cấp bằng chứng hội tụ mạnh rằng điểm uốn vùng giữa là
một quy luật cấu trúc chứ không phải tạo tác của một mẫu đơn lẻ.

Ở biên trên của phổ thể chế, hình thái gần tuyến tính dương của Nhóm I
được làm rõ về cơ chế nhờ bằng chứng đơn quốc gia. Tại Singapore, nhánh
suy giảm bên phải không được nhận diện vì khoảng \(82\%\) doanh nghiệp
không xuất khẩu và chỉ \(3\%\) vượt ngưỡng \(50\%\) FSTS, khiến điểm uốn
\(\approx 88{,}6\%\) rơi vào vùng dữ liệu thưa với khoảng tin cậy
bootstrap rộng \([53\%; 253\%]\). Cơ chế đứng sau biên trên là chọn lọc
theo năng suất kiểu Melitz (2003): ở chế độ thể chế mạnh nhất, câu hỏi
ràng buộc không phải doanh nghiệp xuất khẩu \emph{bao nhiêu} mà là
\emph{có} xuất khẩu hay không, và các doanh nghiệp vượt ngưỡng tham gia
là những doanh nghiệp được trang bị năng lực, đúng như logic tự chọn lọc
mà mô hình doanh nghiệp dị biệt dự đoán cho biên thâm dụng tri thức
(Helpman và cộng sự, 2004). Ba vùng do đó không chỉ mô tả ba hình thái
đường cong mà ánh xạ ba cơ chế kinh tế khác nhau: chi phí điều phối bám
ngưỡng ở vùng giữa, chọn lọc theo năng suất ở biên trên, và hình phạt
quốc tế hóa cưỡng bức ở biên dưới.

\hypertarget{kiux1ec3m-ux111ux1ecbnh-ux111ux1ed9-vux1eefng-bux1ed5-sung-cux1ee5m-nhux1ecf-ux111a-kiux1ec3m-ux111ux1ecbnh-bux1ea5t-biux1ebfn-ux111o-lux1b0ux1eddng-vuxe0-chux1ecdn-lux1ecdc}{%
\subsection{4.6.8 Kiểm định độ vững bổ sung: cụm nhỏ, đa kiểm định, bất
biến đo lường và chọn
lọc}\label{kiux1ec3m-ux111ux1ecbnh-ux111ux1ed9-vux1eefng-bux1ed5-sung-cux1ee5m-nhux1ecf-ux111a-kiux1ec3m-ux111ux1ecbnh-bux1ea5t-biux1ebfn-ux111o-lux1b0ux1eddng-vuxe0-chux1ecdn-lux1ecdc}}

Bốn kiểm định độ vững bổ sung được thực hiện trực tiếp trên khung 50 nền
để củng cố ranh giới suy luận của kết quả gộp (chi tiết tái lập ở Phụ
lục dữ liệu, tệp \texttt{p7\_robustness\_suite\_2026-06.md}).

\textbf{Thứ nhất, suy luận cụm nhỏ.} Vì ước lượng điểm uốn theo từng
nhóm ICRV chỉ dựa trên 6--16 nền kinh tế (cụm), sai số chuẩn cụm CRV1 có
xu hướng bác bỏ quá mức. Một kiểm định bootstrap cụm hoang dã
(wild-cluster, Cameron--Gelbach--Miller, trọng số Rademacher, B = 9.999,
cụm theo nền kinh tế) cho thấy độ cong cấp nhóm nhạy cảm với số cụm: số
hạng bậc hai của vùng chuyển đổi (Nhóm IV) vẫn mạnh nhất nhưng giá trị
\(p\) nới rộng từ \(0{,}001\) (CRV1) lên khoảng \(0{,}10\), trong khi
các nhóm Đang nổi và SIDS giữ nguyên không có ý nghĩa
(\(p_{\text{wild}} \approx 0{,}33\) và \(0{,}30\)). Kết quả này xác nhận
rằng suy luận hình thức về tính phi tuyến được neo đúng ở \textbf{cấp
gộp} (mô hình M2 với 49 cụm, \(\beta_2 = -1{,}399\), \(p < {,}001\)) và
tam giác hóa với phân tích tổng hợp P6, chứ không đọc ra từ từng nhóm
con ít cụm; các điểm uốn theo nhóm ICRV vì vậy được diễn giải như một
phổ mô tả chứ không phải các kiểm định độc lập đủ công suất.

\textbf{Thứ hai, đa kiểm định.} Áp dụng hiệu chỉnh Holm và Bonferroni
cho cả họ bốn số hạng tương tác năng lực số (FSTS×TCI, FSTS²×TCI,
FSTS×DAI, FSTS²×DAI): không số hạng nào có ý nghĩa ở mức thô (mọi
\(p \ge 0{,}14\)) và tất cả vẫn không có ý nghĩa sau hiệu chỉnh. Kết
luận năng lực là yếu tố nâng mặt bằng chứ không uốn đường cong do đó
không phải một âm tính giả sinh ra từ kiểm định lặp, mà vững trước hiệu
chỉnh bảo thủ nhất.

\textbf{Thứ ba, bất biến đo lường.} Các mục đo của DAI Tầng 1 (hiện diện
web) và TCI hiện diện cho khoảng \(99\%\) doanh nghiệp ở cả ba thế hệ
lược đồ (PICS3, Standardized, BREADY/BEE), nên các cấu trúc hài hòa là
so sánh được giữa các đợt; trung bình DAI tăng đơn điệu qua ba thế hệ
(\(0{,}30 \to 0{,}46 \to 0{,}53\)), nhất quán với khuếch tán số thực
chất chứ không phải đứt gãy lược đồ. Các so sánh chéo đợt (P5 2012↔2024;
P9 2014↔2025) do đó phản ánh thay đổi thực chất trên một mục bất biến.

\textbf{Thứ tư, chọn lọc tham gia xuất khẩu.} Một kiểm tra kiểu Heckman
bổ sung tỷ số Mills nghịch đảo (từ probit tham gia xuất khẩu) vào phương
trình kết quả: độ cong chữ U ngược và điểm uốn gần như không đổi
(\(\text{TP}\) \(43{,}6\% \to 43{,}5\%\); \(\beta_2 = -0{,}732\),
\(p < {,}001\)), và bản thân tỷ số Mills không có ý nghĩa (\(-0{,}414\),
\(p = {,}317\)). Không có bằng chứng về thiên lệch tự chọn lọc kiểu
Melitz trong độ cong ước lượng: chữ U ngược không phải tạo tác của việc
doanh nghiệp năng suất cao tự chọn vào xuất khẩu.

\begin{center}\rule{0.5\linewidth}{0.5pt}\end{center}

\hypertarget{kux1ebft-quux1ea3-trux1b0ux1eddng-hux1ee3p-biuxean-sids-cuxe1c-nux1ec1n-kinh-tux1ebf-ux111ux1ea3o-nhux1ecf-ux111ang-phuxe1t-triux1ec3n-p8}{%
\section{4.7 Kết quả trường hợp biên SIDS: Các nền kinh tế đảo nhỏ đang
phát triển
(P8)}\label{kux1ebft-quux1ea3-trux1b0ux1eddng-hux1ee3p-biuxean-sids-cuxe1c-nux1ec1n-kinh-tux1ebf-ux111ux1ea3o-nhux1ecf-ux111ang-phuxe1t-triux1ec3n-p8}}

\hypertarget{ux111ux1eb7c-ux111iux1ec3m-mux1eabu-cuxe1c-nux1ec1n-kinh-tux1ebf-ux111ux1ea3o-nhux1ecf}{%
\subsection{4.7.1 Đặc điểm mẫu các nền kinh tế đảo
nhỏ}\label{ux111ux1eb7c-ux111iux1ec3m-mux1eabu-cuxe1c-nux1ec1n-kinh-tux1ebf-ux111ux1ea3o-nhux1ecf}}

Phân tích P8 tập trung vào bảy nền kinh tế đảo nhỏ đang phát triển Thái
Bình Dương thuộc Nhóm VI, gồm Fiji, Kiribati, Papua New Guinea, Samoa,
Quần đảo Solomon, Tonga và Vanuatu, với dữ liệu WBES các đợt khảo sát
giai đoạn 2009--2025. Bảy nền này là các quốc đảo Thái Bình Dương theo
phân loại của Ngân hàng Thế giới và danh sách SIDS Thái Bình Dương của
Liên Hợp Quốc; Timor-Leste, dù từng bị một số bản dựng dữ liệu xếp nhầm
vào nhóm, \textbf{không} thuộc nhóm quốc đảo Thái Bình Dương (Ngân hàng
Thế giới xếp là một nền kinh tế Đông Á--Thái Bình Dương riêng) nên được
loại khỏi phân tích. Mẫu phân tích, gồm các quan sát có đủ giá trị về
năng suất lao động và cường độ xuất khẩu, có tổng \(N = 1.450\) doanh
nghiệp trên bảy cụm quốc gia, trong đó 212 doanh nghiệp (14,6\% mẫu) có
xuất khẩu với cường độ dương. Cường độ xuất khẩu trung bình rất thấp, ở
mức 0,054 (tức 5,4\%), phản ánh đúng đặc trưng của nhóm: phần lớn doanh
nghiệp định hướng nội địa, còn một thiểu số xuất khẩu thì làm vậy do bắt
buộc cấu trúc chứ không phải lựa chọn chiến lược.

Ba đặc điểm cấu trúc của các nền này định khung cho toàn bộ phân tích.
Thứ nhất, thị trường nội địa siêu nhỏ, với quy mô dân số từ khoảng
94.000 (Micronesia, ngoài mẫu) và Tonga 104.000, Kiribati 121.000 đến
tối đa khoảng 10 triệu (Papua New Guinea, nền lớn nhất trong mẫu), khiến
nhiều doanh nghiệp không thể đạt quy mô hiệu quả tối thiểu nếu chỉ phục
vụ thị trường trong nước. Thứ hai, sự cô lập địa lý cực độ khuếch đại
chi phí thương mại thêm khoảng 30 đến 210\% so với mặt bằng. Thứ ba, các
khoảng trống thể chế toàn diện hạn chế khả năng doanh nghiệp chiếm giữ
lợi ích từ hoạt động quốc tế. Trong bối cảnh này, cường độ xuất khẩu
không còn là chỉ báo của cam kết tự nguyện với thị trường nước ngoài mà
trở thành đại lượng đại diện cho \emph{mức độ bất cập của thị trường nội
địa}: doanh nghiệp có cường độ xuất khẩu cao hơn không hẳn là doanh
nghiệp có năng lực quốc tế vượt trội, mà có thể chỉ là doanh nghiệp chịu
ràng buộc thị trường nội địa gay gắt hơn. Chính sự đảo nghĩa của thước
đo này là nền tảng cho dự đoán hình phạt quốc tế hóa cưỡng bức.

\hypertarget{tan-ruxe3-chux1eef-u-ngux1b0ux1ee3c-tux1ea1i-cuxe1c-nux1ec1n-kinh-tux1ebf-ux111ux1ea3o-nhux1ecf}{%
\subsection{4.7.2 Tan rã chữ U ngược tại các nền kinh tế đảo
nhỏ}\label{tan-ruxe3-chux1eef-u-ngux1b0ux1ee3c-tux1ea1i-cuxe1c-nux1ec1n-kinh-tux1ebf-ux111ux1ea3o-nhux1ecf}}

Phát hiện trọng tâm và khác biệt nhất của bối cảnh SIDS so với mọi nhóm
ICRV khác là \textbf{việc dạng chữ U ngược mất cấu trúc}. Trên mẫu đầy
đủ bảy nền kinh tế đảo nhỏ Thái Bình Dương (\(N = 1.450\) doanh nghiệp,
bảy cụm quốc gia), ước lượng bằng đúng quy trình hài hòa hóa của P7
(biến phụ thuộc là năng suất lao động chuẩn hóa z trong từng cặp nền ×
năm, hiệu ứng cố định nền và năm) với suy luận few-cluster hợp lệ bằng
bootstrap cụm hoang dã (Cameron, Gelbach và Miller, 2008; B = 9.999),
\textbf{cả độ cong lẫn độ dốc của FSTS đều không có ý nghĩa thống kê}:
độ dốc tuyến tính \(\hat{\beta}(\text{FSTS}) = -0{,}085\)
(\(p_{\text{wild}} = {,}66\)) ở mô hình tuyến tính, và số hạng bậc hai
\(\hat{\beta}(\text{FSTS}^2) = +0{,}696\) chỉ đạt ý nghĩa biên
(\(p_{\text{wild}} = {,}082\)). Các điều kiện Lind--Mehlum cho một chữ U
ngược thực sự không được thỏa, và không tồn tại cực đại nội vùng. Điều
này tương phản sắc nét với chế độ chuyển đổi lục địa (Nhóm IV) nơi cùng
mô hình cho chữ U ngược mạnh và rất có ý nghĩa
(\(\hat{\beta}(\text{FSTS}^2) = -1{,}012\), \(p < {,}001\)):
\textbf{dạng hàm chữ U ngược, vốn tổ chức cả tài liệu I--P, không khái
quát được tới biên cực đoan.}

Một dấu hiệu \textbf{gợi ý} (chỉ ở mức biên ý nghĩa) xuất hiện khi giữ
cố định năng lực doanh nghiệp (mô hình M3): độ cong chuyển dương,
\(\hat{\beta}(\text{FSTS}^2) = +1{,}051\) (\(p_{\text{wild}} = {,}056\))
kèm số hạng tuyến tính âm \(\hat{\beta}(\text{FSTS}) = -0{,}852\)
(\(p_{\text{wild}} = {,}024\)), hàm ý một \textbf{vùng phạt năng suất
trong dải cường độ xuất khẩu thấp đến trung bình} --- đúng vùng mà gần
như toàn bộ doanh nghiệp đảo nhỏ hoạt động --- phù hợp với cơ chế hình
phạt quốc tế hóa cưỡng bức (FIP). Tuy nhiên, do chỉ đạt ý nghĩa biên và
số cụm nhỏ, luận án \textbf{không khẳng định một sự đảo chiều dạng hàm}
mà chỉ ghi nhận xu hướng này một cách thận trọng. Năng lực công nghệ
nâng mặt bằng năng suất (\(\hat{\beta}(\text{TCI}_z) = +0{,}064\),
\(p_{\text{wild}} = {,}036\)) nhưng không khôi phục dạng chữ U ngược.

\begin{figure}
\centering
\includegraphics{figures/fig_4_7_1_sids_inversion.png}
\caption{Hình 4.7.1. Việc chữ U ngược mất cấu trúc ở biên cực đoan: lục
địa (∩, chế độ chuyển đổi, độ cong âm có ý nghĩa) so với Pacific SIDS
(độ cong không có ý nghĩa trên mẫu đầy đủ; chỉ chuyển dương ở mức biên
khi giữ cố định năng lực). Phần lớn doanh nghiệp đảo nhỏ nằm trong dải
cường độ xuất khẩu thấp đến trung bình. Nguồn: ước lượng của tác giả từ
WBES.}
\end{figure}

Trên một bản dựng dữ liệu \textbf{hạn chế hơn} (chỉ ba nền có đủ biến
kiểm soát đợt trước 2018: Fiji 2009, Papua New Guinea 2015, Vanuatu
2009; \(N = 209\)), độ dốc tuyến tính âm mạnh và có ý nghĩa,
\(\hat{\beta}(\text{FSTS}_c) = -1{,}339\) (\(SE = 0{,}386\),
\(p < {,}001\)) --- chữ ký FIP đậm nét nhất. Tuy nhiên kết quả này dựa
trên \textbf{ba cụm} và \textbf{nhạy cảm với phiên bản dữ liệu} (xem Mục
4.7.8): trên mẫu đầy đủ, dấu và độ lớn không ổn định. Do đó hệ số −1,339
được đọc như bằng chứng minh họa cho \textbf{trường hợp giới hạn lý
thuyết FIP} chứ không phải một hiệu ứng được ước lượng vững, và phát
hiện vững của bối cảnh SIDS là \textbf{việc chữ U ngược mất cấu trúc};
độ cong dương có điều kiện chỉ là một dấu hiệu gợi ý ở mức biên ý nghĩa.

\hypertarget{giux1ea3i-thuxedch-cux1a1-chux1ebf-huxecnh-phux1ea1t-quux1ed1c-tux1ebf-huxf3a-cux1b0ux1ee1ng-bux1ee9c}{%
\subsection{4.7.3 Giải thích cơ chế hình phạt quốc tế hóa cưỡng
bức}\label{giux1ea3i-thuxedch-cux1a1-chux1ebf-huxecnh-phux1ea1t-quux1ed1c-tux1ebf-huxf3a-cux1b0ux1ee1ng-bux1ee9c}}

Luận án đặt tên cho hiện tượng này là \textbf{Forced
Internationalization Penalty (FIP)}, gánh nặng quốc tế hóa bắt buộc, một
cấu trúc lý thuyết mới chưa có trong tài liệu IB trước đây. Cơ chế giải
thích gồm ba yếu tố cấu trúc đặc thù của SIDS: (1) thị trường nội địa
quá nhỏ, buộc doanh nghiệp SIDS xuất khẩu không phải vì có lợi thế cạnh
tranh mà vì nhu cầu nội địa không đủ để duy trì hoạt động; (2) chi phí
hậu cần và tiếp cận thị trường cực cao do vị trí địa lý xa xôi và phân
tán; và (3) thiếu hỗ trợ thể chế và năng lực xuất khẩu.

FIP là một đóng góp \textbf{khái niệm} của luận án, lần đầu được đặt tên
trong tài liệu quốc tế hóa--hiệu quả hoạt động kinh doanh của doanh
nghiệp cho khu vực Thái Bình Dương; bằng chứng thực nghiệm đi kèm mang
tính \textbf{thăm dò} và phải được đọc cùng giới hạn dữ liệu trình bày ở
Mục 4.7.8.

\hypertarget{ux111ux1ecbnh-vux1ecb-fip-trong-bux1eb1ng-chux1ee9ng-vux129-muxf4-vux1ec1-tux1ed5n-thux1b0ux1a1ng-cux1ee7a-nux1ec1n-kinh-tux1ebf-nhux1ecf}{%
\subsection{4.7.4 Định vị FIP trong bằng chứng vĩ mô về tổn thương của
nền kinh tế
nhỏ}\label{ux111ux1ecbnh-vux1ecb-fip-trong-bux1eb1ng-chux1ee9ng-vux129-muxf4-vux1ec1-tux1ed5n-thux1b0ux1a1ng-cux1ee7a-nux1ec1n-kinh-tux1ebf-nhux1ecf}}

Cơ chế FIP ở cấp doanh nghiệp ăn khớp với đặc trưng tổn thương cấu trúc
của các nền kinh tế nhỏ và thị trường biên được ghi nhận ở cấp vĩ mô.
Bằng chứng của Ngân hàng Thế giới cho thấy bản chất phơi nhiễm bên ngoài
của các thị trường biên là bất lợi so với thị trường mới nổi: dòng vốn
danh mục biến động hơn, rổ xuất khẩu tập trung hơn, và nợ chủ yếu được
định danh bằng ngoại tệ với mức dự trữ thấp, tất cả làm tăng rủi ro bảng
cân đối (World Bank, 2026c). Đối với các nền kinh tế quy mô siêu nhỏ như
SIDS Thái Bình Dương, những bất lợi này được khuếch đại bởi chi phí
logistics cực cao và sự tập trung xuất khẩu cao độ, khiến việc gia tăng
cường độ xuất khẩu không đi kèm phần thưởng năng suất mà thậm chí trở
thành gánh nặng. Như vậy, FIP không phải một ngoại lệ thống kê mà là
biểu hiện cấp doanh nghiệp của một quy luật cấu trúc rộng hơn: ở những
bối cảnh thị trường nội địa quá nhỏ và đệm thể chế quá mỏng, quốc tế hóa
mang tính bắt buộc chứ không phải lựa chọn chiến lược tối ưu hóa năng
suất. Phát hiện này bổ sung một viên gạch vi mô còn thiếu vào bức tranh
vĩ mô của Ngân hàng Thế giới về các nền kinh tế nhỏ và dễ tổn thương.

\hypertarget{cux1ea5u-truxfac-mux1eabu-bux1ea3y-nux1ec1n-kinh-tux1ebf-vuxe0-phux1b0ux1a1ng-sai-giux1eefa-quux1ed1c-gia}{%
\subsection{4.7.5 Cấu trúc mẫu bảy nền kinh tế và phương sai giữa quốc
gia}\label{cux1ea5u-truxfac-mux1eabu-bux1ea3y-nux1ec1n-kinh-tux1ebf-vuxe0-phux1b0ux1a1ng-sai-giux1eefa-quux1ed1c-gia}}

Cấu trúc bảy nền kinh tế cho thấy mức độ phơi nhiễm xuất khẩu rất mỏng
và không đồng đều: tỷ lệ doanh nghiệp xuất khẩu dao động từ \(2{,}6\%\)
(Kiribati, FSTS trung bình \(0{,}011\)) đến \(27{,}9\%\) (Fiji, FSTS
trung bình \(0{,}109\)), với mức gộp toàn mẫu là \(13{,}8\%\) và FSTS
trung bình \(0{,}048\). Chỉ ba nền kinh tế (Fiji, Papua New Guinea,
Vanuatu) có hai đợt khảo sát WBES, phần còn lại chỉ có một đợt, nên nhận
dạng chủ yếu dựa trên biến thiên chéo trong từng ô quốc gia--năm. Hiệu
ứng cố định quốc gia hấp thụ một phần phương sai rất lớn: \(R^2\) tăng
từ \(0{,}109\) (chỉ hiệu ứng cố định năm) lên \(0{,}657\) khi bổ sung
hiệu ứng cố định quốc gia, nghĩa là khoảng \(83\%\) phương sai năng suất
được giải thích là phương sai giữa các quốc gia, một đặc trưng nhất quán
với tính dị biệt cấu trúc cực đoan giữa các nền kinh tế đảo nhỏ.

\hypertarget{kiux1ec3m-ux111ux1ecbnh-bux1ed1n-ux111iux1ec1u-kiux1ec7n-ux111ux1b0ux1eddng-cong-vuxe0-hux1ec7-sux1ed1-bux1eadc-hai-ux111ux1a1n-ux111iux1ec7u-uxe2m}{%
\subsection{4.7.6 Kiểm định bốn điều kiện đường cong và hệ số bậc hai
đơn điệu
âm}\label{kiux1ec3m-ux111ux1ecbnh-bux1ed1n-ux111iux1ec1u-kiux1ec7n-ux111ux1b0ux1eddng-cong-vuxe0-hux1ec7-sux1ed1-bux1eadc-hai-ux111ux1a1n-ux111iux1ec7u-uxe2m}}

Quan hệ đơn điệu âm được xác nhận qua việc cả bốn điều kiện cần của một
đường cong chữ U ngược (Haans và cộng sự, 2016) đều thất bại, đúng như
chữ ký kỳ vọng của hiện tượng gánh nặng quốc tế hóa bắt buộc. Ở mô hình
bậc hai M2, cả hai số hạng đều mất ý nghĩa khi đưa vào đồng thời:
\(\hat{\beta}(\text{FSTS}_c) = -2{,}177\) (\(SE = 1{,}603\);
\(p = {,}176\)) và \(\hat{\beta}(\text{FSTS}_c^2) = +1{,}125\)
(\(SE = 1{,}943\); \(p = {,}563\)). Điều đáng chú ý là dấu dương của hệ
số bậc hai (nếu có ý nghĩa) sẽ hàm ý một đường cong chữ U xuôi chứ không
phải chữ U ngược, tức một sự đảo ngược hoàn toàn của dạng hàm cổ điển,
bản thân điều này nhất quán với FIP. Cụ thể, điều kiện C1 (độ dốc dương
ở FSTS thấp) không thỏa vì hệ số bậc nhất âm và không có ý nghĩa; C2 (độ
cong âm \(\beta_2 < 0\)) không thỏa vì dấu dương; C3 (điểm uốn nằm trong
miền dữ liệu) cho điểm uốn đại số \(\text{TP}^* = 0{,}968\) nhưng không
số hạng nào có ý nghĩa thống kê; và C4 (kiểm định Lind--Mehlum) không
bác bỏ giả thuyết quan hệ đơn điệu âm. Sự suy giảm ý nghĩa khi đưa hai
số hạng vào cùng lúc phản ánh đa cộng tuyến cao trong mẫu FSTS thấp
(FSTS trung bình \(0{,}048\)).

\hypertarget{muxf4-thux1ee9c-dux1ea5u-cux1ee7a-hux1ec7-sux1ed1-fip-qua-cuxe1c-muxf4-huxecnh-phiuxean-bux1ea3n-dux1eef-liux1ec7u-gux1ed1c-vuxe0-cux1a1-chux1ebf-nux103ng-lux1ef1c-khuxf4ng-ux111iux1ec1u-tiux1ebft}{%
\subsection{4.7.7 Mô thức dấu của hệ số FIP qua các mô hình (phiên bản
dữ liệu gốc) và cơ chế năng lực không điều
tiết}\label{muxf4-thux1ee9c-dux1ea5u-cux1ee7a-hux1ec7-sux1ed1-fip-qua-cuxe1c-muxf4-huxecnh-phiuxean-bux1ea3n-dux1eef-liux1ec7u-gux1ed1c-vuxe0-cux1a1-chux1ebf-nux103ng-lux1ef1c-khuxf4ng-ux111iux1ec1u-tiux1ebft}}

Trên phiên bản dữ liệu gốc dùng cho nghiên cứu thành phần P8, hệ số FIP
duy trì dấu âm và quy mô đáng kể qua một chuỗi mô hình, một mô thức nhất
quán nội bộ trên dữ liệu đó (lưu ý độ bền này phụ thuộc phiên bản dữ
liệu, xem Mục 4.7.8). Khi loại bỏ hiệu ứng cố định quốc gia (chỉ giữ
hiệu ứng cố định năm), hệ số âm còn lớn hơn rõ rệt:
\(\hat{\beta}(\text{FSTS}) = -3{,}351\) (\(SE = 0{,}808\);
\(t = -4{,}15\); \(p < {,}001\)), cho thấy bất lợi cấu trúc cấp quốc gia
cộng dồn vào FIP cấp doanh nghiệp. Ở mô hình song biến không kiểm soát,
quan hệ vẫn âm và ở biên ý nghĩa (\(\hat{\beta} = -0{,}864\);
\(p = {,}050\)). Hai kiểm định loại trừ bổ sung xác nhận FIP không phải
là sản phẩm phụ của mẫu hoàn chỉnh hay của định nghĩa bảy nền kinh tế:
bỏ biến kiểm soát sở hữu nước ngoài (vốn chỉ có ở \(22\%\) doanh nghiệp)
để khôi phục mẫu đầy đủ cho \(\hat{\beta}(\text{FSTS}_c) = -0{,}835\)
(\(p = {,}010\); \(N = 950\)), và mở rộng sang nhóm chín nền kinh tế đảo
nhỏ (thêm Comoros và Timor-Leste) vẫn cho
\(\hat{\beta}(\text{FSTS}_c) = -0{,}510\) (\(p = {,}008\);
\(N = 1.469\)).

Cơ chế năng lực không điều tiết được kiểm chứng ở M3. Năng lực công nghệ
có hiệu ứng trực tiếp dương và có ý nghĩa,
\(\hat{\beta}(\text{TCI}_z) = +0{,}299\) (\(SE = 0{,}099\);
\(p = {,}003\)), trong khi chấp nhận số không có ý nghĩa,
\(\hat{\beta}(\text{DAI}_z) = +0{,}094\) (\(SE = 0{,}088\);
\(p = {,}285\)). Cả hai năng lực đều không làm đổi độ dốc âm của FSTS,
nên giả thuyết không điều tiết được giữ vững. Mô thức năng lực nâng mặt
bằng năng suất nhưng không uốn được độ dốc này chính là điểm phân biệt
then chốt giữa bối cảnh SIDS và các bối cảnh châu Á lục địa: khi ba điều
kiện tiên quyết cấu trúc đồng thời vắng mặt, năng lực cấp doanh nghiệp
không thể chuyển hướng quan hệ quốc tế hóa và hiệu quả sang vùng dương.
Hệ số FSTS trong M3 vẫn âm và lớn hơn về độ lớn
(\(\hat{\beta} = -2{,}979\); \(p = {,}056\)), với việc mất ý nghĩa quy
cho cỡ mẫu giảm mạnh (còn \(N = 205\) do thiếu biến năng lực).

\hypertarget{giux1edbi-hux1ea1n-dux1eef-liux1ec7u-vuxe0-cuxe1ch-ux111ux1ecdc-thux1eadn-trux1ecdng-hux1ec7-sux1ed1-fip}{%
\subsection{4.7.8 Giới hạn dữ liệu và cách đọc thận trọng hệ số
FIP}\label{giux1edbi-hux1ea1n-dux1eef-liux1ec7u-vuxe0-cuxe1ch-ux111ux1ecdc-thux1eadn-trux1ecdng-hux1ec7-sux1ed1-fip}}

Cần đọc các hệ số FIP trên với một \textbf{giới hạn dữ liệu quan trọng}
được phát hiện trong quá trình kiểm tra tái lập. Mẫu ước lượng cốt lõi
của M1 (\(N = 209\)) gồm các doanh nghiệp có đủ biến kiểm soát ở đợt
khảo sát trước năm 2018, và trên thực tế tập trung vào \textbf{ba nền
kinh tế có dữ liệu hoàn chỉnh nhất} (Fiji 2009, Papua New Guinea 2015,
Vanuatu 2009); các mô hình độ vững mở rộng ở Mục 4.7.7 (mẫu đầy đủ
\(N = 950\), nhóm chín nền \(N = 1.469\)) được tính trên cùng phiên bản
dữ liệu gốc đó. Khi bộ dữ liệu SIDS được dựng lại để nạp đầy đủ hơn các
nền có tên tệp gạch nối và bổ sung các đợt khảo sát gần đây, hệ số cường
độ xuất khẩu trên mẫu lớn hơn trở nên \textbf{không có ý nghĩa thống kê
và không ổn định về dấu}. Như vậy độ lớn của FIP \textbf{nhạy cảm với
phiên bản dữ liệu} và với số nền đảo nhỏ có dữ liệu doanh nghiệp đầy đủ;
hơn nữa, với số cụm rất nhỏ (ba đến chín nền), suy luận dựa trên sai số
chuẩn cụm tự thân đã không đáng tin cậy. Vì những lý do đó, kết quả FIP
trong luận án được đọc đúng nhất như \textbf{bằng chứng thăm dò} cho một
\textbf{điều kiện biên lý thuyết} --- khái niệm gánh nặng quốc tế hóa
bắt buộc có giá trị định hướng --- chứ \textbf{chưa phải một hiệu ứng
được ước lượng vững} trên toàn bộ nhóm SIDS Thái Bình Dương. Việc kiểm
định FIP một cách dứt khoát đòi hỏi dữ liệu doanh nghiệp đầy đủ hơn ở
nhiều nền đảo nhỏ và một số cụm đủ lớn cho suy luận cụm hợp lệ, đây là
một hướng nghiên cứu tương lai trực tiếp (xem thêm Mục 5.5).

\begin{center}\rule{0.5\linewidth}{0.5pt}\end{center}

\hypertarget{kux1ebft-quux1ea3-bux1ed1i-cux1ea3nh-biux1ebfn-ux111ux1ed9ng-thux1ec3-chux1ebf-cux1ef1c-ux111oan-ux1ea5n-ux111ux1ed9-20142025-p9}{%
\section{4.8 Kết quả bối cảnh biến động thể chế cực đoan: Ấn Độ,
2014--2025
(P9)}\label{kux1ebft-quux1ea3-bux1ed1i-cux1ea3nh-biux1ebfn-ux111ux1ed9ng-thux1ec3-chux1ebf-cux1ef1c-ux111oan-ux1ea5n-ux111ux1ed9-20142025-p9}}

Phân tích trên ba đợt WBES Ấn Độ (2014 PICS3,
\(N_{\text{phân tích}} = 8.941\); 2022 BEE,
\(N_{\text{phân tích}} = 9.300\); 2025 BREADY,
\(N_{\text{phân tích}} = 10.476\)) cho mẫu gộp \(N = 28.717\), lớn nhất
trong các nghiên cứu đơn quốc gia của luận án. Khoảng thời gian
2014--2025 bao trùm một loạt cú sốc thể chế chưa từng có trong lịch sử
Ấn Độ hiện đại: bãi bỏ tiền mặt tháng 11/2016, ban hành Luật Phá sản và
Mất khả năng thanh toán (IBC) tháng 5/2016, áp dụng Thuế Hàng hóa và
Dịch vụ (GST) tháng 7/2017, ra mắt Giao diện Thanh toán Hợp nhất (UPI)
tháng 4/2016 đạt mức 12 tỷ giao dịch/tháng năm 2024, triển khai sáng
kiến Sản xuất tại Ấn Độ và chương trình Liên kết Khuyến khích Sản xuất
(PLI) năm 2020, cùng chính sách Tự chủ Ấn Độ (Atmanirbhar Bharat) định
hướng lại chiến lược quốc gia theo hướng nội địa hóa giai đoạn 2020.

Kết quả thực nghiệm cho thấy một dạng phá vỡ ngưỡng (threshold collapse)
không quan sát được ở bất kỳ bối cảnh quốc gia nào khác trong portfolio.
Đợt 2014 ghi nhận quan hệ U ngược cổ điển với hệ số bậc nhất dương
(\(\beta_1 = +1{,}865\), \(p < {,}001\)) và bậc hai âm
(\(\beta_2 = -1{,}508\), \(p < {,}001\)), điểm uốn ước lượng tại
\(TP = 61{,}8\%\) FSTS (khoảng tin cậy delta-method 95\%:
\([55{,}0\%; 68{,}6\%]\)), Lind--Mehlum xác nhận đường cong U ngược ở
mức \(p < {,}001\). Đợt 2022, sau ba cú sốc thể chế đầu tiên (bãi bỏ
tiền mặt, GST, IBC), đường cong U ngược vẫn duy trì nhưng điểm uốn đã
dịch chuyển đáng kể vào trong: \(\beta_1 = +1{,}542\),
\(\beta_2 = -1{,}893\), \(TP = 40{,}7\%\) (khoảng tin cậy 95\%:
\([38{,}0\%; 43{,}5\%]\)). Đợt 2025, hệ số bậc hai mất ý nghĩa thống kê
(\(\beta_2 = -0{,}160\), \(p = {,}420\)), quan hệ trở thành đơn điệu âm
trên toàn dải cường độ xuất khẩu thực tế.

Kiểm định Paternoster so sánh hệ số giữa 2014 và 2025 bác bỏ giả thuyết
cân bằng với cường độ cực mạnh: \(z(\text{FSTS}) = -7{,}94\),
\(p < {,}0001\) và \(z(\text{FSTS}^2) = +4{,}17\), \(p < {,}0001\) dưới
sai số chuẩn HC1; dưới sai số chuẩn cụm theo bang (cluster-robust,
\(G = 24\) bang), kết quả vẫn có ý nghĩa thống kê:
\(z(\text{FSTS}) = -3{,}50\), \(p = {,}001\) và
\(z(\text{FSTS}^2) = +2{,}19\), \(p = {,}029\). Mô hình gộp ba đợt với
biến giả đợt khảo sát xác nhận sự dịch chuyển có hệ thống:
\(\beta(\text{FSTS} \times \text{đợt}\_2025) = -1{,}63\)
(\(p < {,}0001\)), cho thấy độ dốc đường cong quốc tế hóa và hiệu quả đã
thay đổi căn bản giữa 2014 và 2025.

\textbf{Bảng 4.8.} \emph{Điểm uốn theo đợt khảo sát, Ấn Độ.}

\begin{longtable}[]{@{}
  >{\raggedright\arraybackslash}p{(\columnwidth - 8\tabcolsep) * \real{0.2000}}
  >{\raggedright\arraybackslash}p{(\columnwidth - 8\tabcolsep) * \real{0.2000}}
  >{\raggedright\arraybackslash}p{(\columnwidth - 8\tabcolsep) * \real{0.2000}}
  >{\raggedright\arraybackslash}p{(\columnwidth - 8\tabcolsep) * \real{0.2000}}
  >{\raggedright\arraybackslash}p{(\columnwidth - 8\tabcolsep) * \real{0.2000}}@{}}
\toprule\noalign{}
\begin{minipage}[b]{\linewidth}\raggedright
Đợt khảo sát
\end{minipage} & \begin{minipage}[b]{\linewidth}\raggedright
\(N\) phân tích
\end{minipage} & \begin{minipage}[b]{\linewidth}\raggedright
Điểm uốn (TP)
\end{minipage} & \begin{minipage}[b]{\linewidth}\raggedright
Lind--Mehlum
\end{minipage} & \begin{minipage}[b]{\linewidth}\raggedright
Nhận xét
\end{minipage} \\
\midrule\noalign{}
\endhead
\bottomrule\noalign{}
\endlastfoot
2014 (PICS3) & 8.941 & 61,8\% FSTS & \(p < {,}001\) & U ngược cổ điển \\
2022 (BEE) & 9.300 & 40,7\% FSTS & \(p < {,}001\) & U ngược, TP dịch
trong \\
2025 (BREADY) & 10.476 & Không xác định & \(p = {,}99\) & Đường cong sụp
đổ \\
\end{longtable}

\emph{Nguồn: phân tích Ấn Độ của tác giả (nghiên cứu thành phần P9), mở
rộng công trình đã công bố Do và Phan (2025, IntechOpen, phân tích 380
doanh nghiệp WBES Ấn Độ, ước lượng FGLS); bản thảo mở rộng kiểm định độ
bền ngưỡng I--P qua ba đợt 2014--2025 đang chuẩn bị cho MIR/IJOEM.}

Phát hiện ngưỡng phá vỡ ở Ấn Độ trái ngược trực tiếp với kết quả Trung
Quốc thuộc chế độ thu nhập trung bình cao (Nhóm III ICRV), trong khi Ấn
Độ thuộc Nhóm IV (chuyển đổi thu nhập trung bình--thấp), nơi điểm uốn
của Trung Quốc được chứng minh là bền vững cấu trúc qua 12 năm tăng
trưởng tích lũy (P5 China: \(TP = 49{,}4\%\) năm 2012 so với
\(TP = 47{,}2\%\) năm 2024, Paternoster \(z = +{0,}82\), \(p = {,}412\)
chấp nhận giả thuyết cân bằng). Sự đối lập này cho thấy bền vững ngưỡng
không phải đặc tính phổ quát của chế độ thể chế ICRV III mà phụ thuộc
vào mức độ và tốc độ thay đổi thể chế tích lũy trong cửa sổ quan sát.
Cần lưu ý thêm rằng điểm uốn Ấn Độ năm 2014 (\(61{,}8\%\)) cao bất
thường so với dải đặc trưng của Nhóm IV (điểm uốn vùng giữa khoảng
\(43\%\) trong mẫu gộp P7), gần với vùng của các chế độ thể chế mạnh
hơn; chính sự \textbf{hội tụ} của điểm uốn xuống \(40{,}7\%\) năm 2022
mới là dữ kiện nhất quán với chuẩn Nhóm IV, chứ không phải mức 2014.
Cách đọc này biến một mâu thuẫn bề ngoài (một nền Nhóm IV khởi đầu với
điểm uốn kiểu Nhóm I) thành bằng chứng ủng hộ: Ấn Độ tiền cú sốc vận
hành như một chế độ thể chế cao hơn vị trí thu nhập của nó, và chỉ chuỗi
cú sốc thể chế sau 2016 mới kéo điểm uốn về dải Nhóm IV trước khi làm
mất cấu trúc hẳn cấu trúc vào 2025.

Vai trò của TCI tại Ấn Độ cho thấy hiện tượng đảo dấu (sign flip) chưa
từng quan sát ở các bối cảnh quốc gia khác. Năm 2014, hiệu ứng trực tiếp
TCI dương đáng kể (\(\beta = +0{,}12\), \(p < {,}001\)), năm 2022 tiếp
tục dương và mạnh hơn (\(\beta = +0{,}19\), \(p < {,}001\)), nhưng năm
2025 đảo sang âm (\(\beta = -0{,}08\), \(p < {,}001\)) đi kèm tương tác
FSTS\(\times\)TCI dương (\(\beta = +0{,}20\), \(p = {,}03\)). Diễn giải
nhất quán nhất với khung Mizik và Jacobson (2003) về đánh đổi chú trọng
chiến lược (strategic emphasis trade-off): khi chính sách PLI và
Atmanirbhar Bharat sau 2020 chuyển hướng năng lực công nghệ sang sản
xuất nội địa, đầu tư TCI của doanh nghiệp định hướng nội địa cạnh tranh
nguồn lực với hoạt động xuất khẩu, làm triệt tiêu lợi ích năng suất TCI
vốn quan sát được khi tích lũy năng lực không thiên kiến giữa nội địa và
xuất khẩu.

Vai trò của DAI tại Ấn Độ làm rõ một sự tinh chỉnh lý thuyết quan trọng
cho khung CDCM. Tại đợt 2025, ngoài chỉ báo Tầng 1 (trang web) chung cho
cả ba đợt, WBES BREADY cung cấp thêm chỉ báo Tầng 2 (tỷ trọng doanh thu
nhận qua thanh toán điện tử, biến \texttt{k33}) làm bối cảnh thí nghiệm
gần tự nhiên cho hạ tầng số công cộng UPI. Tương tác DAI Tầng 2 với
cường độ xuất khẩu lại âm và có ý nghĩa thống kê:
\(\text{FSTS} \times \text{DAI}_{\text{Tầng 2}} = -4{,}02\),
\(p = {,}004\), ngược với dự đoán phụ phẩm cộng hưởng năng lực số công
cộng. Diễn giải cộng hưởng (complementarity reading): hạ tầng số công
cộng UPI vận hành trên đường ray rupee nội địa và yêu cầu KYC Aadhaar
cho đối tác cư trú tại Ấn Độ; doanh nghiệp xuất khẩu cần hạ tầng tài
chính xuyên biên giới khác hẳn (SWIFT, ngân hàng đại lý, tài trợ thương
mại, bảo hiểm tín dụng xuất khẩu) mà UPI không thể thay thế. Phát hiện
này cho thấy hạ tầng công cộng số chỉ thay thế năng lực tư nhân khi hai
loại hạ tầng có cùng định hướng mục đích, không phải thay thế phổ quát
cho mọi hoạt động sản xuất doanh nghiệp.

\hypertarget{muxf4-tux1ea3-mux1eabu-vuxe0-phux1b0ux1a1ng-truxecnh-hux1ed3i-quy-tux1eebng-ux111ux1ee3t}{%
\subsection{4.8.1 Mô tả mẫu và phương trình hồi quy từng
đợt}\label{muxf4-tux1ea3-mux1eabu-vuxe0-phux1b0ux1a1ng-truxecnh-hux1ed3i-quy-tux1eebng-ux111ux1ee3t}}

Mô tả mẫu cho thấy một sự thu hẹp xuất khẩu mang tính hệ thống đồng hành
với sự sụp đổ ngưỡng. Cường độ xuất khẩu trung bình dịch chuyển từ
\(7{,}2\%\) (2014) lên \(8{,}0\%\) (2022) rồi rơi xuống \(2{,}7\%\)
(2025), nhất quán với bước ngoặt hướng nội sau 2020 dưới chính sách
Atmanirbhar Bharat; mặt bằng năng suất tăng đều, với \(\ln\text{LP}\)
trung bình lần lượt là \(13{,}93\); \(14{,}37\) và \(15{,}01\). Tỷ lệ áp
dụng website (DAI Tầng 1) đạt \(52{,}7\%\) (2014), \(60{,}6\%\) (2022)
rồi giảm còn \(41{,}8\%\) (2025), trong khi tỷ trọng doanh thu qua thanh
toán điện tử (DAI Tầng 2, chỉ đo năm 2025) trung bình \(65{,}8\%\), phản
ánh mức bão hòa thanh toán hợp nhất (Unified Payments Interface ---
UPI).

Các phương trình hồi quy bậc hai từng đợt làm rõ động học sụp đổ ngưỡng.
Năm 2014, đường cong là chữ U ngược cổ điển:
\(\ln\text{LP} = 13{,}51 + 1{,}86 \cdot \text{FSTS} - 1{,}51 \cdot \text{FSTS}^2 + \mathbf{X}\boldsymbol{\beta}\),
với điểm uốn \(61{,}8\%\). Năm 2022, đường cong vẫn duy trì nhưng điểm
uốn dịch vào trong \(21\) điểm phần trăm:
\(\ln\text{LP} = 13{,}47 + 1{,}54 \cdot \text{FSTS} - 1{,}89 \cdot \text{FSTS}^2 + \mathbf{X}\boldsymbol{\beta}\),
với điểm uốn \(40{,}7\%\) và đây là đường chữ U ngược mạnh nhất trong ba
đợt (Lind--Mehlum liên hợp \(p < 1\mathrm{e}{-}12\)). Năm 2025, độ cong
mất ý nghĩa:
\(\ln\text{LP} = 14{,}24 - 0{,}36 \cdot \text{FSTS} - 0{,}16 \cdot \text{FSTS}^2 + \mathbf{X}\boldsymbol{\beta}\),
quan hệ trở nên đơn điệu âm trên toàn dải, với độ dốc tại
\(\text{FSTS}_{\min}\) là \(-0{,}36\) (\(p = {,}02\)) và tại
\(\text{FSTS}_{\max}\) là \(-0{,}68\) (\(p = {,}005\)).

\hypertarget{muxf4-huxecnh-gux1ed9p-ba-ux111ux1ee3t-vuxe0-bux1eb1ng-chux1ee9ng-dux1ecbch-chuyux1ec3n-lux169y-tiux1ebfn}{%
\subsection{4.8.2 Mô hình gộp ba đợt và bằng chứng dịch chuyển lũy
tiến}\label{muxf4-huxecnh-gux1ed9p-ba-ux111ux1ee3t-vuxe0-bux1eb1ng-chux1ee9ng-dux1ecbch-chuyux1ec3n-lux169y-tiux1ebfn}}

Mô hình gộp ba đợt với biến giả đợt khảo sát (\(N = 28.717\); \(R^2\)
hiệu chỉnh \(= 0{,}197\)) phân giải sự dịch chuyển thành các thành phần
có thể kiểm định riêng. Hệ số \(\text{Wave}\_2022 = +0{,}46\)
(\(p < {,}0001\)) và \(\text{Wave}\_2025 = +1{,}12\) (\(p < {,}0001\))
xác nhận đà nâng mặt bằng năng suất. Quan trọng hơn, tương tác
\(\text{FSTS} \times \text{đợt}\_2022 = +0{,}45\) (\(p = {,}15\)) không
cho thấy dịch chuyển bền vững so với mốc 2014, trong khi
\(\text{FSTS} \times \text{đợt}\_2025 = -1{,}63\) (\(p < {,}0001\)) xác
nhận sự sụp đổ độ dốc năm 2025; ở vế bậc hai,
\(\text{FSTS}^2 \times \text{đợt}\_2022 = -1{,}09\) (\(p = {,}001\))
phản ánh điểm uốn 2022 dịch vào trong qua việc độ cong dốc thêm, còn
\(\text{FSTS}^2 \times \text{đợt}\_2025 = +0{,}86\) (\(p = {,}007\))
phản ánh độ cong làm phẳng năm 2025. Mô hình xu hướng tuyến tính cho
\(\text{FSTS} \times \text{trend} = -0{,}10\) (\(p < {,}001\)), nhất
quán với một sự dịch chuyển cấu trúc lũy tiến chứ không phải đứt gãy đột
ngột do thay đổi công cụ khảo sát.

\hypertarget{nux103m-kiux1ec3m-ux111ux1ecbnh-ux111ux1ed9-bux1ec1n-vuxe0-sux1ef1-cuxf4-ux111ux1eb7c-cux1ee7a-chux1eef-u-ngux1b0ux1ee3c-2014-ux1edf-doanh-nghiux1ec7p-lux1edbn}{%
\subsection{4.8.3 Năm kiểm định độ bền và sự cô đặc của chữ U ngược 2014
ở doanh nghiệp
lớn}\label{nux103m-kiux1ec3m-ux111ux1ecbnh-ux111ux1ed9-bux1ec1n-vuxe0-sux1ef1-cuxf4-ux111ux1eb7c-cux1ee7a-chux1eef-u-ngux1b0ux1ee3c-2014-ux1edf-doanh-nghiux1ec7p-lux1edbn}}

Năm mô hình độ vững củng cố diễn giải cấu trúc. (R1) Mẫu chỉ doanh
nghiệp chế biến chế tạo làm mô thức mạnh thêm đôi chút: điểm uốn
\(57{,}5\%\) (2014), \(41{,}3\%\) (2022), độ cong 2025 không có ý nghĩa
(\(\beta_2 = -0{,}46\); \(p = {,}017\)); kiểm định Paternoster chéo đợt
trên mẫu này cho \(z(\text{FSTS}) = -6{,}89\) (\(p < {,}0001\)) và
\(z(\text{FSTS}^2) = +3{,}64\) (\(p = {,}0003\)), với độ lớn hiệu ứng
nằm trong khoảng \(12\%\) của ước lượng toàn mẫu, bác bỏ khả năng kết
quả do thành phần doanh nghiệp dịch vụ. (R2) Cắt bỏ doanh nghiệp có FSTS
\(> 0{,}95\) cho điểm uốn \(48{,}7\%\) (2014) và \(39{,}7\%\) (2022), độ
cong 2025 vẫn không có ý nghĩa (\(\beta_2 = -0{,}30\); \(p = {,}37\)).

Phân tích theo quy mô (R3) phát hiện một sắc thái tinh tế. Ở mẫu doanh
nghiệp nhỏ (dưới 100 lao động), chữ U ngược năm 2014 không được hỗ trợ
(\(\beta_2 = -0{,}07\); \(p = {,}84\)), chỉ định hình ở 2022
(\(\beta_2 = -1{,}64\); \(p < {,}0001\); điểm uốn \(41{,}2\%\)); ngược
lại ở mẫu doanh nghiệp lớn (từ 100 lao động), chữ U ngược hiện diện ở
2014 (điểm uốn \(50{,}2\%\)) và 2022 (điểm uốn \(40{,}8\%\)) với độ cong
rất lớn (\(\beta_2\) giữa \(-2{,}4\) và \(-3{,}0\)). Điều này cho thấy
chữ U ngược năm 2014 cô đặc ở doanh nghiệp lớn và chỉ lan tới phân khúc
nhỏ vào 2022, trùng khớp với sự mở rộng tham gia xuất khẩu của doanh
nghiệp nhỏ và vừa dưới chính sách thúc đẩy xuất khẩu giai đoạn
2014--2020. (R4) Thay biến phụ thuộc bằng năng suất lao động chuẩn hóa
theo mức cho điểm uốn \(54{,}4\%\) (2014, Lind--Mehlum ở biên
\(p = {,}048\)), \(41{,}5\%\) (2022) và độ cong sụp đổ ở 2025
(\(\beta_2 = +0{,}03\); \(p = {,}90\)).

Một lưu ý đo lường quan trọng cần đi kèm: đợt 2025 trùng với sự co hẹp
mạnh tham gia xuất khẩu, khi số doanh nghiệp có FSTS trên \(50\%\) giảm
từ \(536\) xuống còn \(141\), nên việc mất độ cong bậc hai năm 2025 một
phần phản ánh sự suy giảm biến thiên nhận dạng ở đuôi xuất khẩu cao. Tuy
vậy, chính dấu của độ dốc ở vùng xuất khẩu thấp cũng đảo chiều, từ nhánh
đi lên dương năm 2014 sang đơn điệu âm năm 2025, một thay đổi định tính
mà riêng việc đuôi mỏng đi không thể tạo ra. Đối với điều tiết DAI Tầng
2, bên cạnh tương tác bậc nhất âm đã nêu, số hạng bậc hai
\(\text{FSTS}^2 \times \text{DAI}_{\text{epay}} = +3{,}39\)
(\(p = {,}10\)); và DAI Tầng 1 (website) có hiệu ứng trực tiếp dương ở
cả ba đợt (\(\beta \approx +0{,}38\) đến \(+0{,}39\)), với tương tác
\(\text{FSTS} \times \text{DAI}\) trở thành dương có ý nghĩa ở 2022
(\(+1{,}43\); \(p = {,}02\)) rồi mất ý nghĩa ở 2025, nhất quán với vai
trò biên của website giảm dần khi UPI hấp thụ phần tương tác số hướng
tới khách hàng.

\hypertarget{bux1eb1ng-chux1ee9ng-tux1eadp-trung-vuxe0o-quux1ea3n-trux1ecb-cux1ea5p-cao-tux1ea1i-ux1ea5n-ux111ux1ed9-vai-truxf2-ux111iux1ec1u-tiux1ebft-cux1ee7a-kinh-nghiux1ec7m-vuxe0-giux1edbi-tuxednh-chux1b0ux1a1ng-suxe1ch}{%
\subsection{4.8.4 Bằng chứng tập trung vào quản trị cấp cao tại Ấn Độ:
vai trò điều tiết của kinh nghiệm và giới tính (chương
sách)}\label{bux1eb1ng-chux1ee9ng-tux1eadp-trung-vuxe0o-quux1ea3n-trux1ecb-cux1ea5p-cao-tux1ea1i-ux1ea5n-ux111ux1ed9-vai-truxf2-ux111iux1ec1u-tiux1ebft-cux1ee7a-kinh-nghiux1ec7m-vuxe0-giux1edbi-tuxednh-chux1b0ux1a1ng-suxe1ch}}

Bổ sung cho phân tích đường cong theo thời gian của P9, một nghiên cứu
thành phần khác cùng bối cảnh Ấn Độ, được công bố dưới dạng chương sách
(Do và Phan, 2025), tập trung riêng vào vai trò điều tiết của đặc điểm
nhà quản trị cấp cao. Nghiên cứu phân tích \textbf{380 doanh nghiệp}
thuộc cả khu vực chế biến chế tạo và dịch vụ (760 quan sát doanh nghiệp
và năm) từ dữ liệu WBES, sử dụng hồi quy bình phương nhỏ nhất tổng quát
khả thi (FGLS) để xử lý phương sai sai số thay đổi. Biến phụ thuộc ở đây
là tỷ suất lợi nhuận trên doanh thu (ROS, trung bình 30,34\%), khác với
thước đo năng suất lao động dùng ở các phần còn lại của luận án, và mô
hình là tuyến tính thay vì bậc hai.

Kết quả xác nhận ba mệnh đề. Thứ nhất, mức độ quốc tế hóa tác động âm
lên hiệu quả ở mô hình tuyến tính (\(\hat{\beta} = -36{,}66\),
\(p < {,}001\)), nhất quán với logic chi phí gia nhập giai đoạn đầu của
mô hình Uppsala, tức nhánh đi xuống bên trái của quan hệ phi tuyến khi
quan sát ở dải cường độ quốc tế hóa thấp đặc trưng của mẫu này. Thứ hai,
kinh nghiệm của nhà quản trị điều tiết dương quan hệ này
(\(\hat{\beta} = +1{,}55\), \(p < {,}001\)), làm dịu tác động bất lợi
của quốc tế hóa lên lợi nhuận. Thứ ba, lãnh đạo nữ cũng điều tiết dương
quan hệ (\(\hat{\beta} = +77{,}03\), \(p < {,}001\)), củng cố giả thuyết
H4 ở khía cạnh điều tiết.

Phát hiện về giới tính ở đây cần được đối chiếu thận trọng với hệ số âm
của biến quản trị nữ trong các ước lượng dựa trên năng suất (P7 toàn mẫu
và P10 Nhật Bản). Sự khác dấu này không phải mâu thuẫn nội tại mà phản
ánh ba khác biệt thiết kế: biến phụ thuộc là lợi nhuận trên doanh thu
chứ không phải năng suất lao động; tham số ước lượng là hệ số điều tiết
độ dốc của quan hệ quốc tế hóa và hiệu quả chứ không phải hiệu ứng mức
chính; và mẫu là 380 doanh nghiệp Ấn Độ chứ không phải mẫu gộp đa quốc
gia. Cách đọc hài hòa, nhất quán với phần tổng hợp giả thuyết của luận
án, là đặc điểm nhà quản trị cấp cao có thể nâng chất lượng quyết định
quốc tế hóa (biểu hiện thành điều tiết dương trên thước đo lợi nhuận
trong một bối cảnh đơn quốc gia), trong khi hệ số mức âm trên thước đo
năng suất ở mẫu gộp phản ánh chọn lọc cấu thành chứ không phải hiệu quả
lãnh đạo. Đặt cạnh nhau, hai mảnh bằng chứng cho thấy vai trò của đặc
trưng quản trị là phụ thuộc thước đo và phụ thuộc bối cảnh, đúng như kết
luận chung của luận án rằng H4 không được hỗ trợ như một hiệu ứng mức
phổ quát mà cần phân tích nhận biết chọn lọc trong nội bộ từng bối cảnh.

\hypertarget{cux1ea5u-truxfac-mux1eabu-ba-ux111ux1ee3t-vuxe0-khoux1ea3ng-tin-cux1eady-ux111iux1ec3m-uux1ed1n-theo-phux1b0ux1a1ng-phuxe1p-delta}{%
\subsection{4.8.5 Cấu trúc mẫu ba đợt và khoảng tin cậy điểm uốn theo
phương pháp
delta}\label{cux1ea5u-truxfac-mux1eabu-ba-ux111ux1ee3t-vuxe0-khoux1ea3ng-tin-cux1eady-ux111iux1ec3m-uux1ed1n-theo-phux1b0ux1a1ng-phuxe1p-delta}}

Bộ ba khoảng tin cậy điểm uốn theo phương pháp delta cung cấp bằng chứng
định lượng trực tiếp cho luận điểm dịch chuyển ngưỡng thay vì chỉ dựa
vào giá trị điểm. Mẫu phân tích thu được sau khi loại bỏ theo danh sách
trên bộ biến trọng tâm và kiểm tra tính hợp lệ của tuổi doanh nghiệp là
\(8.941\) (2014 PICS3, từ \(9.281\) thô), \(9.300\) (2022 BEE, từ
\(9.376\) thô) và \(10.476\) (2025 BREADY, từ \(10.479\) thô), cho mẫu
gộp phân tích \(28.717\) trên tổng thô \(29.136\). Điểm uốn 2014 ở mức
\(61{,}8\%\) với khoảng tin cậy delta-method \([55{,}0\%; 68{,}6\%]\),
một mức đỉnh năng suất rất cao đặc trưng cho bối cảnh thị trường mới nổi
nơi doanh nghiệp xuất khẩu chiết xuất lợi suất học hỏi qua xuất khẩu
đáng kể trước khi chi phí mở rộng quá mức bám chặt. Điểm uốn 2022 dịch
vào trong còn \(40{,}7\%\) với khoảng tin cậy \([38{,}0\%; 43{,}5\%]\),
và đây là chữ U ngược mạnh nhất ba đợt với kiểm định Lind--Mehlum liên
hợp \(p < 1\mathrm{e}{-}12\). Điểm cốt lõi là hai khoảng tin cậy 2014 và
2022 \textbf{không chồng lấn} (\([55{,}0; 68{,}6]\) so với
\([38{,}0; 43{,}5]\)), bằng chứng định lượng trực tiếp rằng dịch chuyển
ngưỡng \(21\) điểm phần trăm là thực chứ không phải nhiễu ước lượng,
tương phản sắc nét với ba khoảng tin cậy chồng lấn chặt của Trung Quốc
đã trình bày ở Mục 4.4.4 (nơi điểm uốn ổn định cấu trúc).

\hypertarget{phux1ea3n-buxe1c-giux1ea3-thuyux1ebft-tux1ea1o-tuxe1c-cuxf4ng-cux1ee5-khux1ea3o-suxe1t-qua-ux111ux1ee3t-cux1ea7u-nux1ed1i-2022}{%
\subsection{4.8.6 Phản bác giả thuyết tạo tác công cụ khảo sát qua đợt
cầu nối
2022}\label{phux1ea3n-buxe1c-giux1ea3-thuyux1ebft-tux1ea1o-tuxe1c-cuxf4ng-cux1ee5-khux1ea3o-suxe1t-qua-ux111ux1ee3t-cux1ea7u-nux1ed1i-2022}}

Một quan ngại tự nhiên là liệu sự dịch chuyển biểu kiến có phải tạo tác
của việc thay đổi công cụ khảo sát hay không, vì đợt 2014 dùng công cụ
PICS3 còn đợt 2025 dùng công cụ BREADY mới, hai bộ không đồng nhất về
diễn đạt câu hỏi và thiết kế chọn mẫu. Đợt cầu nối 2022 (BEE trung gian)
cung cấp kiểm định quyết định: nếu sự sụp đổ 2025 thuần túy là tạo tác
công cụ, một bước gián đoạn sẽ phải xuất hiện giữa 2022 và 2025 (chuyển
tiếp BEE sang BREADY), trong khi 2014 và 2022 không phân biệt được. Dữ
liệu cho thấy mô thức ngược lại: điểm uốn dịch đơn điệu vào trong từ
\(61{,}8\%\) (2014) xuống \(40{,}7\%\) (2022) qua cửa sổ bãi bỏ tiền
mặt, GST và IBC trên cùng công cụ phái sinh từ PICS, rồi quan hệ mới mất
cấu trúc vào 2025; mô hình xu hướng tuyến tính cho
\(\mathrm{FSTS} \times \mathrm{trend} = -0{,}10\) (\(p < {,}001\)) khẳng
định một tiến trình đơn điệu thay vì đứt gãy đột ngột. Mô hình gộp ba
đợt (\(N = 28.717\); \(R^2\) hiệu chỉnh \(= 0{,}197\)) định lượng các
thành phần dịch chuyển:
\(\mathrm{FSTS}^2 \times \text{đợt}\_2022 = -1{,}09\) (\(p = {,}001\))
xác nhận điểm uốn 2022 dịch vào trong qua độ cong dốc thêm, còn
\(\mathrm{FSTS}^2 \times \text{đợt}\_2025 = +0{,}86\) (\(p = {,}007\))
xác nhận độ cong làm phẳng năm 2025. Tiến trình dịch chuyển cấu trúc
khớp dữ liệu; tạo tác công cụ thì không.

\hypertarget{phuxe2n-ruxe3-theo-quy-muxf4-vuxe0-nux103m-kiux1ec3m-ux111ux1ecbnh-ux111ux1ed9-vux1eefng-ux111ux1ecbnh-vux1ecb-nguux1ed3n-dux1ecbch-chuyux1ec3n}{%
\subsection{4.8.7 Phân rã theo quy mô và năm kiểm định độ vững định vị
nguồn dịch
chuyển}\label{phuxe2n-ruxe3-theo-quy-muxf4-vuxe0-nux103m-kiux1ec3m-ux111ux1ecbnh-ux111ux1ed9-vux1eefng-ux111ux1ecbnh-vux1ecb-nguux1ed3n-dux1ecbch-chuyux1ec3n}}

Năm kiểm định độ vững định vị chính xác phân khúc doanh nghiệp nơi chữ U
ngược hình thành rồi mất cấu trúc. Mẫu chỉ chế biến chế tạo (ISIC 15--37
cho 2014; ISIC 10--33 cho 2022 và 2025; \(N = 6.974\), \(5.366\),
\(5.706\)) làm mô thức mạnh thêm với điểm uốn \(57{,}5\%\) (2014) và
\(41{,}3\%\) (2022), kiểm định Paternoster chéo đợt cho
\(z(\mathrm{FSTS}) = -6{,}89\) (\(p < {,}0001\)) và
\(z(\mathrm{FSTS}^2) = +3{,}64\) (\(p = {,}0003\)), với độ lớn hiệu ứng
nằm trong khoảng \(12\%\) của ước lượng toàn mẫu, bác bỏ khả năng kết
quả do thành phần doanh nghiệp dịch vụ. Phân rã theo quy mô phát hiện
sắc thái then chốt: ở mẫu doanh nghiệp nhỏ (dưới 100 lao động;
\(N = 6.844\), \(6.497\), \(7.383\)), chữ U ngược 2014 không được hỗ trợ
(\(\beta_2 = -0{,}07\), \(p = {,}84\)), chỉ định hình ở 2022
(\(\beta_2 = -1{,}64\), \(p < {,}0001\); điểm uốn \(41{,}2\%\)); ngược
lại ở mẫu doanh nghiệp lớn (từ 100 lao động; \(N = 2.097\), \(2.803\),
\(3.093\)), chữ U ngược hiện diện ở 2014 (điểm uốn \(50{,}2\%\)) và 2022
(điểm uốn \(40{,}8\%\)) với độ cong rất lớn (\(\beta_2\) giữa \(-2{,}4\)
và \(-3{,}0\)). Điều này cho thấy chữ U ngược 2014 cô đặc ở doanh nghiệp
lớn và chỉ lan tới phân khúc nhỏ vào 2022, trùng khớp với sự mở rộng
tham gia xuất khẩu của doanh nghiệp nhỏ và vừa dưới chính sách thúc đẩy
xuất khẩu giai đoạn 2014--2020 rồi đảo chiều sau 2020 dưới Atmanirbhar
Bharat. Bốn trong năm mô hình (R1, R2, R3b, R4) xác nhận mô thức định
tính chữ U ngược hiện diện 2014--2022 rồi biến mất 2025; mô hình thứ năm
(R3a) xác nhận hiện diện 2022 và biến mất 2025.

\hypertarget{ux111ux1ed9-vux1eefng-suy-luux1eadn-trux1b0ux1edbc-sai-sux1ed1-chuux1ea9n-cux1ee5m-vuxe0-quan-ngux1ea1i-nux1ed9i-sinh-upi}{%
\subsection{4.8.8 Độ vững suy luận trước sai số chuẩn cụm và quan ngại
nội sinh
UPI}\label{ux111ux1ed9-vux1eefng-suy-luux1eadn-trux1b0ux1edbc-sai-sux1ed1-chuux1ea9n-cux1ee5m-vuxe0-quan-ngux1ea1i-nux1ed9i-sinh-upi}}

Tính bền vững của phát hiện dịch chuyển ngưỡng được kiểm tra trước hai
cấp độ suy luận. Dưới sai số chuẩn vững HC1, kiểm định Paternoster bác
bỏ cân bằng hệ số 2014--2025 với cường độ cực mạnh
(\(z(\mathrm{FSTS}) = -7{,}94\); \(z(\mathrm{FSTS}^2) = +4{,}17\); cả
hai \(p < {,}0001\)); dưới sai số chuẩn cụm theo bang bảo thủ hơn
(khoảng 24 cụm, theo Cameron và cộng sự), kết quả vẫn có ý nghĩa
(\(z(\mathrm{FSTS}) = -3{,}50\), \(p = {,}0005\);
\(z(\mathrm{FSTS}^2) = +2{,}19\), \(p = {,}029\)) bất chấp cỡ mẫu hiệu
dụng nhỏ hơn nhiều mà 24 cụm bang hàm ý. Với số cụm ở mức biên này, biện
minh tiệm cận cho suy luận cụm-vững là cận biên, nên đây được đọc là
kiểm định bảo thủ chứ không phải mô hình chính. Đối với tương tác UPI
Tầng 2 âm
(\(\mathrm{FSTS} \times \mathrm{DAI}_{\text{epay}} = -4{,}02\),
\(p = {,}004\)), một quan ngại nội sinh được thừa nhận rõ: ở các doanh
nghiệp 2025 có tỉ trọng thanh toán UPI cao, phân phối cường độ xuất khẩu
bị cắt cụt về phía thấp một cách cơ học, vì doanh nghiệp có khách hàng
nước ngoài không thanh toán qua UPI nên tỉ trọng UPI cao và tỉ trọng
xuất khẩu cao nghịch biến ngay trong quá trình sinh dữ liệu. Do đó việc
bác bỏ H4b (Tầng 2 lấn át Tầng 1) được đọc là gợi ý chứ không dứt khoát,
với diễn giải cộng hưởng là cách giải thích nhất quán lý thuyết nhất: hạ
tầng số công cộng chỉ thay thế năng lực tư nhân khi cùng định hướng mục
đích, và UPI định hướng nội địa không thể thay thế hạ tầng tài chính
xuyên biên giới (SWIFT, ngân hàng đại lý, tài trợ thương mại) mà doanh
nghiệp xuất khẩu cần.

\begin{center}\rule{0.5\linewidth}{0.5pt}\end{center}

\hypertarget{kux1ebft-quux1ea3-bux1ed1i-cux1ea3nh-tiuxean-tiux1ebfn-trux1b0ux1edfng-thuxe0nh-nhux1eadt-bux1ea3n-p10}{%
\section{4.9 Kết quả bối cảnh tiên tiến trưởng thành: Nhật Bản
(P10)}\label{kux1ebft-quux1ea3-bux1ed1i-cux1ea3nh-tiuxean-tiux1ebfn-trux1b0ux1edfng-thuxe0nh-nhux1eadt-bux1ea3n-p10}}

Nghiên cứu thành phần P10 khai thác đợt khảo sát WBES khai mở của Nhật
Bản năm 2025 (\(N = 2.168\) cơ sở; 339 biến hài hòa; chọn mẫu ngẫu nhiên
phân tầng theo ngành, quy mô và vùng) để kiểm định trực tiếp \textbf{dự
đoán biên trên} của khung phụ thuộc thể chế: tại chế độ thể chế mạnh
nhất, nhánh đi xuống của chữ U ngược bị đẩy ra ngoài miền cường độ xuất
khẩu quan sát được, khiến quan hệ quốc tế hóa và hiệu quả trong dữ liệu
trở nên dương và xấp xỉ tuyến tính. Bối cảnh mang sức nặng lý thuyết đặc
biệt: chính các tập đoàn đa quốc gia Nhật Bản là nơi chữ U ngược được
thiết lập ban đầu, song Nhật Bản chưa từng được quan sát trên một bộ
công cụ hài hòa cho phép định vị nó so với các nền kinh tế cùng khu vực.
Năng suất lao động (ln doanh thu trên một lao động thường xuyên,
winsorize tại phân vị 1/99) được hồi quy theo FSTS và bình phương của nó
với hiệu ứng cố định ngành và sai số chuẩn vững HC1, theo chuỗi phân
tầng M1--M5.

\hypertarget{quan-hux1ec7-cux1b0ux1eddng-ux111ux1ed9-duux1ed7i-thux1eb3ng-vux1ec1-gux1ea7n-tuyux1ebfn-tuxednh}{%
\subsection{4.9.1 Quan hệ cường độ duỗi thẳng về gần tuyến
tính}\label{quan-hux1ec7-cux1b0ux1eddng-ux111ux1ed9-duux1ed7i-thux1eb3ng-vux1ec1-gux1ea7n-tuyux1ebfn-tuxednh}}

Hiệu ứng tuyến tính của cường độ xuất khẩu lên năng suất \textbf{dương
và có ý nghĩa cao} ở mô hình nền (\(M1: \beta_1 = +0{,}671\);
\(p < 0{,}001\)). Khi bổ sung số hạng bậc hai (M2), số hạng tuyến tính
vẫn dương (\(+1{,}042\); \(p = 0{,}002\)) trong khi bậc hai âm nhưng
\textbf{xa ý nghĩa} (\(-0{,}606\); \(p = 0{,}323\)); điểm uốn đại số hàm
ý 90,1\% nằm sát rìa trên cùng của miền FSTS và không phân biệt được với
trạng thái ``không có điểm uốn''. Khi lần lượt thêm kiểm soát nhân khẩu
(M3), năng lực và chấp nhận số (M4), rồi tương tác uốn đường cong (M5),
\textbf{số hạng bậc hai không bao giờ tiệm cận ý nghĩa} và điểm uốn hàm
ý trôi tới các giá trị phi lý (108\% ở M3, vượt 700\% khi kiểm soát năng
lực), dấu hiệu điển hình của một hệ số không phân biệt được với 0. Theo
tiêu chí Lind--Mehlum (2010), \textbf{không một mô hình nào xác nhận chữ
U ngược}: trong miền quan sát (FSTS trung bình 4,1\%, trung bình có điều
kiện 24,6\%), quốc tế hóa nâng hiệu quả một cách đơn điệu.

\textbf{Bảng 4.9.} \emph{Hiệu ứng cường độ xuất khẩu (FSTS) lên năng
suất qua chuỗi mô hình phân tầng, Nhật Bản 2025 (\(N = 2.168\); HC1).}

\begin{longtable}[]{@{}
  >{\raggedright\arraybackslash}p{(\columnwidth - 10\tabcolsep) * \real{0.1667}}
  >{\raggedleft\arraybackslash}p{(\columnwidth - 10\tabcolsep) * \real{0.1667}}
  >{\raggedleft\arraybackslash}p{(\columnwidth - 10\tabcolsep) * \real{0.1667}}
  >{\raggedleft\arraybackslash}p{(\columnwidth - 10\tabcolsep) * \real{0.1667}}
  >{\raggedleft\arraybackslash}p{(\columnwidth - 10\tabcolsep) * \real{0.1667}}
  >{\raggedleft\arraybackslash}p{(\columnwidth - 10\tabcolsep) * \real{0.1667}}@{}}
\toprule\noalign{}
\begin{minipage}[b]{\linewidth}\raggedright
Hệ số
\end{minipage} & \begin{minipage}[b]{\linewidth}\raggedleft
M1
\end{minipage} & \begin{minipage}[b]{\linewidth}\raggedleft
M2
\end{minipage} & \begin{minipage}[b]{\linewidth}\raggedleft
M3
\end{minipage} & \begin{minipage}[b]{\linewidth}\raggedleft
M4
\end{minipage} & \begin{minipage}[b]{\linewidth}\raggedleft
M5
\end{minipage} \\
\midrule\noalign{}
\endhead
\bottomrule\noalign{}
\endlastfoot
FSTS (định tâm) & +0,671*** & +1,042** & +0,872* & +0,476 & +0,486 \\
FSTS² & --- & −0,606 (ns) & (ns) & (ns) & (ns) \\
Điểm uốn đại số & --- & {[}90,1\%{]} & {[}108,3\%{]} & ngoài miền &
ngoài miền \\
Bổ sung kiểm soát & nền & +bậc hai & +nhân khẩu & +TCI/DAI & +tương
tác \\
\end{longtable}

\emph{Ghi chú: sai số chuẩn vững HC1. * \(p<0{,}05\); ** \(p<0{,}01\);
*** \(p<0{,}001\); (ns) không có ý nghĩa thống kê. Không mô hình nào có
FSTS² âm có ý nghĩa nên không xác nhận chữ U ngược (Lind \& Mehlum,
2010). Độ vững với d3c riêng: \(\beta_1 = +1{,}255\)**;
\(\beta_2 = -0{,}674\) (\(p = 0{,}445\)), cùng kết luận. Nguồn: P10,
\texttt{p10\_japan/replication/}.}

Kết luận dương đơn điệu \textbf{vững qua sáu mô hình độ vững}: thêm kiểm
soát quy mô (\(+0{,}656\); \(p = 0{,}046\)), giới hạn cơ sở nội địa
(\(+0{,}917\); \(p = 0{,}012\) --- loại trừ khả năng số ít doanh nghiệp
vốn nước ngoài dẫn dắt), và winsorize năng suất tại phân vị 5/95
(\(+0{,}837\); \(p = 0{,}006\)); ở mọi mô hình số hạng bậc hai vô nghĩa.
Riêng mô hình chỉ trên doanh nghiệp xuất khẩu (\(N = 344\)), hiệu ứng
cường độ dương nhưng không còn ý nghĩa (\(+0{,}200\); \(p = 0{,}410\)),
nhất quán với việc trọng tâm nằm ở biên rộng chứ không phải cường độ.
\textbf{H1 được ủng hộ} --- Nhật Bản nằm trên một nhánh đi lên dài mà
nhánh đi xuống vắng mặt về mặt thực nghiệm, trùng khớp kết quả Singapore
(Mục 4.3) và kết quả cấp nhóm Advanced của ước lượng 50 nền (Mục 4.6.2),
hoàn tất bức tranh ba vùng của phổ thể chế: chữ U ngược ở các chế độ
chuyển đổi, gần tuyến tính ở biên trên, mất cấu trúc ở biên dưới.

\hypertarget{nux103ng-lux1ef1c-sux1ed1-huxf3a-vuxe0-ux111ux1eb7c-ux111iux1ec3m-quux1ea3n-trux1ecb-nuxe2ng-mux1eb7t-bux1eb1ng-khuxf4ng-uux1ed1n-ux111ux1b0ux1eddng-cong}{%
\subsection{4.9.2 Năng lực, số hóa và đặc điểm quản trị: nâng mặt bằng,
không uốn đường
cong}\label{nux103ng-lux1ef1c-sux1ed1-huxf3a-vuxe0-ux111ux1eb7c-ux111iux1ec3m-quux1ea3n-trux1ecb-nuxe2ng-mux1eb7t-bux1eb1ng-khuxf4ng-uux1ed1n-ux111ux1b0ux1eddng-cong}}

Năng lực công nghệ có \textbf{hiệu ứng mặt bằng chính xác}
(\(+0{,}100\); \(p < 0{,}001\)) và hấp thụ phần lớn hiệu ứng FSTS thô
khi được đưa vào M4, nhất quán với chọn lọc theo năng lực vào hoạt động
xuất khẩu; tương tác uốn đường cong không có ý nghĩa (M5: \(-0{,}126\);
\(p = 0{,}328\)) --- năng lực nâng sàn, không uốn đường cong (\textbf{H2
được ủng hộ}). Chấp nhận số giữ phần bù dương khiêm tốn (\(+0{,}110\);
\(p = 0{,}028\)) dù tỷ lệ website đã bão hòa ở 83,8\%, không có tương
tác uốn đường cong; phần bù bằng khoảng một nửa mức của ước lượng gộp 50
nền, đúng mức nén mà trạng thái \textbf{bão hòa Tầng 1} dự đoán
(\textbf{H3 được ủng hộ}). Hai quy luật đặc thù Nhật Bản kỷ luật hóa
cách đọc hiệu ứng thượng tầng quản trị: phần bù \textbf{trường thọ doanh
nghiệp} dương (\(+0{,}073\); \(p = 0{,}007\)), mỗi đơn vị log tuổi thêm
khoảng 7\% năng suất, cộng hưởng với truyền thống \emph{shinise} ở nền
kinh tế có cơ sở lâu đời nhất mẫu gộp (tuổi trung bình 50,4 năm)
(\textbf{H5 được ủng hộ}); và hệ số \textbf{âm} của nữ quản lý cấp cao
(\(-0{,}170\); \(p = 0{,}010\)) trong nền kinh tế có biên lãnh đạo nữ
mỏng nhất khu vực (7,3\%), được đọc thận trọng như chọn lọc số ít cơ sở
do nữ điều hành vào các ngách biên thấp chứ \textbf{không diễn giải nhân
quả}.

\hypertarget{dux1ecbch-chuyux1ec3n-trux1ecdng-tuxe2m-sang-biuxean-rux1ed9ng-ai-xuux1ea5t-khux1ea9u}{%
\subsection{4.9.3 Dịch chuyển trọng tâm sang biên rộng: ai xuất
khẩu}\label{dux1ecbch-chuyux1ec3n-trux1ecdng-tuxe2m-sang-biuxean-rux1ed9ng-ai-xuux1ea5t-khux1ea9u}}

Vì đường cường độ phẳng trong miền quan sát, biến thiên có ý nghĩa nằm ở
\textbf{ai xuất khẩu}. Mô hình xác suất tuyến tính của tham gia xuất
khẩu cho thấy tham gia được chi phối bởi \textbf{sở hữu nước ngoài}
(\(+0{,}408\); \(p < 0{,}001\) --- một cơ sở có ít nhất 10\% vốn nước
ngoài có xác suất xuất khẩu cao hơn khoảng bốn mươi điểm phần trăm),
\textbf{năng lực công nghệ} (\(+0{,}083\) trên mỗi độ lệch chuẩn;
\(p < 0{,}001\)) và \textbf{chấp nhận số} (\(+0{,}046\);
\(p = 0{,}005\)); mô hình giải thích \(R^2 = 0{,}228\) biến thiên cắt
ngang của tham gia. Doanh nghiệp xuất khẩu hưởng \textbf{phần bù năng
suất 25,8\%} so với doanh nghiệp không xuất khẩu (\(p < 0{,}001\), kiểm
soát sở hữu, tuổi, giới tính và ngành). Đây chính là cơ chế \textbf{tự
chọn lọc Melitz (2003)} ở một nền kinh tế biên giới: tại biên trên của
phổ thể chế, câu hỏi ràng buộc không phải doanh nghiệp xuất khẩu
\emph{bao nhiêu} mà là \emph{có} xuất khẩu hay không --- lợi ích năng
suất của quốc tế hóa vận hành chủ yếu qua ngưỡng tham gia, và các doanh
nghiệp vượt ngưỡng là những doanh nghiệp được trang bị năng lực, công
nghệ, số và gắn với sở hữu nước ngoài (\textbf{H4 được ủng hộ}). Bằng
chứng Nhật Bản do đó vừa hoàn tất phổ ba vùng của khung CDCM ở biên
trên, vừa làm rõ rằng cơ chế quốc tế hóa và hiệu quả chuyển từ biên
cường độ (ở các chế độ chuyển đổi) sang biên rộng (ở chế độ tiên tiến
trưởng thành).

\hypertarget{tam-giuxe1c-huxf3a-kux1ebft-quux1ea3-cux1ea5p-doanh-nghiux1ec7p-vux1edbi-bux1eb1ng-chux1ee9ng-vux129-muxf4-ux111ux1ed9c-lux1eadp}{%
\section{4.10 Tam giác hóa kết quả cấp doanh nghiệp với bằng chứng vĩ mô
độc
lập}\label{tam-giuxe1c-huxf3a-kux1ebft-quux1ea3-cux1ea5p-doanh-nghiux1ec7p-vux1edbi-bux1eb1ng-chux1ee9ng-vux129-muxf4-ux111ux1ed9c-lux1eadp}}

Một thử thách then chốt đối với tính giá trị của các phát hiện cấp doanh
nghiệp là liệu chúng có nhất quán với bằng chứng vĩ mô độc lập hay
không. Mục này thực hiện tam giác hóa (triangulation) bốn kết quả trụ
cột của luận án với các báo cáo chuyên đề mới nhất của Ngân hàng Thế
giới (2026), qua đó củng cố tính giá trị hội tụ (convergent validity)
của khung ICRV.

\textbf{Thứ nhất, phổ thể chế của điểm uốn (H5) hội tụ với quy luật vĩ
mô ``thể chế tốt hơn đi liền tăng trưởng nhanh hơn''.} Luận án phát hiện
chế độ thể chế quyết định chính sự tồn tại và vị trí của chữ U ngược:
đường cong định hình rõ nét nhất ở vùng giữa phổ thể chế (Nhóm IV, điểm
uốn 43,0\%), duỗi thẳng gần tuyến tính ở biên trên và mất cấu trúc ở
biên dưới (Mục 4.6.2). Ở cấp vĩ mô, Ngân hàng Thế giới ghi nhận rằng các
thị trường biên tăng trưởng nhanh nhất là những nền đã cải thiện quản
trị, thu hẹp chênh lệch lãi suất ngân hàng và kiểm soát nợ công, trong
khi các điểm yếu thể chế đi kèm tăng trưởng vốn đáng thất vọng (World
Bank, 2026c). Hai tầng bằng chứng cùng chỉ về một cơ chế: chất lượng thể
chế quyết định mức độ và điều kiện mà tích lũy vốn và quốc tế hóa chuyển
hóa thành kết quả thực. Phổ điểm uốn cấp doanh nghiệp của luận án vì vậy
có thể đọc như biểu hiện vi mô của quy luật phân hóa tăng trưởng vĩ mô.

\textbf{Thứ hai, hai cơ chế FDI tương phản giữa nhóm tài nguyên và nhóm
đổi mới ăn khớp với cú sốc hàng hóa 2026.} Luận án cho thấy nhóm
Advanced tài nguyên dẫn dắt (Nhóm II, các nền GCC) có phân tán năng suất
hẹp đặc trưng của kinh tế nhà nước tô, nơi phân phối lại tô lợi tài
nguyên san bằng khác biệt năng suất giữa doanh nghiệp (Mục 4.1, Chuyên
đề 1). Cú sốc xung đột Trung Đông năm 2026, với giá năng lượng dự báo
tăng 24\% và giá hàng hóa nói chung tăng 16\%, đẩy giá kim loại cơ bản
và quý lên mức cao kỷ lục (World Bank, 2026d), củng cố vị thế tô lợi của
nhóm này trong ngắn hạn nhưng đồng thời khắc sâu rủi ro phụ thuộc tài
nguyên mà mô hình tăng trưởng dựa trên đổi mới (Nhóm I) không gặp phải.
Bối cảnh hàng hóa 2026 vì vậy làm rõ hơn vì sao việc gộp hai loại
Advanced thành một nhóm duy nhất, như tài liệu trước thường làm, che
khuất hai cơ chế chuyển hóa quốc tế hóa hoàn toàn khác nhau.

\textbf{Thứ ba, Nghịch lý bão hòa số và vai trò DAI ăn khớp với chương
trình nghị sự về thể chế công trong kỷ nguyên AI.} Luận án cho thấy phần
thưởng của hiện diện số Tầng 1 co lại ở các nền kinh tế tiên tiến nơi nó
đã tiệm cận bão hòa, trong khi vẫn còn dư địa ở các nền đang chuyển đổi
(Mục 4.3.2, Mục 4.6.3). Phân tích của Ngân hàng Thế giới về thể chế công
trong kỷ nguyên trí tuệ nhân tạo nhấn mạnh rằng hạ tầng số công và quản
trị dữ liệu là điều kiện nền tảng quyết định năng lực số được chuyển hóa
thành giá trị, và rằng giai đoạn trưởng thành số khác nhau đòi hỏi công
cụ chính sách khác nhau (World Bank, 2026e). Sự tương đồng này củng cố
luận điểm cốt lõi của khung CDCM: phần thưởng của áp dụng số là hàm của
mức độ phát triển thể chế và tầng năng lực số được đo, chứ không phải
một ``phần bù số'' đồng nhất.

\textbf{Thứ tư, FIP và thách thức việc làm của các nền kinh tế biên định
vị đóng góp chính sách của luận án.} Phát hiện FIP tại SIDS (Mục 4.7) và
việc xác nhận rằng quốc tế hóa chỉ tạo phần thưởng năng suất trong điều
kiện thể chế và năng lực đủ mạnh, đặt trong bối cảnh hơn 230 triệu người
sẽ bước vào tuổi lao động ở các thị trường biên giai đoạn 2025--2035
(World Bank, 2026c), cho thấy chính sách thúc đẩy xuất khẩu đơn thuần là
chưa đủ và có thể phản tác dụng ở những bối cảnh dễ tổn thương. Hàm ý
chính sách của luận án, ưu tiên nâng cấp năng lực công nghệ trước khi mở
rộng cường độ quốc tế hóa, vì vậy nhất quán với khuyến nghị vĩ mô của
Ngân hàng Thế giới về việc đầu tư vào hạ tầng nền tảng và năng lực hấp
thụ như điều kiện để các nền kinh tế biên hiện thực hóa tiềm năng (World
Bank, 2026c).

\textbf{Thứ năm, cú sốc thương mại Việt Nam đầu năm 2026 minh họa trực
tiếp tính hệ quả của nhánh đi xuống đối với nhóm chuyển đổi.} Dữ liệu
giám sát vĩ mô của Ngân hàng Thế giới (2026) cho Việt Nam, nền kinh tế
Nhóm IV lớn nhất trong mẫu, ghi nhận chỉ số nhà quản trị mua hàng ngành
chế tạo rơi xuống 45,3 điểm (vùng thu hẹp) vào tháng 3 năm 2026 khi đơn
hàng xuất khẩu mới sụt giảm sau khi một cuộc điều tra của cơ quan đại
diện thương mại Hoa Kỳ được khởi động; nhập khẩu hàng hóa đã tăng 27,8\%
so với cùng kỳ năm trước tính đến cùng tháng, một phần do đứt gãy chuỗi
cung ứng năng lượng từ xung đột tại eo biển Hormuz. Đây chính là điều
kiện mà vế phải của đường cong chữ U ngược trở nên hệ quả: các doanh
nghiệp bị đẩy vượt ngưỡng cường độ xuất khẩu tối ưu bởi cam kết xuất
khẩu trước đó phải gánh áp lực chi phí điều phối và chi phí đầu vào
khuếch đại. Vì Việt Nam nằm trong chế độ chuyển đổi nơi chữ U ngược sắc
nét nhất (điểm uốn Nhóm IV khoảng 43\%), các doanh nghiệp xuất khẩu Việt
Nam hiện ở trên ngưỡng này thuộc nhóm chịu phơi nhiễm cao nhất trước các
cú sốc vĩ mô. Bằng chứng năng lực hàm ý rằng xây dựng năng lực công nghệ
và số nâng sàn năng suất cho các doanh nghiệp này, dù trên ước lượng gộp
năng lực không tự nó làm thoải nhánh đi xuống; ràng buộc quyết định là
độ rộng của tham gia xuất khẩu có năng lực, không phải một biện pháp
dịch chuyển đường cong.

Năm điểm hội tụ trên cho thấy các phát hiện cấp doanh nghiệp của luận án
không đứng tách biệt mà ăn khớp một cách hệ thống với bằng chứng vĩ mô
độc lập gần nhất. Sự tam giác hóa này tăng cường tính giá trị của khung
ICRV và đặt đóng góp của luận án vào đúng dòng chảy nghiên cứu phát
triển đương đại.

\begin{center}\rule{0.5\linewidth}{0.5pt}\end{center}

\hypertarget{tux1ed5ng-hux1ee3p-kux1ebft-quux1ea3-theo-hux1ec7-giux1ea3-thuyux1ebft}{%
\section{4.11 Tổng hợp kết quả theo hệ giả
thuyết}\label{tux1ed5ng-hux1ee3p-kux1ebft-quux1ea3-theo-hux1ec7-giux1ea3-thuyux1ebft}}

Trước khi tổng hợp theo hệ giả thuyết, Bảng 4.10 hợp nhất các ước lượng
trụ cột của chín nghiên cứu thành phần thực nghiệm, cho phép đối chiếu
trực tiếp dạng hàm, điểm uốn, vai trò của hai cấu trúc số (TCI, DAI) và
quy mô mẫu xuyên các bối cảnh ICRV.

\textbf{Bảng 4.10.} \emph{Tổng hợp ước lượng các nghiên cứu thành phần
thực nghiệm.}

\begin{longtable}[]{@{}
  >{\raggedright\arraybackslash}p{(\columnwidth - 12\tabcolsep) * \real{0.1429}}
  >{\raggedright\arraybackslash}p{(\columnwidth - 12\tabcolsep) * \real{0.1429}}
  >{\raggedleft\arraybackslash}p{(\columnwidth - 12\tabcolsep) * \real{0.1429}}
  >{\raggedright\arraybackslash}p{(\columnwidth - 12\tabcolsep) * \real{0.1429}}
  >{\raggedright\arraybackslash}p{(\columnwidth - 12\tabcolsep) * \real{0.1429}}
  >{\raggedright\arraybackslash}p{(\columnwidth - 12\tabcolsep) * \real{0.1429}}
  >{\raggedright\arraybackslash}p{(\columnwidth - 12\tabcolsep) * \real{0.1429}}@{}}
\toprule\noalign{}
\begin{minipage}[b]{\linewidth}\raggedright
NC
\end{minipage} & \begin{minipage}[b]{\linewidth}\raggedright
Bối cảnh (nhóm ICRV)
\end{minipage} & \begin{minipage}[b]{\linewidth}\raggedleft
N
\end{minipage} & \begin{minipage}[b]{\linewidth}\raggedright
Dạng hàm
\end{minipage} & \begin{minipage}[b]{\linewidth}\raggedright
Điểm uốn (FSTS)
\end{minipage} & \begin{minipage}[b]{\linewidth}\raggedright
Vai trò TCI
\end{minipage} & \begin{minipage}[b]{\linewidth}\raggedright
Vai trò DAI
\end{minipage} \\
\midrule\noalign{}
\endhead
\bottomrule\noalign{}
\endlastfoot
P2 & Trung Quốc, doanh nghiệp nhỏ và vừa (III) & 4.290 & U ngược &
47,8\% & nâng mặt bằng & nâng mặt bằng \\
P3 & Việt Nam (IV) & 2.958 & U ngược & 39--46\% & dương, có ý nghĩa
(β=0,179, biến công cụ) & không có ý nghĩa (Tầng 1) \\
P4 & Singapore (I) & 623 & gần tuyến tính dương & \textasciitilde88,6\%
(không định vị chắc) & yếu tố vệ sinh & FSTS²×DAI=+3,119** \\
P5 & Trung Quốc (III) & 4.544 & U ngược & 47--49\% (ổn định 2012--2024)
& dương (+0,28 lên +0,43) & không áp dụng \\
P6 & Phân tích tổng hợp (toàn cầu, 49 nền) & k=238 & dương nhỏ & không
áp dụng & không áp dụng & r̄=0,074; I²=62,4\% \\
P7 & Đa quốc gia (50 nền, I--VI) & 81.022 & U ngược & 43,6\% (M5); rõ
nhất ở Nhóm IV (43,0\%) & nâng mặt bằng (+0,141 sd) & nâng mặt bằng
(+0,201 sd) \\
P8 & Đảo nhỏ Thái Bình Dương (VI) & 1.450 & chữ U ngược mất cấu trúc
(FIP: trường hợp giới hạn) & không có & nâng mặt bằng, không điều tiết &
độ cong \& độ dốc n.s. (7 nền, mẫu đầy đủ); −1,339 chỉ ở bản dựng ba
cụm \\
P9 & Ấn Độ (IV) & 28.717 & U ngược chuyển sang sụp đổ ngưỡng & 61,8\%
rồi 40,7\% rồi tan biến & không áp dụng & tương tác Tầng 2 âm \\
P10 & Nhật Bản (I) & 2.168 & gần tuyến tính dương &
\textasciitilde90,1\% (không định vị) & nâng mặt bằng & nâng mặt bằng \\
\end{longtable}

\emph{Nguồn: tổng hợp từ các nghiên cứu thành phần P2--P10. ** p
\textless,01; *** p \textless,001. Điểm uốn xác nhận bằng kiểm định
Lind--Mehlum (2010). Điểm uốn P7 = 43,6\% ở mô hình đầy đủ M5 trên khung
50 nền (dải 43,6--51,5\% giữa các mô hình).}

Bảng 4.10 làm nổi bật ba quy luật xuyên suốt. Thứ nhất, \textbf{dạng hàm
chữ U ngược là quy luật chứ không phải ngoại lệ}, được xác nhận ở năm
trong chín nghiên cứu (P2, P3, P5, P7 và một phần P9), với các trường
hợp biên có ý nghĩa lý thuyết riêng: quan hệ gần tuyến tính dương ở biên
thể chế mạnh (P4 Singapore và P10 Nhật Bản, điểm uốn nằm ngoài miền quan
sát), việc chữ U ngược mất cấu trúc tại các nền đảo nhỏ (P8; FIP là
trường hợp giới hạn lý thuyết), và sự tan biến ngưỡng dưới biến động thể
chế cực đoan tại Ấn Độ (P9). Thứ hai, \textbf{điểm uốn dịch theo chiều
thể chế}: từ khoảng 88--90\% ở Singapore và Nhật Bản (Nhóm I) xuống
47--49\% ở Trung Quốc (III) và 39--46\% ở Việt Nam (IV), nhất quán với
phổ thể chế toàn mẫu của P7. Thứ ba, \textbf{TCI và DAI vận hành khác
nhau}: TCI nhất quán nâng mặt bằng năng suất, trong khi vai trò DAI phụ
thuộc bối cảnh (nén đường cong ở thể chế yếu, khuếch đại ở Singapore),
củng cố sự phân biệt cốt lõi của khung CDCM.

\begin{longtable}[]{@{}
  >{\raggedright\arraybackslash}p{(\columnwidth - 4\tabcolsep) * \real{0.3333}}
  >{\raggedright\arraybackslash}p{(\columnwidth - 4\tabcolsep) * \real{0.3333}}
  >{\raggedright\arraybackslash}p{(\columnwidth - 4\tabcolsep) * \real{0.3333}}@{}}
\toprule\noalign{}
\begin{minipage}[b]{\linewidth}\raggedright
Giả thuyết
\end{minipage} & \begin{minipage}[b]{\linewidth}\raggedright
Nội dung
\end{minipage} & \begin{minipage}[b]{\linewidth}\raggedright
Kết quả
\end{minipage} \\
\midrule\noalign{}
\endhead
\bottomrule\noalign{}
\endlastfoot
H1 (Phi tuyến chữ U ngược) & quốc tế hóa và hiệu quả có dạng chữ U ngược
& Xác nhận tại Trung Quốc, Việt Nam và toàn mẫu P7; xác nhận một phần
tại Singapore và Nhật Bản (gần tuyến tính, điểm uốn ngoài miền); không
áp dụng tại SIDS \\
H2 (TCI điều tiết dương) & TCI khuếch đại quốc tế hóa và hiệu quả & Xác
nhận một phần: TCI nâng mặt bằng năng suất tại tất cả bối cảnh
(+0,141***, 50 nền); uốn đường cong xác nhận tại Việt Nam nhưng không ý
nghĩa tại toàn mẫu P7 \\
H3 (DAI điều tiết) & DAI thay đổi độ dốc quốc tế hóa và hiệu quả & Hỗ
trợ một phần: hiệu ứng mức phổ quát (+0,201*\textbf{, 50 nền); tương tác
đường cong cùng chiều nén nhưng không ý nghĩa toàn mẫu; bằng chứng điều
tiết rõ nhất ở bối cảnh đơn quốc gia (P4: FSTS²×DAI = +3,119}); không có
ý nghĩa tại Việt Nam (Tầng 1) tương thích với hạn chế đo lường \\
H4 (Nhà quản trị điều tiết) & Kinh nghiệm/giới tính nhà quản trị & Không
được hỗ trợ trên toàn mẫu 50 nền (nữ quản lý: mức −0,104,
p\textless,001; tương tác FSTS không ý nghĩa); hiệu ứng phụ thuộc bối
cảnh \\
H5 (ICRV phổ thể chế) & Thể chế điều tiết phổ thể chế quốc tế hóa và
hiệu quả & Xác nhận một phần: \(Q_M = 17{,}35^{**}\) (\emph{df} = 4,
\emph{p} =,002) từ P6 MARA xác nhận dị biệt giữa chế độ; phổ thể chế
định hướng (E1a/E1b) chưa xác nhận do \emph{k} = 3 tại nhóm Frontier
(bin FR, hệ 5-regime của P6); phổ điểm uốn theo ICRV được xác nhận từ
P7 \\
H6 (Temporal heterogeneity) & Hình dạng quốc tế hóa và hiệu quả thay đổi
theo thời gian & Không xác nhận tại Trung Quốc (\(F = 1{,}83\),
\(p = {,}176\)); TP dịch chuyển tại Việt Nam giữa 2009 và 2015 \\
H1b (FIP, SIDS boundary condition) & Chữ U ngược mất cấu trúc tại SIDS;
FIP là trường hợp giới hạn & Ủng hộ theo hướng mất cấu trúc: trên mẫu
đầy đủ bảy nền Thái Bình Dương (\(N=1.450\), bảy cụm) độ dốc
(\(-0{,}085\), \(p_\text{wild}=,66\)) không có ý nghĩa và độ cong
(\(+0{,}696\), \(p_\text{wild}=,082\)) chỉ đạt mức biên. Hệ số âm mạnh
(\(-1{,}339^{***}\)) chỉ xuất hiện trên bản dựng dữ liệu hạn chế ba cụm,
mang tính minh họa và nhạy cảm với phiên bản dữ liệu (Mục 4.7.8) \\
H1c (Tan biến ngưỡng tại Ấn Độ, P9) & Ngưỡng U ngược \emph{tan biến}
dưới biến động thể chế cực đoan 2014--2025 & Xác nhận tại Ấn Độ: TP
61,8\% (2014) rồi 40,7\% (2022) rồi tan biến (2025, \(\beta_2\) NS);
Paternoster \(z = -7{,}94\), \(p < {,}0001\); tương tác DAI Tầng 2 (UPI)
âm (\(-4{,}02\), \(p = {,}004\)) \\
\end{longtable}

\textbf{Đối chiếu với ba phản bác chính.} Để khẳng định độ tin cậy của
các phát hiện trụ cột, cần xử lý tường minh ba phản bác hiển nhiên nhất.
Thứ nhất là nhân quả ngược và nội sinh: doanh nghiệp năng suất cao có
thể tự chọn lọc vào xuất khẩu, làm thổi phồng hệ số nhánh đi lên. Tuy
nhiên, định vị điểm uốn phụ thuộc vào dấu và độ lớn của hệ số bậc hai,
vốn ít nhạy với chiều của chọn lọc hơn hệ số bậc nhất; thêm vào đó, chữ
U ngược tái lập qua nhiều đợt khảo sát và nhiều bối cảnh có môi trường
chọn lọc khác nhau, và nghiên cứu Việt Nam (P3) dùng biến công cụ cho hệ
số TCI, củng cố cách đọc rằng phát hiện không phải tạo tác của nội sinh.
Thứ hai là thiên lệch chọn lọc do dùng mẫu chỉ doanh nghiệp xuất khẩu:
phản bác này được xử lý bằng việc ước lượng trên toàn mẫu gồm cả doanh
nghiệp có cường độ xuất khẩu bằng không, bằng phân rã hai biên tham gia
và cường độ, và bằng độ vững khi loại doanh nghiệp siêu nhỏ. Thứ ba là
tạo tác đo lường, nhất là nhiễu tiền tệ: việc trung tâm hóa cường độ
xuất khẩu, cắt đuôi phân vị trong từng cặp nền kinh tế và năm, và áp
dụng hiệu ứng cố định hai chiều đã hấp thụ chênh lệch mặt bằng giá; quan
trọng hơn, cấu trúc ba vùng được xác nhận độc lập bởi hai phương pháp có
điểm yếu không trùng nhau, là phân tích tổng hợp P6 và ước lượng gộp
toàn mẫu P7. Sự hội tụ này khiến ba phản bác trên khó có thể đồng thời
giải thích toàn bộ mô thức quan sát được.

\begin{center}\rule{0.5\linewidth}{0.5pt}\end{center}

\hypertarget{tuxe0i-liux1ec7u-tham-khux1ea3o-truxedch-dux1eabn-trong-chux1b0ux1a1ng-4}{%
\section{Tài liệu tham khảo (trích dẫn trong Chương
4)}\label{tuxe0i-liux1ec7u-tham-khux1ea3o-truxedch-dux1eabn-trong-chux1b0ux1a1ng-4}}

Arte, P., \& Larimo, J. (2022). Moderating influence of product
diversification on the international diversification-performance
relationship: A meta-analysis. \emph{Journal of Business Research}, 139,
281--295.

Bausch, A., \& Krist, M. (2007). The effect of context-related
moderators on the internationalization-performance relationship:
Evidence from meta-analysis. \emph{Management International Review,
47}(3), 319--347.

Begg, C. B., \& Mazumdar, M. (1994). Operating characteristics of a rank
correlation test for publication bias. \emph{Biometrics, 50}(4),
1088--1101.

Cameron, A. C., Gelbach, J. B., \& Miller, D. L. (2008). Bootstrap-based
improvements for inference with clustered errors. \emph{Review of
Economics and Statistics, 90}(3), 414--427.

Cohen, J. (1988). \emph{Statistical power analysis for the behavioral
sciences} (2nd ed.). Lawrence Erlbaum.

Cohen, W. M., \& Levinthal, D. A. (1990). Absorptive capacity: A new
perspective on learning and innovation. \emph{Administrative Science
Quarterly, 35}(1), 128--152.

Do, T. H., \& Phan, A. T. (2024, December). \emph{Internationalization
and firm performance: A meta-analysis review} {[}Paper presentation{]}.
The Sixth International Conference on Sustainable Development in
Economics, Business, and Finance (ICBEF), Can Tho University, Vietnam.

Do, T. H., \& Phan, A. T. (2025). Internationalization and firm
performance of firms in India: The role of top management. In M.
Bartekova (Ed.), \emph{International business research: Traditional and
creative approaches}. IntechOpen.
\url{https://doi.org/10.5772/intechopen.1011012}

Haans, R. F. J., Pieters, C., \& He, Z.-L. (2016). Thinking about U:
Theorizing and testing U- and inverted U-shaped relationships in
strategy research. \emph{Strategic Management Journal, 37}(7),
1177--1195.

Helpman, E., Melitz, M. J., \& Yeaple, S. R. (2004). Export versus FDI
with heterogeneous firms. \emph{American Economic Review, 94}(1),
300--316.

Hennart, J.-F. (2007). The theoretical rationale for a
multinationality--performance relationship. \emph{Management
International Review, 47}(3), 423--452.

Hsieh, C.-T., \& Klenow, P. J. (2009). Misallocation and manufacturing
TFP in China and India. \emph{Quarterly Journal of Economics, 124}(4),
1403--1448.

Johanson, J., \& Vahlne, J.-E. (1977). The internationalization process
of the firm --- A model of knowledge development and increasing foreign
market commitments. \emph{Journal of International Business Studies,
8}(1), 23--32.

Kaufmann, D., Kraay, A., \& Mastruzzi, M. (2011). The Worldwide
Governance Indicators: Methodology and analytical issues. \emph{Hague
Journal on the Rule of Law, 3}(2), 220--246.

Khanna, T., \& Palepu, K. G. (2010). \emph{Winning in emerging markets:
A road map for strategy and execution}. Harvard Business Press.

Kirca, A. H., Roth, K., Hult, G. T. M., \& Cavusgil, S. T. (2012). The
role of context in the multinationality-performance relationship: A
meta-analytic review. \emph{Global Strategy Journal, 2}(2), 108--121.

Lall, S. (1992). Technological capabilities and industrialization.
\emph{World Development, 20}(2), 165--186.

Lind, J. T., \& Mehlum, H. (2010). With or without U? The appropriate
test for a U-shaped relationship. \emph{Oxford Bulletin of Economics and
Statistics, 72}(1), 109--118.

Marano, V., Arregle, J.-L., Hitt, M. A., Spadafora, E., \& Van Essen, M.
(2016). Home country institutions and the
internationalization-performance relationship: A meta-analytic review.
\emph{Journal of Management, 42}(5), 1075--1110.

Melitz, M. J. (2003). The impact of trade on intra-industry
reallocations and aggregate industry productivity. \emph{Econometrica,
71}(6), 1695--1725.

Mizik, N., \& Jacobson, R. (2003). Trading off between value creation
and value appropriation: The financial implications of shifts in
strategic emphasis. \emph{Journal of Marketing, 67}(1), 63--76.

North, D. C. (1990). \emph{Institutions, institutional change and
economic performance}. Cambridge University Press.

Oster, E. (2019). Unobservable selection and coefficient stability:
Theory and evidence. \emph{Journal of Business \& Economic Statistics,
37}(2), 187--204.

Paternoster, R., Brame, R., Mazerolle, P., \& Piquero, A. (1998). Using
the correct statistical test for the equality of regression
coefficients. \emph{Criminology, 36}(4), 859--866.

Pisani, N., Garcia-Bernardo, J., \& Heemskerk, E. (2020). Does it pay to
be a multinational? A large-sample, cross-national replication assessing
the multinationality--performance relationship. \emph{Strategic
Management Journal, 41}(1), 152--172.

Verhoef, P. C., Broekhuizen, T., Bart, Y., Bhattacharya, A., Dong, J.
Q., Fabian, N., \& Haenlein, M. (2021). Digital transformation: A
multidisciplinary reflection and research agenda. \emph{Journal of
Business Research}, 122, 889--901.

Wagner, J. (2007). Exports and productivity: A survey of the evidence
from firm-level data. \emph{The World Economy, 30}(1), 60--82.

Williamson, O. E. (1985). \emph{The economic institutions of
capitalism}. Free Press.

World Bank. (2025b, September). \emph{Viet Nam economic update:
Nurturing high-tech talents} (Taking Stock, September 2025). World Bank
Group.

World Bank Enterprise Surveys. (2025). \emph{Enterprise surveys data}.
World Bank Group. \url{https://www.enterprisesurveys.org}

Wu, J., Xu, L., \& Yang, Z. (2022). Multinationality and performance of
emerging economy multinational enterprises: A twenty-year study.
\emph{International Business Review, 31}(5), 102003.

\begin{center}\rule{0.5\linewidth}{0.5pt}\end{center}

\end{document}
